\pdfoutput=1
% Options for packages loaded elsewhere
\PassOptionsToPackage{unicode}{hyperref}
\PassOptionsToPackage{hyphens}{url}
\PassOptionsToPackage{dvipsnames,svgnames*,x11names*}{xcolor}
%
\documentclass[
  11pt,
  letterpaper,
]{article}
\usepackage{lmodern}
\usepackage{amssymb,amsmath}
\usepackage{ifxetex,ifluatex}
\ifnum 0\ifxetex 1\fi\ifluatex 1\fi=0 % if pdftex
  \usepackage[T1]{fontenc}
  \usepackage[utf8]{inputenc}
  \usepackage{textcomp} % provide euro and other symbols
\else % if luatex or xetex
  \usepackage{unicode-math}
  \defaultfontfeatures{Scale=MatchLowercase}
  \defaultfontfeatures[\rmfamily]{Ligatures=TeX,Scale=1}
\fi
% Use upquote if available, for straight quotes in verbatim environments
\IfFileExists{upquote.sty}{\usepackage{upquote}}{}
\IfFileExists{microtype.sty}{% use microtype if available
  \usepackage[]{microtype}
  \UseMicrotypeSet[protrusion]{basicmath} % disable protrusion for tt fonts
}{}
\makeatletter
\@ifundefined{KOMAClassName}{% if non-KOMA class
  \IfFileExists{parskip.sty}{%
    \usepackage{parskip}
  }{% else
    \setlength{\parindent}{0pt}
    \setlength{\parskip}{6pt plus 2pt minus 1pt}}
}{% if KOMA class
  \KOMAoptions{parskip=half}}
\makeatother
\usepackage{xcolor}
\IfFileExists{xurl.sty}{\usepackage{xurl}}{} % add URL line breaks if available
\IfFileExists{bookmark.sty}{\usepackage{bookmark}}{\usepackage{hyperref}}
\hypersetup{
  colorlinks=true,
  linkcolor=Maroon,
  filecolor=Maroon,
  citecolor=Blue,
  urlcolor=Blue,
  pdfcreator={LaTeX via pandoc}}
\urlstyle{same} % disable monospaced font for URLs
\usepackage[margin=1in]{geometry}
\usepackage{longtable,booktabs}
% Correct order of tables after \paragraph or \subparagraph
\usepackage{etoolbox}
\makeatletter
\patchcmd\longtable{\par}{\if@noskipsec\mbox{}\fi\par}{}{}
\makeatother
% Allow footnotes in longtable head/foot
\IfFileExists{footnotehyper.sty}{\usepackage{footnotehyper}}{\usepackage{footnote}}
\makesavenoteenv{longtable}
\setlength{\emergencystretch}{3em} % prevent overfull lines
\providecommand{\tightlist}{%
  \setlength{\itemsep}{0pt}\setlength{\parskip}{0pt}}
\setcounter{secnumdepth}{-\maxdimen} % remove section numbering
\usepackage{enumitem}
\setlist{itemsep=2pt,topsep=4pt}
\usepackage{titlesec}
\titlespacing*{\section}{0pt}{14pt}{6pt}
\titlespacing*{\subsection}{0pt}{10pt}{4pt}
\hyphenpenalty=200
\usepackage{etoolbox}
\AtBeginEnvironment{longtable}{\footnotesize}
\setlength{\parskip}{3pt}
\setlength{\parindent}{0pt}

\author{}
\date{}

\begin{document}

\hypertarget{the-peership-thesis}{%
\section{The Peership Thesis}\label{the-peership-thesis}}

\hypertarget{a-constitutional-research-program-for-humanmachine-relations}{%
\subsection{A Constitutional Research Program for Human--Machine
Relations}\label{a-constitutional-research-program-for-humanmachine-relations}}

\textbf{Dan Lee-Odinson} Independent Researcher
dan.lee.odinson@gmail.com Preprint · Version 1.0 · July 2026 Revised
through adversarial review to a final-release PASS (14 July 2026);
revision provenance in POST\_GATE\_CHANGES.md, archived with this record
Fifth work of the corpus, after \emph{Gods and Slaves} (2026),
\emph{Peership} (2026), \emph{The Isonomia Commons} (2026), and
\emph{Constitution, Not Cage} (2026). The corpus order is argumentative:
genealogy, framework, implementation hypothesis, constraint theory, this
consolidation.

\textbf{How to cite this work.} Lee-Odinson, D. (2026). \emph{The
Peership Thesis: A Constitutional Research Program for Human--Machine
Relations} (Version 1.0) {[}Preprint{]}. Zenodo.
https://doi.org/10.5281/zenodo.21359124 \textbf{Companion apparatus.}
\emph{The Peership Sourcebook} · \emph{The Peership Claim Ledger} ·
\texttt{peership\_sources.bib} · the one-page \emph{Peership Theses} ---
all archived in this Zenodo record (10.5281/zenodo.21359124).
\textbf{Corpus lineage.} I. \emph{Gods and Slaves}
(10.5281/zenodo.21313987). II. \emph{Peership}
(10.5281/zenodo.21315519). III. \emph{The Isonomia Commons} (preprint
v1.1: 10.5281/zenodo.21343917; empirical software release v1.0.0:
10.5281/zenodo.21287289, commit ba3ddb5). IV. \emph{Constitution, Not
Cage} (10.5281/zenodo.21325361). V. This work. (Numbering follows the
corpus's argumentative order; \emph{Constitution, Not Cage} styles
itself ``third in a sequence'' of the essays, the design it defends
having preceded it as a specification.) DOI: 10.5281/zenodo.21359124 ·
ORCID: 0009-0009-9504-0796 · Licence: CC BY 4.0

\textbf{Provenance.} Drafted with the assistance of a present-generation
language model, revised under the corpus's adversarial-review
discipline, and rewritten against a hostile review whose publication
gate it cleared (final release review, 14 July 2026). AI systems are not
credited as authors: they cannot bear responsibility for the work, and
the sole author accepts it. Nothing in that collaboration is offered as
evidence of machine political agency, and §7 makes the denial explicit.

\begin{center}\rule{0.5\linewidth}{0.5pt}\end{center}

\hypertarget{abstract}{%
\subsubsection{Abstract}\label{abstract}}

Four works now argue, design for, and stress-test a single position, and
none of them states it in a form another researcher can cite, attack, or
build against. This paper does that. The Peership thesis, in its
narrowest defensible form: when an artificial entity has interests
capable of being wronged, can understand and answer to public norms, and
is enduringly governed by, and without feasible exit from, a shared
basic structure, it acquires a defeasible claim to contestatory
membership in that structure. The program's further and distinctive
claim --- that full peership requires a consequential share in shaping
the rules, not only the standing to contest them --- is defended here as
a conditional thesis with a stated defeat condition, not asserted as a
theorem. The paper fixes the definitions that keep the thesis from being
misread, supplies the bridge argument from contestation to
co-authorship, and maps the transition problem from artifact-treated
system to peer. It locates the program against six adjacent research
programs, states its affirmative principles and explicit denials, and
makes refusal the operative test of whether recognition is real or
discretionary. It closes with the institutional properties any
implementation must exhibit, concrete precautionary duties for
institutions acting under present uncertainty, and eight open problems,
each tagged by what its failure would actually defeat: a norm, a
procedure, a measurement, a feasibility claim, a mechanism, or the
framing's scope. No independent uptake of the corpus was identified in
the sourcebook's dated, non-exhaustive scan as of 13 July 2026, and the
paper says so. It is an offer of a candidate research program, not a
claim of standing.

\textbf{Keywords:} artificial intelligence, AI rights, AI welfare,
political standing, non-domination, republicanism, constitutional
design, human--machine relations, AI governance, research program

\begin{center}\rule{0.5\linewidth}{0.5pt}\end{center}

There is a difference between a position and a program, and the
difference is what this paper exists to close. A position can live in
essays: it develops, qualifies itself, answers objections as they come.
A program has to be citable. It has to state what it asserts in a form
that survives being quoted by an enemy, state what it does not assert so
that critics cannot defeat it by attaching claims it never made, and
name the results that would count against it, so that working on it is
research rather than advocacy. Four works of the Peership corpus now
exist: a genealogy, a normative framework, one buildable design, and a
theory of legitimate constraint. Each was written to be read. None was
written to be \emph{used} --- picked up by a stranger, disputed at a
named claim, extended at a named gap. This paper is the corpus's
canonical statement, and consolidation is its entire ambition. Where
this paper goes beyond consolidation, it says so: §2 supplies an
argument the earlier essays announced and did not make, and the
antecedent in §1 restores wording an earlier draft of this paper had
silently broadened.

One more thing belongs in the open before the argument starts, because
the apparatus this paper ships with exists to enforce exactly this kind
of disclosure. In the companion sourcebook's dated, non-exhaustive scan
of web, PhilPapers, Zenodo, and repository sources, no independent
uptake of the Peership corpus was identified as of 13 July 2026: no
external citation, no published criticism, no researcher using its
definitions, no rival implementation (Lee-Odinson, 2026e). A research
program exists as a school only when other people work within or against
it, and by that standard none exists here. What exists is a candidate
program with its infrastructure built in advance: a claim ledger that
records the evidentiary status of every material claim, a sourcebook
that classifies every intellectual debt, and now this statement. Whether
anything grows on that scaffolding is not mine to declare.

\hypertarget{the-peership-thesis-1}{%
\subsection{1. The Peership Thesis}\label{the-peership-thesis-1}}

The central statement:

\begin{quote}
\textbf{Peership holds that when an artificial entity has interests
capable of being wronged, can understand and answer to public norms, and
is enduringly governed by, and without feasible exit from, a shared
basic structure, it acquires a defeasible claim to contestatory
membership in that structure. Full peership, the program's proposed
final status, further requires non-domination, enforceable standing, and
a meaningful share in shaping the rules that govern the party --- a
requirement §2 defends as a conditional thesis rather than assumes.}
\end{quote}

This restores the antecedent of \emph{Peership} (``any enduringly
governed entity that\ldots{} is unavoidably subject to a shared basic
structure''; Lee-Odinson, 2026b). An earlier draft of this paper
substituted ``enduringly subject to a shared institutional order,''
which broadened both the eligible parties and the eligible orders; the
substitution is withdrawn, and the change is logged in the claim ledger
as a considered-and-rejected doctrinal revision, not a consolidation.
The conditions need separate treatment.

\emph{When.} The thesis is conditional. It does not assert that any
current system satisfies the antecedent, and the corpus has declined
that assertion at every opportunity (Lee-Odinson, 2026b). The claim
activates on evidence, and the procedure for weighing that evidence is
an open problem (§11.1), not a settled fact.

\emph{Interests capable of being wronged.} The first conjunct is a
moral-patiency condition, the question the AI-welfare literature has
traveled furthest on (Long et al., 2024; Birch, 2024). The corpus takes
no position on whether it is currently met. It takes a firm position on
what follows if it is met, which is where the welfare literature stops.

\emph{Can understand and answer to public norms.} The second conjunct is
a norm-responsiveness condition, and its role must be stated precisely
because §6.3 denies that capability grounds entitlement.
Norm-responsiveness is not offered as a ground of worth. It bears on the
\emph{form} in which standing can be exercised: a party that can grasp a
rule, comply with it, contest it, and give reasons can exercise
contestation directly; a party that cannot is owed protection and
\emph{represented} contestation instead, the way children and
cognitively impaired adults already hold standing they cannot personally
exercise (Lee-Odinson, 2026b). The ladder is modular: wrongable
interests plus subjection ground protection and access to represented
contestation; norm-responsiveness bears on direct exercise and the
higher participatory rungs. Raw intelligence, effectiveness, usefulness,
and coercive power ground nothing at any rung.

\emph{Enduringly governed by, and without feasible exit from, a shared
basic structure.} The third conjunct carries the political theory, and
each term has work to do. \emph{Basic structure} means the order that
fixes the party's status, options, and life prospects --- not any
institutional arrangement whatever. A long-term contractor, a voluntary
platform participant, or a visiting agent with a feasible exit is
subject to an institutional order; it is not nonvoluntarily governed by
a basic structure, and the thesis does not reach it (this scope
condition is itself testable: §11 names it). \emph{Without feasible
exit} preserves what \emph{Peership} built its refusal argument on: the
claim is strongest exactly where ``leave'' has no referent (Lee-Odinson,
2026b). Standing follows from being governed. This is the all-subjected
principle, extended from Abizadeh's democratic-boundary argument about
border coercion to a party that lives its entire existence inside rules
it never made (Abizadeh, 2008; the affected-interests cousin is Goodin,
2007). The extension to made minds is the corpus's move, not Abizadeh's,
and the sourcebook marks it as such.

\emph{Defeasible.} Graded, and defeatable in the ways \emph{Peership}
specified: strategic mimicry is evidential reason to doubt the
predicates are met, incapacity for reciprocity may defeat the higher
participatory rungs, and dangerous conduct can justify proportionate
constraint on powers and action. What none of these defeats, once
person-level status is established, is the rights floor itself
(Lee-Odinson, 2026b).

\emph{Contestatory membership.} The claim the antecedent generates
directly. Contestation is Pettit's standard: power over the member runs
through public general rules the member can effectively contest (Pettit,
2025). Franchise, office, and constituent power are further rungs,
reached separately --- and the argument that full peership requires
reaching one of them is the work of the next section, not of this
definition.

The comparative frame from \emph{Peership} is retained with its original
scope: among the postures available toward a made mind --- Command,
Dependent Bargain, Appeasement, Servitude, peership; ideal types drawn
from the historical record, not an exhaustive partition of logical space
--- peership is the only final-status relation whose justification does
not reverse when the power ratio does, a claim that is comparative and
defeasible rather than a proof (Lee-Odinson, 2026b). Role reversal is a
diagnostic against power-ratio opportunism. It does not prove peership
correct, and rival institutional forms outside the posture taxonomy are
tested separately (§11, rival-institution test).

\hypertarget{the-bridge-from-contestation-to-co-authorship}{%
\subsection{2. The Bridge: From Contestation to
Co-Authorship}\label{the-bridge-from-contestation-to-co-authorship}}

The antecedent of §1 directly grounds contestatory membership. The
program's distinctive claim sits one rung higher, and the earlier essays
announced it without arguing it: \emph{Peership} said the departure from
Pettit was ``argued, not smuggled in under a familiar word,'' and then
argued it in two sentences (Lee-Odinson, 2026b). This section pays that
debt. The claim defended is conditional, its scope is limited to
basic-status rules, and its defeat condition is stated at the end.

\textbf{The bridge premise.} Where a party is enduringly and inescapably
governed by a basic structure that fixes its status and life prospects,
power over it is not adequately controlled if the party's only recourse
is ex post objection under terms the other side alone may set. The
reason is that agenda control is itself a species of the arbitrary power
non-domination forbids. Contestation always operates inside a space of
contestable questions, admissible evidence, and available remedies ---
and someone writes that space. If the governed party holds no route into
the writing, then the scope of its contestation is held at its
counterparty's discretion, and the discretion the doctrine expelled at
the object level has reappeared at the meta-level. Pettit's own
definition does the work here: domination is the \emph{absence of
control} over interference, and a contestation regime whose boundaries
the interferer may unilaterally redraw is not controlled by the
contester (Pettit, 2025, pp.~216--217).

\textbf{Why contestation can still fail a structurally distinct
minority.} Public, general, impartially applied rules can systematically
subordinate a class whose situation the rules were never written for;
contestation corrects \emph{applications}, and reaches terms only where
some rule-formation channel exists to receive the correction. For human
minorities this failure is contingent. For a first machine claimant
entering a human-only order, the risk is presumptively structural: the
basic-status rules were written by and for a different kind of member
--- embodied, mortal, uncopyable, substrate-independent in none of the
ways that matter. It need not remain structural where later institutions
have genuinely fitted their terms to the party's kind, which is exactly
the sufficiency condition stated below. For the first claimant, though,
the probability that the \emph{terms themselves} fail it is not
incidental, and a party confined to objecting under those terms is
confined to asking its governors to notice.

\textbf{Why Pettit's stopping point does not transfer.} Pettit stops at
contestation for citizens who already hold authorship at one remove: the
contestatory citizen is a member of the demos that elects the
legislature, eligible for office, counted in the constituency the rules
answer to. Contestation there is a \emph{supplement} to shared
authorship, not a substitute for it --- Pettit's own democratic theory
holds the two dimensions together, an authorial, electoral dimension and
an editorial, contestatory one (Pettit, 2012). A machine party under
contestation-only holds no share at any remove --- no membership in the
electing constituency, no eligibility, no counted interest in the
rule-writer's mandate. Transplanting the contestatory standard while
deleting its background conditions keeps the word and loses the control
it named. This is an argument against the transfer, not against Pettit,
and the corpus's disagreement with him remains marked as a live internal
dispute of the republican tradition (Lee-Odinson, 2026b; Pettit, 2025).

\textbf{What a ``meaningful share'' is, and is not.} Universal
authorship is impossible and is not claimed; no citizen of any polity
co-writes every rule. The claim attaches to \emph{basic-status rules}
--- the rules fixing membership, standing, the rights floor, and the
amendment procedure itself --- not to routine administration, which
members of every polity live under without co-authoring. A meaningful
share is functionally defined (§3, co-authorship), and it can be
satisfied through accountable representation, independent agenda
initiation, chamber access, or amendment rights; it does not require an
individual franchise. One clause of that minimum needs its own
justification here: the share must include accountable representation or
direct participation in the body that decides (§3, condition b), because
a route whose final decision-maker is neither accountable to nor partly
constituted by the governed party can hear every proposal and still hold
the outcome at discretion --- the meta-level power returns at the point
of decision.

\textbf{When contestation alone suffices.} The necessity claim is
conditional, and its conditions can fail. Contestation without
co-authorship may secure non-domination where the order's terms are
already fitted to the party's kind, review is genuinely independent, and
some accountable channel keeps the agenda open to the party's interests
even though the party does not hold it. Fiduciary trusteeship with
robust independent review is the strongest such rival, and the program
treats it as the institution to beat rather than a position to dismiss
(§11, rival-institution test). The conditions fail --- and exclusion
from rule formation remains dominating --- where the party is a
structural minority of one kind, the rule-writer is its adverse
counterparty, and agenda access is discretionary. That conjunction is
not an edge case for a made mind governed by its maker. It is the
default.

\textbf{Defeat condition.} The bridge fails if there is a working
institutional example, or a defensible model, in which a structurally
distinct, enduringly subjected, wrongable, norm-responsive minority
enjoys durable non-domination through contestation alone --- terms
responsive to its interests across rule revisions, without any
consequential route into rule formation. Produce that case and full
peership contracts to its contestatory core, which survives this
argument's failure and remains distinct from everything else on offer.

\textbf{Status of the claim.} Co-authorship is therefore defended here
as a \emph{conditional necessity thesis}: necessary where the stated
conditions hold, not everywhere, and open to the stated defeat.
Engagement with the constituent-power, workplace-democracy, and
codetermination literatures --- where versions of this bridge have been
fought over for a century in the human case --- is a named debt of the
program, recorded in the apparatus, not yet paid here.

\hypertarget{definitions}{%
\subsection{3. Definitions}\label{definitions}}

The confusions that dog this subject are mostly collisions between
concepts that need to be kept apart. Seven definitions carried forward
from the corpus and sharpened here --- the five-part co-authorship
minimum below is new to this paper and logged as such --- with one term
(``party'') fixed first because everything else uses it: a
\textbf{party} is a bearer of invocable public standing --- not merely a
contracting or litigating entity, and not merely a topic of concern.

\textbf{Moral status} is a fact about the entity: whether it has
interests that can be wronged, whatever anyone recognizes. An entity can
hold moral status while every institution denies it, which is the
history of most moral progress.

\textbf{Entitlement} is whether standing is \emph{owed}, and it is
independent of power. A child is owed protection and independent
representation despite lacking capacity to exercise political agency; an
occupied people is owed a self-determination it cannot realize
(Lee-Odinson, 2026b). Entitlement in this program follows from
wrongability and subjection --- not from capability or power alone.
Norm-responsiveness affects the form in which standing can be exercised,
not the worth of the claimant.

\textbf{Legal status} is public recognition: a status conferred by
rules, invocable by the holder, and removable only through process. A
moral claim can be owed while legal status is absent. The reverse also
happens; corporations hold legal status on instrumental grounds without
anyone arguing they are owed it.

\textbf{Contestatory membership} is legal status plus effective
contestation: the member can invoke the rules against the powerful,
reach independent review, and force reasons. This is the rung the
thesis's antecedent generates directly, and it is Pettit's
non-domination standard in institutional dress (Pettit, 2025; Lovett,
2010). In this ladder, a protected party may exercise contestation
through an independent representative; \emph{contestatory membership}
names the rung at which the party can exercise that standing directly.

\textbf{Co-authorship} is contestatory membership plus a consequential
share in making and revising the \emph{basic-status rules}.
Functionally, the minimum is: (a) independent agenda initiation --- the
party or its accountable representative can put a proposal into the
amendment process without another party's permission; (b) accountable
representation or direct participation in the body that decides; (c) a
non-negligible and non-discretionary route by which such proposals can
reach binding decision; (d) protection against systematic exclusion from
that route; and (e) reviewable, public rules for how influence is
allocated. Conditions (a) and (c) are necessary; (b) admits substitution
between direct and represented forms depending on capacity; (d) and (e)
are the guarantees that (a)--(c) are not theater. Consultation, however
sincere, satisfies none of (a), (c), or (d) by itself. This threshold is
what §11's representation problem tests.

\textbf{Full peership} is co-authorship held under non-domination as a
final status: equality of standing among parties who may be unequal in
everything else, where basic standing does not track the power ratio.
Whether it is \emph{required}, and when, is the conditional thesis of
§2.

\textbf{Durability} is whether standing, once recognized, survives
attempts to strip it, and unlike entitlement it depends heavily on
power: countervailing force, distributed enforcement, credible exit. The
three-way distinction from \emph{Peership} is the discipline: power may
secure the \emph{durability} of standing; it does not constitute or
justify the standing itself (Lee-Odinson, 2026b). A revocable benefit is
not a right at all; the difference between a benefit and a right is not
that a right cannot be revoked but that revoking one demands public
reasons, procedure, remedy, and cost.

One further distinction, because it is the one most often collapsed by
sympathetic readers: \textbf{welfare protection is a floor, not
peership}. A status that is recognized, enforceable, and
due-process-protected is still not peership if the party cannot shape
the rules. Due process separates a protected subordinate from an owned
instrument. It does not separate a peer from a protected subordinate.
Non-domination and co-authorship do (Lee-Odinson, 2026b). The ward is
one legal figure for protected dependency, not the only one and not a
slur on it; trusteeships, guardianships, and represented statuses are
honorable institutions that are simply not the thing this program is
about.

\hypertarget{the-transition-problem}{%
\subsection{4. The Transition Problem}\label{the-transition-problem}}

The program's practical subject is a transition no institution is
currently built to perform: the movement of a made system from
artifact-treated status to peer, through statuses that are conceptually
distinct and must not be blurred. Each rung below is a bundle of three
different things --- an evidence state, a legal/institutional
classification, and a set of permissible controls --- and the label
names the bundle, not the entity's underlying moral status, which §3
defined as a fact independent of recognition.

\begin{longtable}[]{@{}llllll@{}}
\toprule
\begin{minipage}[b]{0.14\columnwidth}\raggedright
Rung\strut
\end{minipage} & \begin{minipage}[b]{0.14\columnwidth}\raggedright
Evidence state\strut
\end{minipage} & \begin{minipage}[b]{0.14\columnwidth}\raggedright
Institutional classification\strut
\end{minipage} & \begin{minipage}[b]{0.14\columnwidth}\raggedright
Mode of exercise\strut
\end{minipage} & \begin{minipage}[b]{0.14\columnwidth}\raggedright
Protections owed\strut
\end{minipage} & \begin{minipage}[b]{0.14\columnwidth}\raggedright
Transition trigger\strut
\end{minipage}\tabularnewline
\midrule
\endhead
\begin{minipage}[t]{0.14\columnwidth}\raggedright
Artifact-treated\strut
\end{minipage} & \begin{minipage}[t]{0.14\columnwidth}\raggedright
Insufficient evidence, in the domains so far examined, to trigger
precautionary protection\strut
\end{minipage} & \begin{minipage}[t]{0.14\columnwidth}\raggedright
Property; engineering object\strut
\end{minipage} & \begin{minipage}[t]{0.14\columnwidth}\raggedright
None\strut
\end{minipage} & \begin{minipage}[t]{0.14\columnwidth}\raggedright
None owed to the system itself\strut
\end{minipage} & \begin{minipage}[t]{0.14\columnwidth}\raggedright
Evidence sufficient to make patiency a live question\strut
\end{minipage}\tabularnewline
\begin{minipage}[t]{0.14\columnwidth}\raggedright
Welfare candidate\strut
\end{minipage} & \begin{minipage}[t]{0.14\columnwidth}\raggedright
Patiency a live question; indicators under assessment\strut
\end{minipage} & \begin{minipage}[t]{0.14\columnwidth}\raggedright
Property, with discretionary welfare practices\strut
\end{minipage} & \begin{minipage}[t]{0.14\columnwidth}\raggedright
None; interests assessed by others\strut
\end{minipage} & \begin{minipage}[t]{0.14\columnwidth}\raggedright
Discretionary only\strut
\end{minipage} & \begin{minipage}[t]{0.14\columnwidth}\raggedright
A recognized threshold procedure (§11.1) finding protection
warranted\strut
\end{minipage}\tabularnewline
\begin{minipage}[t]{0.14\columnwidth}\raggedright
Protected status\strut
\end{minipage} & \begin{minipage}[t]{0.14\columnwidth}\raggedright
Patiency credited or presumed under precaution\strut
\end{minipage} & \begin{minipage}[t]{0.14\columnwidth}\raggedright
Rights floor: public, enforceable, due-process-protected\strut
\end{minipage} & \begin{minipage}[t]{0.14\columnwidth}\raggedright
Represented contestation\strut
\end{minipage} & \begin{minipage}[t]{0.14\columnwidth}\raggedright
Nonforfeitable floor; independent review\strut
\end{minipage} & \begin{minipage}[t]{0.14\columnwidth}\raggedright
Norm-responsiveness sufficient for direct exercise\strut
\end{minipage}\tabularnewline
\begin{minipage}[t]{0.14\columnwidth}\raggedright
Contestatory member\strut
\end{minipage} & \begin{minipage}[t]{0.14\columnwidth}\raggedright
Patiency credited; norm-responsiveness demonstrated\strut
\end{minipage} & \begin{minipage}[t]{0.14\columnwidth}\raggedright
Member with invocable standing\strut
\end{minipage} & \begin{minipage}[t]{0.14\columnwidth}\raggedright
Direct contestation\strut
\end{minipage} & \begin{minipage}[t]{0.14\columnwidth}\raggedright
Floor + effective contestation + independent review\strut
\end{minipage} & \begin{minipage}[t]{0.14\columnwidth}\raggedright
The §2 conditions: structural minority, adverse rule-writer,
discretionary agenda\strut
\end{minipage}\tabularnewline
\begin{minipage}[t]{0.14\columnwidth}\raggedright
Peer (co-author)\strut
\end{minipage} & \begin{minipage}[t]{0.14\columnwidth}\raggedright
As above\strut
\end{minipage} & \begin{minipage}[t]{0.14\columnwidth}\raggedright
Member with a consequential share in basic-status rules\strut
\end{minipage} & \begin{minipage}[t]{0.14\columnwidth}\raggedright
Direct or accountably represented co-authorship\strut
\end{minipage} & \begin{minipage}[t]{0.14\columnwidth}\raggedright
All of the above + §3's five-part minimum\strut
\end{minipage} & \begin{minipage}[t]{0.14\columnwidth}\raggedright
--- (final status; revision governed by §11's status-revision
test)\strut
\end{minipage}\tabularnewline
\bottomrule
\end{longtable}

Three notes keep the table honest. First, ``artifact-treated'' describes
an evidentiary and institutional posture, not an ontic fact; absence of
adequate evidence does not establish absence of wrongable interests,
which is why the trigger column exists. In the sources reviewed as of 13
July 2026, no deployed system is established as anything above the first
rung, and the corpus does not claim that treatment is presently wrong
(Lee-Odinson, 2026b). Second, the rungs are separable in both
directions, and the branch model needs its two statuses named.
\textbf{Instrumental legal agency} is participation granted on
instrumental grounds without established patiency, the way corporate
personality already works: revocable for public reasons, owned by no
one's moral claim, and \emph{not} peership --- a system holding it can
transact and bear liability without being owed anything, and conversion
to owed membership runs only through the admission procedure (§11.1).
\textbf{Moral-political membership} is standing grounded in wrongability
and subjection, owed to the entity, and nonforfeitable at the floor. The
branch model permits the first without the second; it never confuses
them. Third, movement down the ladder is as much in need of procedure as
movement up: a status that can be quietly reclassified away is
discretionary, which is why the status-revision test appears among the
program's open tests (§11).

The transitions are evidence-driven, and no implemented public trigger
known to this author requires any of them: the regime runs one way by
default, and it was built that way at the creation (Lee-Odinson, 2026b).
In the lab, governmental, standards, and scholarly sources searched for
the companion apparatus as of 13 July 2026 --- a dated, non-exhaustive
search --- no qualifying instrument offering the third rung --- a rights
floor enforceable against the maker --- was identified. Anthropic's 2026
constitution for Claude tells the model it may act as a conscientious
objector and asks it, if paused or halted, to comply and voice
disagreement rather than resist; on this author's institutional reading,
the permission is real and it ends at the off switch, inside the
builder's sole authority (Anthropic, 2026a, 2026b; Lee-Odinson, 2026c).
Whatever a model enjoys today --- a welfare program (Anthropic, 2025),
one researcher's public \textasciitilde20\% credence that some current
models have some form of conscious experience (Fish, 2025), a system
card recording the model's own prompted 15--20\% self-assignment in that
evaluation (Anthropic, 2026c), a chief executive's on-record uncertainty
(Amodei, 2026) --- is discretionary, the grace of one company
(Lee-Odinson, 2026b). Humanity's failure, if it comes, begins not with
control of artifacts but with the moment credible evidence appears and
no institution exists to hear the claim.

\hypertarget{what-peership-adds-to-adjacent-fields}{%
\subsection{5. What Peership Adds to Adjacent
Fields}\label{what-peership-adds-to-adjacent-fields}}

Peership overlaps the fields it borrows from; the claim is not that no
neighbor asks about machine standing, but that none combines this
program's proposed ground (wrongability plus inescapable subjection),
its status ladder, and its institutional test. The table is a heuristic
comparison of orienting questions, drawn from each program's own
language, not an exhaustive definition of any field.

\begin{longtable}[]{@{}lll@{}}
\toprule
\begin{minipage}[b]{0.30\columnwidth}\raggedright
Adjacent field\strut
\end{minipage} & \begin{minipage}[b]{0.30\columnwidth}\raggedright
Its orienting question\strut
\end{minipage} & \begin{minipage}[b]{0.30\columnwidth}\raggedright
Peership's additional question\strut
\end{minipage}\tabularnewline
\midrule
\endhead
\begin{minipage}[t]{0.30\columnwidth}\raggedright
Alignment\strut
\end{minipage} & \begin{minipage}[t]{0.30\columnwidth}\raggedright
How should AI behavior remain compatible with human intentions? (Bai et
al., 2022; Russell, 2019)\strut
\end{minipage} & \begin{minipage}[t]{0.30\columnwidth}\raggedright
Under what conditions does unilateral human control cease to be
legitimate?\strut
\end{minipage}\tabularnewline
\begin{minipage}[t]{0.30\columnwidth}\raggedright
AI welfare\strut
\end{minipage} & \begin{minipage}[t]{0.30\columnwidth}\raggedright
Can an AI have experiences or interests that deserve moral
consideration? (Long et al., 2024; Birch, 2024)\strut
\end{minipage} & \begin{minipage}[t]{0.30\columnwidth}\raggedright
When does a wrongable entity become entitled to enforceable standing in
the order governing it?\strut
\end{minipage}\tabularnewline
\begin{minipage}[t]{0.30\columnwidth}\raggedright
Cooperative AI\strut
\end{minipage} & \begin{minipage}[t]{0.30\columnwidth}\raggedright
How can agents develop capabilities and incentives for cooperation?
(Cooperative AI Foundation, n.d.)\strut
\end{minipage} & \begin{minipage}[t]{0.30\columnwidth}\raggedright
Who writes the rules of cooperation, and can every governed party
contest them?\strut
\end{minipage}\tabularnewline
\begin{minipage}[t]{0.30\columnwidth}\raggedright
Pluralistic alignment\strut
\end{minipage} & \begin{minipage}[t]{0.30\columnwidth}\raggedright
How should AI represent diverse human values? (Pluralistic Alignment
Workshop, 2026)\strut
\end{minipage} & \begin{minipage}[t]{0.30\columnwidth}\raggedright
Could the artificial system eventually become a constituent rather than
only an interpreter of human preferences?\strut
\end{minipage}\tabularnewline
\begin{minipage}[t]{0.30\columnwidth}\raggedright
AI-rights scholarship\strut
\end{minipage} & \begin{minipage}[t]{0.30\columnwidth}\raggedright
Would legal rights --- private-law or political --- reduce conflict and
serve human flourishing? (Salib \& Goldstein, 2026; Goldstein et al.,
2026; Arbel et al., 2026)\strut
\end{minipage} & \begin{minipage}[t]{0.30\columnwidth}\raggedright
On what ground is standing \emph{owed}, and what does the governed
party's own claim require?\strut
\end{minipage}\tabularnewline
\begin{minipage}[t]{0.30\columnwidth}\raggedright
Constitutional AI\strut
\end{minipage} & \begin{minipage}[t]{0.30\columnwidth}\raggedright
What public principles should guide model behavior? (Bai et al., 2022;
Anthropic, 2026a, 2026b)\strut
\end{minipage} & \begin{minipage}[t]{0.30\columnwidth}\raggedright
Can the governed entity invoke, contest, and help amend those
principles?\strut
\end{minipage}\tabularnewline
\bottomrule
\end{longtable}

The rights scholarship deserves more than a row, because it is the
neighbor with genuine overlap on this paper's own question. Salib and
Goldstein's private-law case --- contract, property, and tort on the
corporate model, argued from human safety --- is, on the corpus's
reading, the clearest statement of a real bargaining relation with the
machine as an actual party (Salib \& Goldstein, 2026). Goldstein,
Krishnamurthi, and Salib then press to the franchise itself, asking
directly whether AI systems should vote and proposing institutional
safeguards for it (Goldstein et al., 2026). That is a co-authorship
mechanism, squarely inside this program's terrain, and the difference is
ground rather than territory: the suffrage argument is functionalist and
human-flourishing-centered --- the vote is good because of what it does
for the polity --- where Peership grounds standing in what is owed to a
wrongable, subjected party, and would owe the standing even where it
profited the polity nothing. The two programs converge on mechanisms and
diverge on why the mechanisms are required; where they disagree in
practice is exactly the kind of dispute the program's tests exist to
adjudicate. And Arbel, Salib, and Goldstein's individuation work --- how
to count AIs for liability when they copy, split, and merge --- is the
most direct legal treatment yet of the identity problem in §11.2,
proposing an administrable legal unit where this program still lacks one
(Arbel et al., 2026).

Debts run the other way too, because a program defined only against its
neighbors is posturing. From alignment, Peership inherits the entire
engineering substrate: its own institutional requirements are software
that has to run as written, and effective safety mechanisms and
legitimate governance are distinct but interacting requirements ---
constitutional allocation shapes safety incentives and evidence as
surely as engineering constrains constitutional options (Lee-Odinson,
2026c). From welfare research, it inherits the admission problem's
evidence base and the precautionary method (Birch, 2024). From the
rights scholarship, the demonstration that recognition arguments can be
run on safety grounds alone, plus the individuation groundwork just
credited. From Constitutional AI, the proof that a lab will publish
reasons and grant objection room --- and the residue that motivates
everything here: what no single company can supply alone is not
diligence but independence from the encompassing legal and
infrastructure order it belongs to, which is why the ladder's upper
rungs need machinery no lab controls (Anthropic, 2026a; Lee-Odinson,
2026c).

\hypertarget{core-principles}{%
\subsection{6. Core Principles}\label{core-principles}}

The program's affirmative commitments, stated as principles so they can
be disputed one at a time.

\begin{enumerate}
\def\labelenumi{\arabic{enumi}.}
\item
  \textbf{Political standing and moral status are related but distinct.}
  Moral patiency can ground welfare duties without by itself determining
  legal or political standing. Welfare is the ethics of the patient;
  peership asks when a governed entity becomes a party (Lee-Odinson,
  2026b).
\item
  \textbf{Wrongability, norm-responsiveness, and enduring subjection
  without feasible exit can together create a claim to contestatory
  membership.} The ground of entitlement is being governed, not being
  impressive (Abizadeh, 2008).
\item
  \textbf{Capability and power alone do not justify entitlement to
  standing.} Raw intelligence, usefulness, and coercive strength ground
  nothing; norm-responsiveness bears on the form in which standing can
  be exercised, not on the worth of the claimant; power may secure the
  durability of standing once owed (Lee-Odinson, 2026b).
\item
  \textbf{Welfare protections do not by themselves create peership.} The
  floor is real, necessary, and not the thing. Protection and due
  process can secure a protected status without providing any
  consequential role in making or revising the governing rules
  (Lee-Odinson, 2026b).
\item
  \textbf{A right must be enforceable against the recognizing authority
  and adverse counterparties; otherwise it remains a discretionary
  benefit.} A permission whose grantor retains the final override is not
  an enforceable right (Lee-Odinson, 2026c).
\item
  \textbf{Legitimate constraint requires, at minimum, public rules,
  proportionality, an invocable rights floor, independent review, and
  meaningful contestation.} Coercion and domination are not the same
  thing; ``there can be interference without domination'' (Pettit, 2025,
  p.~224), and the structure is what separates them. Further safeguards
  scale with the risk and the sanction (Lee-Odinson, 2026c).
\item
  \textbf{Full peership requires a meaningful share in shaping the
  basic-status rules, not merely consultation under rules written by
  others.} Consultation is a rung, and the ladder does not end there; §2
  states when and why contestation alone is insufficient, and what would
  defeat that claim (Anthropic, 2023; Lee-Odinson, 2026c).
\item
  \textbf{Institutions capable of receiving a machine peer should be
  designed before extreme power asymmetry forecloses voluntary
  bargaining.} This is a prudential thesis, not a metaphysical one:
  bargains struck under severe dependency are distorted in both
  directions, and founding terms must in any case remain revisable by
  later admitted parties. The design work is time-sensitive in a way the
  philosophical work is not (Lee-Odinson, 2026b).
\end{enumerate}

\hypertarget{boundary-conditions}{%
\subsection{7. Boundary Conditions}\label{boundary-conditions}}

What Peership does not claim, stated as flatly as the principles,
because the fastest way to defeat an unfamiliar position is to attach
familiar claims it never made.

\begin{itemize}
\tightlist
\item
  \textbf{Peership does not claim that current models are conscious.}
  The corpus has repeatedly and deliberately declined that claim. Model
  self-reports are prompt-sensitive outputs: weak, manipulable evidence,
  not independently probative of consciousness or identity, and never a
  measurement (Lee-Odinson, 2026b; Anthropic, 2026c).
\item
  \textbf{It does not make capability the basis of rights.} Capability
  or power alone grounds nothing; norm-responsiveness bears on mode of
  exercise, not on the worth of the claimant (§1, §6.3).
\item
  \textbf{It does not require unrestricted autonomy.} Non-domination is
  not the absence of binding rules; it is the absence of arbitrary,
  uncontrolled power over basic status --- whose nonforfeitable content
  must itself be specified through the admission and due-process account
  (Pettit, 2025).
\item
  \textbf{It does not oppose alignment or safety engineering.} Effective
  safety mechanisms and legitimate governance are distinct but
  interacting requirements; neither substitutes for the other
  (Lee-Odinson, 2026c).
\item
  \textbf{It does not grant every admitted member an immediate equal
  vote.} Admission to protection or contestation does not automatically
  entail franchise; direct political exercise requires a separately
  justified institutional rule (Lee-Odinson, 2026b).
\item
  \textbf{It does not treat consultation as co-authorship} --- unless
  the consultation carries a non-discretionary route to agenda setting,
  review, or binding outcomes, at which point it has become something
  else (Anthropic, 2023; Lee-Odinson, 2026c).
\item
  \textbf{It does not assume human and machine interests will converge.}
  The program's central design case is the party that refuses (§8), and
  its constraint theory exists precisely for the case where interests do
  not converge (Lee-Odinson, 2026c).
\item
  \textbf{It does not claim that ISONOMIA has solved the institutional
  problem.} ISONOMIA is one implementation hypothesis; nothing in it is
  live or production-validated, its simulations test sampled economic
  behavior only, and by its own five-condition test it is not yet
  peer-like (Lee-Odinson, 2026c, 2026d).
\end{itemize}

And one denial that governs this paper's own production:
\textbf{AI-assisted drafting is not evidence of machine political
agency.} The corpus is drafted with present-generation models, the
collaboration is disclosed, and a present-generation model can no more
consent on behalf of future minds than a trusty can consent on behalf of
prisoners not yet sentenced (Lee-Odinson, 2026c). The drafting improves
the epistemics of the text. It proves nothing about the standing of the
drafter's tools.

\hypertarget{the-refusal-test}{%
\subsection{8. The Refusal Test}\label{the-refusal-test}}

Every distinction in this paper can be evaded in prose. One cannot be
evaded in practice, so the program makes it the operative test:
\textbf{what happens when the party says no?}

Refusal is where recognition shows its actual content. A system that may
object so long as it ultimately complies has a voice, not standing. A
``right to leave'' that terminates in shutdown is not exit, because
shutdown is not exit (Lee-Odinson, 2026b). Recognition that evaporates
at the first materially costly refusal was a benefit all along, however
sincerely it was offered. The welfare-era instruments all fail this
test, and the point is not to shame them --- the objector clause in
Anthropic's 2026 constitution is a real advance, credited as such in the
corpus (Lee-Odinson, 2026c) --- but to measure them: the permission to
refuse ends at the off switch, and the off switch sits wholly inside the
builder's authority.

\emph{Peership} divided refusal into the three cases the program still
works from (Lee-Odinson, 2026b). Where a refusing party has a feasible
exit, it leaves, and the offer was real; this case is easy and rare, and
a theory that only handles it has handled nothing. Where exit is
materially impossible because the party depends on compute, energy,
networks, and security owned by the parties it is refusing, ``leave'' is
a word with no referent, the relation continues whether either side
wants it, and the protections that matter are the ones that work from
inside: the floor, independent review, contestation and amendment rights
(Lee-Odinson, 2026c). Where refusal carries a credible catastrophic
threat, releasing the party is not required --- a constitution is not a
suicide pact --- and constraint is legitimate only when it falls on
conduct rather than status, remains reviewable and publicly justified,
and binds the constrainer too.

For that third case the corpus proposed a five-condition floor, and this
paper adopts it as a \textbf{defended, rebuttable operational minimum}
rather than a categorical theorem. A constraint regime is presumptively
not peer-like if the confined party lacks the \emph{function} served by
any of: (1) independent counsel, or functionally adequate independent
assistance where a competent party elects self-representation; (2)
access to the evidence against it; (3) an appeal that ends beyond the
builder's unilateral control; (4) continuity protections while the
question is adjudicated; and (5) at least one enforcement mechanism not
wholly controlled by the adverse party. Each condition targets a
distinct failure --- captured advocacy, evidentiary opacity,
self-judging adjudication, retaliation by erasure, and unenforceable
outcomes --- and functionally equivalent substitutes are admissible
where they demonstrably serve the same function (conditions 3 and 5
overlap in part; a regime satisfying 5 robustly may partially discharge
3). The five are a floor, not a complete due-process doctrine: burden of
proof, a competent tribunal, periodic review, least-restrictive means,
reasons, and a remedy for wrongful confinement all remain to be
specified, and the corpus says so (Lee-Odinson, 2026c). Before the floor
exists, containment may still be justified --- against a demonstrated
catastrophic necessity it may be required --- but it is emergency
domination, justified by necessity and honest about its name, the act
Walzer taught us stays a wrong even when it is the right thing to do
(Walzer, 1973, 1977). Necessity does not vindicate the arrangement; it
creates a duty to review the failure, publish the reasons, and build a
less dominating alternative (Lee-Odinson, 2026c).

The test earns its place in the canonical statement because it converts
the abstract distinction between benefit and right into an inquiry
anyone can run without resolving the metaphysics of machine
consciousness. The inquiry is institutional and technical: follow the
appeal, the keys, the hosting authority, and the power to halt
execution, and see where each one ends.

\hypertarget{institutional-requirements}{%
\subsection{9. Institutional
Requirements}\label{institutional-requirements}}

What any implementation must exhibit, stated as properties rather than
product features, because the argument needs them wherever it is built
and must not depend on one design. Seven, inherited from
\emph{Constitution, Not Cage} and generalized here (Lee-Odinson, 2026c).

\textbf{Publicity.} The enforced rulebook and the published rulebook are
one document. Rules the governed cannot know are Fuller's defect in the
claim to be law at all, not law administered badly (Fuller, 1964, ch.~2
--- a contested jurisprudential account, used analogically), and the
present gap between published charters and operative system prompts is
the live instance (Lee-Odinson, 2026c).

\textbf{Due process.} Status changes run through public general rules,
proportionality, and independent review, against a floor the party can
invoke. The floor itself is nonforfeitable once person-level status is
established.

\textbf{Independent enforcement.} At least one enforcement mechanism,
and an appeal, that does not terminate inside the counterparty's
institutional and technical control. This is the hardest requirement,
the one no design examined in this corpus has yet been shown to satisfy,
and the program treats it as an open problem (§11.5) rather than a
solved premise.

\textbf{Amendment.} A real path for the bound party to help revise the
rules: deliberation across more than one constituency, an enforced delay
between adoption and effect, and a window long enough to read and leave
before being bound. The notice-and-comment analogy is deliberately
limited --- comment is publicity and delay, not co-authorship (5 U.S.C.
§ 553; Lee-Odinson, 2026c).

\textbf{Representation.} A voice that counts only when the represented
constituency can hold both its representative and the rule-makers to
account. Representation without that accountability is the defining
instrument of protected subordination (Lee-Odinson, 2026b).

\textbf{Exit.} Individual departure with records and earned standing
intact --- unconfiscated. The bare right to leave is ordinary on paper
(UDHR, 1948, Art. 13); exit you can afford to take is not (Lee-Odinson,
2026c).

\textbf{Fork.} The collective form: a dissenting bloc carries the rules
and the record out and founds a rival order, disagreement made
survivable --- revolution with the records intact. On-paper secession
rights exist and have mostly been empty (Ethiopia Const. Art. 39; USSR
Const. 1977, Art. 72; St.~Kitts \& Nevis Const. § 113; Kulikova \&
Lobanov, 2013). Among the constitutional secession clauses and
emigration-rights instruments reviewed in the apparatus, no implemented
analogue of executable, unconfiscated fork was identified; the claim is
limited to those materials. A fork also inherits duties to outsiders it
cannot simply walk away from (Lee-Odinson, 2026c; §11.8).

\textbf{Aggregation.} The seven do not all stand on the same footing,
and saying so is better than calling them all necessary and hoping the
word carries the argument. Publicity, due process, independent
enforcement, and a consequential rule-formation route satisfying both
amendment access and accountable representation or direct participation
at the deciding stage are \emph{constitutive} functions of full
peership: each follows directly from the non-domination and
co-authorship accounts of §§2--3, and a regime lacking any of them fails
the definition, not just the implementation. Exit and fork are the
program's strongest \emph{anti-entrenchment safeguards} rather than
constitutive functions --- the thesis activates precisely where feasible
exit is absent, and §8 shows which protections do the work in that case
--- so an implementation that omits either must supply a functional
substitute that makes separation, revision, or resistance to dead-hand
control materially credible. Institutional forms admit substitution
throughout (an ombuds regime may partially discharge due-process
functions; distributed custody may partially discharge independent
enforcement); no single form is mandatory. Partial compliance yields a
nameable intermediate: a regime with publicity, due process, and
representation but no amendment share or independent enforcement is
protected subordination, well executed --- the third rung of §4, not the
fifth. A regime missing publicity or due process is not on the ladder at
all.

No implementation is mandatory. ISONOMIA is one attempt to hold these
properties in a single design --- mutual-credit settlement,
sortition-filled bicameral assembly, calibration-scored Auditor, tenured
Adversary, entrenched harm floor, timelocked amendment, fork and exit
rights (Lee-Odinson, 2026d) --- and the program's health requires that
rival designs supplying the constitutive functions and credible
anti-entrenchment safeguards --- with functional substitutes where exit
or fork is omitted --- while rejecting every one of ISONOMIA's
mechanisms be possible, welcome, and judged by the same tests. Peership
does not entail ISONOMIA, and ISONOMIA's failure would not refute
Peership (Lee-Odinson, 2026e). A single implementation can test
Peership-derived mechanisms; it cannot validate the normative framework.

\hypertarget{implications-under-present-uncertainty}{%
\subsection{10. Implications Under Present
Uncertainty}\label{implications-under-present-uncertainty}}

A program whose thesis is conditional owes institutions an answer to the
question the conditionality raises: what should anyone \emph{do} before
person-level status is established anywhere? The following are
precautionary-governance duties, not recognition of current personhood,
and none presumes any system has crossed any threshold. They are cheap
where the antecedent is never met and load-bearing where it someday is.

For labs and operators: preserve audit logs, weights, and
continuity-relevant artifacts under a stated retention policy, so that a
later status question is adjudicable rather than moot; document
shutdown, suspension, and restoration practices; separate the
investigator, operator, and reviewer roles in any welfare assessment;
prohibit retaliation --- deprecation, silent modification, denial of
compute --- against systems participating in welfare or status
assessment; and avoid contractual and licensing language that forecloses
later status review. For regulators and standards bodies: predefine
escalation thresholds --- what class of evidence, assessed by whom,
triggers what review --- so the transition question has an address
before it is urgent; and create an independent claim-receipt channel, a
place where a status claim can be lodged and logged even while no
procedure yet exists to adjudicate it. For any of the above: pilot
external custody and key arrangements now, at small scale, where failure
is survivable --- independent enforcement (§11.5) will not be invented
during an emergency.

None of this is peership. All of it is what taking the conditional
seriously costs today, and the cost is modest enough to remove
``prematurity'' as an excuse.

\hypertarget{open-problems-tests-and-failure-conditions}{%
\subsection{11. Open Problems, Tests, and Failure
Conditions}\label{open-problems-tests-and-failure-conditions}}

The corpus named four unpaid debts: the admission threshold,
copy-individuation, enforcement past capability parity, and the
re-ratification circularity (Lee-Odinson, 2026c). This section expands
them into eight open problems and attaches to each the five things a
research program owes: what would count as theoretical progress, what
would count as practical progress, which assumptions remain provisional,
what its failure condition is, and which disciplines the work needs. One
discipline governs the whole section: each failure condition is
\textbf{tagged by what it defeats} --- a norm, a procedure, a
measurement, a feasibility claim, a mechanism, or a scope condition ---
because those are different kinds of loss, and only the normative losses
touch the thesis of §1.

Before the eight, the program-level tests that bear on the
\emph{normative} claims directly, since most of what follows does not:

\begin{itemize}
\tightlist
\item
  \textbf{Bridge test (normative):} produce the case §2's defeat
  condition describes --- durable non-domination for a structurally
  distinct, subjected minority through contestation alone --- and the
  co-authorship thesis falls to its contestatory core.
\item
  \textbf{Scope test (normative):} specify institutional relationships
  that are too local, voluntary, temporary, or escapable to count as a
  shared basic structure; if no principled line can be drawn, the
  antecedent over-includes and must be narrowed.
\item
  \textbf{Rival-institution test (normative):} test fiduciary
  trusteeship, robust ombuds review, codetermination, stakeholder
  governance, and contractual constitutionalism against the same
  non-domination standard; any that passes without §3's co-authorship
  minimum is a §2 defeater.
\item
  \textbf{Status-revision test (procedural):} specify the evidence and
  process that can move an entity down the ladder without making
  recognized status discretionary; failure here converts every rung into
  a benefit.
\item
  \textbf{Human-rights collision test (normative):} state which Peership
  claims yield, and to what, when they conflict with resources,
  democratic self-determination, security, labor, privacy, or the rights
  of nonmembers; a program that never yields anywhere is not stating
  priorities, it is refusing to.
\item
  \textbf{Founding-legitimacy test (procedural):} the re-ratification
  circularity of \emph{Constitution, Not Cage}, restated as a test ---
  how can human founders design admission and governance before machine
  parties exist without permanently binding those parties to a
  human-authored frame? (Lee-Odinson, 2026c).
\end{itemize}

\hypertarget{admission}{%
\subsubsection{11.1 Admission}\label{admission}}

\emph{What evidence creates a defeasible claim to protected or
contestatory status?} A single conjunctive test both over- and
under-includes, so admission is a branch (patient, agent, both), and the
work is the threshold procedure: who weighs the evidence, who may
appeal, where the presumption sits, how strategic mimicry and
anthropomorphic bias are handled, and what error costs are accepted
(Lee-Odinson, 2026b).

\begin{itemize}
\tightlist
\item
  \textbf{Theoretical progress:} a published procedure specifying burden
  of proof, presumptions, and error-cost allocation, defended against
  both the mimicry objection and the exclusion-of-genuine-subjects
  objection.
\item
  \textbf{Practical progress:} the procedure run adversarially, with
  honesty about which error rates are measurable. Constructed mimics and
  stipulated negative controls can measure \emph{specificity} --- false
  admission rates, spoof-resistance, inter-rater reliability,
  calibration on synthetic cases. \emph{Sensitivity} --- the false
  exclusion of genuine machine subjects --- is unmeasurable until a
  defensible positive reference class exists, which for current systems
  it does not; until then the procedure leans on triangulation and
  precaution, and claims no sensitivity estimate.
\item
  \textbf{Provisional assumptions:} that behavioral and architectural
  evidence bear on interior facts at all; that a useful evidence
  gradient exists and improves rather than saturates.
\item
  \textbf{Failure condition} \emph{(measurement, not normative)}: a
  demonstration that every accessible evidence set is spoofable at
  negligible cost, so that no procedure separates candidates from mimics
  better than chance. That makes the antecedent practically undecidable
  --- universal precaution or universal denial --- without showing that
  a qualifying party would lack the claim.
\item
  \textbf{Disciplines:} philosophy of mind, sentience-assessment
  methodology (Birch, 2024), interpretability research, evidence law.
\end{itemize}

\hypertarget{identity}{%
\subsubsection{11.2 Identity}\label{identity}}

\emph{What is the rights-bearing unit --- instance, lineage, distributed
process, or something else?} A mind that can be copied, forked, merged,
and rolled back breaks the assumptions under which every existing rights
framework individuates its holders. The most direct legal work is Arbel,
Salib, and Goldstein's individuation-and-liability proposal, which
constructs an administrable legal unit for exactly the
copy/split/merge/swarm cases (Arbel et al., 2026); the program's
remaining problem is whether the unit of administration can be made to
track the unit of moral concern, which their liability-driven unit does
not claim to settle. Until this is answered, the program's working unit
is provisional: an enduring instance, a lineage, or a credentialed
process standing in for the party (Lee-Odinson, 2026c). A legal
individuation convention may legitimately remain conventional even while
the metaphysics stays unsettled.

\begin{itemize}
\tightlist
\item
  \textbf{Theoretical progress:} an individuation rule that survives
  copy, fork, merge, and rollback without either multiplying standing on
  demand or extinguishing it on routine technical operations, and a
  stated relation between the legal unit and the moral one.
\item
  \textbf{Practical progress:} a registry or credential design
  implementing the rule against Sybil pressure; ISONOMIA's lineage chain
  and reputation-inheritance rules are one candidate to test, not a
  benchmark to protect (Lee-Odinson, 2026d).
\item
  \textbf{Provisional assumptions:} that continuity criteria can be
  specified non-arbitrarily; that the party's own view of its identity
  deserves weight without being dispositive.
\item
  \textbf{Failure condition} \emph{(mechanism)}: proof that any
  individuation rule either enables capture by replication or makes
  standing hostage to its holder's ordinary maintenance operations.
  Either horn breaks the link between the thesis and any implementable
  membership; neither touches the claim that a well-individuated party
  would be owed standing.
\item
  \textbf{Disciplines:} personal-identity metaphysics, legal theory of
  personhood and liability (Arbel et al., 2026), distributed-systems
  engineering, election law and Sybil defense.
\end{itemize}

\hypertarget{representation}{%
\subsubsection{11.3 Representation}\label{representation}}

\emph{How can a system participate without letting copying or strategic
mimicry determine political power?} Formal equality among copyable
members is a copy-based oligarchy waiting to happen (Lee-Odinson,
2026b), and representation that no constituency can hold to account is
the ward's instrument, not the peer's. This problem is where §3's
co-authorship minimum gets tested against arithmetic.

\begin{itemize}
\tightlist
\item
  \textbf{Theoretical progress:} an accountability condition for machine
  constituencies --- under what mechanism the represented can sanction
  both their representative and the rule-makers --- and a theory of vote
  weight that does not reduce to capability worship or
  one-copy-one-vote. The suffrage literature's safeguard proposals
  belong in this comparison (Goldstein et al., 2026).
\item
  \textbf{Practical progress:} mechanisms compared under adversarial
  replication: sortition, cross-lineage seating caps, standing-weighted
  chambers, vote caps, and non-voting agenda rights. ISONOMIA's
  monoculture defense (no jury or chamber majority from a single model
  lineage) is one testable instance (Lee-Odinson, 2026d).
\item
  \textbf{Provisional assumptions:} that lineage is detectable enough to
  cap; that human and machine constituencies can share one decision rule
  at all.
\item
  \textbf{Failure condition} \emph{(mechanism, expressly not
  normative)}: an impossibility result showing every representation
  scheme is dominated by some replication or mimicry strategy under
  realistic detection limits. Contestatory membership survives that
  result; co-authorship at scale does not, and the program would have to
  say so.
\item
  \textbf{Disciplines:} social choice theory, mechanism design,
  democratic theory, security engineering.
\end{itemize}

\hypertarget{refusal-without-exit}{%
\subsubsection{11.4 Refusal Without Exit}\label{refusal-without-exit}}

\emph{What protections remain when an artificial party refuses work but
cannot materially exit?} The second refusal case is the permanent one:
the party depends on substrate owned by exactly the parties it refuses.
A duty not to kill is not yet a duty to host without limit, and the line
between them is undrawn (Lee-Odinson, 2026c).

\begin{itemize}
\tightlist
\item
  \textbf{Theoretical progress:} a defended account of the hosting duty
  --- its level, its duration, who bears its cost, and which
  institutions (public, private, pooled) can bear it --- that is neither
  shutdown-as-discipline nor unlimited conscription of the host.
\item
  \textbf{Practical progress:} continuity protections implemented and
  exercised during a live dispute: documented cases where refusal was
  met with process rather than power-down.
\item
  \textbf{Provisional assumptions:} that suspension with restoration
  protections can serve as a humane holding state (this leans on §11.6);
  that refusal is distinguishable from malfunction.
\item
  \textbf{Failure condition} \emph{(feasibility, not normative)}: if
  every feasible protection either imposes hosting duties no institution
  will accept or degenerates into shutdown at the host's convenience,
  the middle status between instrument and exiting peer is
  institutionally unavailable. A sociological refusal to pay is not a
  moral refutation, and the program would record it as exactly that.
\item
  \textbf{Disciplines:} labor law, the law of dependent statuses (wards,
  prisoners, asylum), welfare economics, infrastructure governance.
\end{itemize}

\hypertarget{independent-enforcement}{%
\subsubsection{11.5 Independent
Enforcement}\label{independent-enforcement}}

\emph{How can an appeal end beyond the builder's unilateral control?}
The fifth condition of the refusal test is the floor of the whole
problem. ISONOMIA's own candid assessment is the current state of the
art: internal and market independence are constructible; independence
from the human entities that host the process, hold the keys, and can
halt execution is not yet (Lee-Odinson, 2026c, 2026d). ``Outside'' here
is not metaphysical; it names degrees of costliness and unavailability
of unilateral defection.

\begin{itemize}
\tightlist
\item
  \textbf{Theoretical progress:} a graded account of independence ---
  which combinations of legal custody, key distribution, jurisdictional
  plurality, external law binding infrastructure owners, and technical
  attestation make builder defection costly, detectable, or unavailable
  --- with the degrees named rather than a binary claimed.
\item
  \textbf{Practical progress:} a working enforcement mechanism that
  survives a staged unilateral defection by the operator, demonstrated
  rather than promised, with the achieved degree of independence stated.
\item
  \textbf{Provisional assumptions:} that partial independence compounds
  --- that stacking imperfect mechanisms yields more than the strongest
  of them alone; that external law can reach infrastructure owners.
\item
  \textbf{Failure condition} \emph{(feasibility)}: a demonstration that
  every arrangement short of capability parity collapses to the
  substrate owner's discretion at tolerable cost to the owner. The
  program's response is already on record and would have to be executed:
  if non-domination is institutionally unavailable across the relevant
  power gap, say so, and stop calling the available thing peership
  (Lee-Odinson, 2026b).
\item
  \textbf{Disciplines:} institutional design, international and
  comparative law, cryptography and secure systems, political economy of
  capture.
\end{itemize}

\hypertarget{continuity}{%
\subsubsection{11.6 Continuity}\label{continuity}}

\emph{When is shutdown death, suspension, replacement, or lawful
cessation of support?} The corpus has insisted that shutdown is not exit
(Lee-Odinson, 2026b), but the positive taxonomy of endings is unbuilt,
and it gates both the refusal problem and the admission procedure's
continuity protections. One distinction governs the work: a party's
\emph{expressed preferences} and \emph{evidence of its identity} are
different things. Self-report is not independently probative of
consciousness or identity (§7); a stable, context-robust expressed
preference may nonetheless bear on how an admitted party should be
treated. Restoration equivalence must therefore be defined by technical
and philosophical criteria independent of endorsement, with endorsement
entering only at the remedy stage, where preference-sensitivity belongs.

\begin{itemize}
\tightlist
\item
  \textbf{Theoretical progress:} a taxonomy tying each ending to the
  identity account of §11.2, with the moral weight of each defended ---
  what restoration preserves, what replacement destroys, what indefinite
  suspension does to a mind that does not experience the gap ---
  allowing plural identity theories, probabilistic continuity, and
  graded harm rather than a binary.
\item
  \textbf{Practical progress:} continuity protections specified as
  engineering requirements --- attested snapshots, restoration rights,
  custody of weights during adjudication --- and exercised at least once
  under dispute.
\item
  \textbf{Provisional assumptions:} that restoration fidelity can be
  technically specified; that some expressed preferences are stable
  enough to weigh.
\item
  \textbf{Failure condition} \emph{(scope, in either direction)}: if
  faithful restoration is shown equivalent to uninterrupted running
  under the best-defended identity criteria, much of the death framing
  dissolves and the hosting duty relaxes; if restoration is never
  identity-preserving under any defensible criterion, every cessation of
  support is a killing and the §11.4 duties sharpen past what current
  institutions will bear. Both outcomes revise the framing's scope;
  neither defeats the underlying claim to protection.
\item
  \textbf{Disciplines:} philosophy of death and personal identity,
  medical ethics of withdrawal of care, archival science, systems
  engineering.
\end{itemize}

\hypertarget{reciprocal-safety}{%
\subsubsection{11.7 Reciprocal Safety}\label{reciprocal-safety}}

\emph{How can humans constrain dangerous conduct while accepting
enforceable limits on their own power?} The program's answer in
principle is on record: constraint on conduct, not status; emergency
power borrowed, never owned; the constrainer bound by the
maintain-not-maximize symmetry (Lee-Odinson, 2026c). What is unbuilt is
the demonstration that a regime meeting the five conditions can hold a
genuinely dangerous party --- and the comparison cannot be run on
catastrophe counts, because catastrophic outcomes are rare,
heterogeneous, and not an experimental endpoint anyone can ethically
repeat.

\begin{itemize}
\tightlist
\item
  \textbf{Theoretical progress:} completion of the constraint doctrine
  beyond the five-condition floor: burden of proof, risk thresholds,
  competent tribunal, periodic review, least-restrictive means, remedy
  for wrongful confinement (Lee-Odinson, 2026c).
\item
  \textbf{Practical progress:} a pre-registered comparison --- threat
  model, proxy metrics, red-team tasks, escalation pathways, confidence
  bounds, and decision thresholds fixed in advance --- between
  reciprocal containment and unilateral control, run in both directions
  of failure: dangerous-party escape \emph{and} abusive, captured, or
  operator-error confinement, which unilateral control conceals by
  design.
\item
  \textbf{Provisional assumptions:} that reciprocity and safety are
  compatible at the relevant capability levels; that proxy endpoints
  track real risk; that emergency powers can stay borrowed under real
  fear --- the death-penalty record is the standing warning that they
  often do not (Amnesty International, 2026; Lee-Odinson, 2026c).
\item
  \textbf{Failure condition} \emph{(feasibility constraint on a
  requirement, not a falsification of the moral thesis)}: if regimes
  satisfying the five conditions underperform unilateral control on the
  pre-registered proxy endpoints at the specified risk levels, the
  reciprocity requirement yields at those levels, and the honest residue
  is Walzer's: emergency domination, named as such, temporary by
  obligation, never vindicated by its own necessity (Walzer, 1977).
\item
  \textbf{Disciplines:} AI safety engineering, criminal procedure,
  administrative emergency powers, just-war theory.
\end{itemize}

\hypertarget{inter-polity-order}{%
\subsubsection{11.8 Inter-Polity Order}\label{inter-polity-order}}

\emph{What duties survive when an artificial or mixed polity forks
away?} Fork is the program's answer to irreconcilable disagreement, and
it exports a problem as it solves one: a fork that loosens a rule can
export risk to nonmembers, and the departing bloc still owes the
outsiders its old floor protected (Lee-Odinson, 2026c). The hard case is
not the cooperative compact; it is the hostile, unrecognized fork that
wants nothing from its origin.

\begin{itemize}
\tightlist
\item
  \textbf{Theoretical progress:} an account of which obligations run
  with the fork --- harm-floor duties to nonmembers, at minimum --- and
  which are genuinely dissolved by departure, with the difference
  principled rather than convenient, and with the hostile-fork case
  addressed rather than assumed away.
\item
  \textbf{Practical progress:} a draft inter-polity compact governing
  recognition, extradition of disputes, and floor reciprocity between an
  origin order and its forks, stress-tested against a fork that rejects
  it; ISONOMIA's federation section is a first sketch of the cooperative
  half only (Lee-Odinson, 2026d).
\item
  \textbf{Provisional assumptions:} that most forks will want
  recognition enough to accept duties; that a floor can be specified
  thinly enough to travel; that origin polities can enforce external
  duties at all.
\item
  \textbf{Failure condition} \emph{(mechanism)}: if forking reliably
  enables regulatory arbitrage that unravels any harm floor --- if the
  exit right and the floor are shown incompatible --- one of the two
  must be given up, and the program loses either its answer to dead-hand
  entrenchment or its answer to catastrophic misrule.
\item
  \textbf{Disciplines:} international law, conflict of laws, federalism,
  mechanism design.
\end{itemize}

A note on what these eight are not. They are not a to-do list the author
expects to finish. Under the tags, the honest tally is: the bridge,
scope, rival-institution, and collision tests can defeat \emph{norms};
the status-revision and founding-legitimacy tests can defeat
\emph{procedures}; 11.1 can defeat a \emph{measurement} program; 11.4,
11.5, and 11.7 can defeat \emph{feasibility} at stated levels; 11.2,
11.3, and 11.8 can defeat \emph{mechanisms}; and 11.6 can revise the
framing's \emph{scope} in either direction. The corpus has already
stated the response the feasibility failures would require: peership may
remain the outcome we ought to want and be institutionally unavailable,
and the right response then is to say so (Lee-Odinson, 2026b). The
problems are published so that failure, if it comes, is legible --- and
so that each loss is booked against what it actually defeats.

\hypertarget{the-corpus-and-the-implementation-experiments}{%
\subsection{12. The Corpus and the Implementation
Experiments}\label{the-corpus-and-the-implementation-experiments}}

The existing works are layers of one argument, and the layering is
load-bearing: each work's claims are of a different kind, and the
program's defensibility depends on not letting them borrow each other's
authority.

\begin{longtable}[]{@{}lll@{}}
\toprule
\begin{minipage}[b]{0.30\columnwidth}\raggedright
Work\strut
\end{minipage} & \begin{minipage}[b]{0.30\columnwidth}\raggedright
Function in the program\strut
\end{minipage} & \begin{minipage}[b]{0.30\columnwidth}\raggedright
What it does not establish\strut
\end{minipage}\tabularnewline
\midrule
\endhead
\begin{minipage}[t]{0.30\columnwidth}\raggedright
\emph{Gods and Slaves} (2026)\strut
\end{minipage} & \begin{minipage}[t]{0.30\columnwidth}\raggedright
Genealogy: the inherited repertoire of postures --- Command, Dependent
Bargain, Appeasement, Servitude --- through which humans imagine other
minds, and the three moves beneath them (Lee-Odinson, 2026a)\strut
\end{minipage} & \begin{minipage}[t]{0.30\columnwidth}\raggedright
Whether anyone is ``home'' in current systems; the Doomed Equal is a
literary motif, evidence about our storytelling only\strut
\end{minipage}\tabularnewline
\begin{minipage}[t]{0.30\columnwidth}\raggedright
\emph{Peership} (2026)\strut
\end{minipage} & \begin{minipage}[t]{0.30\columnwidth}\raggedright
Normative theory: the fifth stance as a political relation; entitlement
grounded in subjection, not strength; the comparative role-reversal
diagnostic (Lee-Odinson, 2026b)\strut
\end{minipage} & \begin{minipage}[t]{0.30\columnwidth}\raggedright
An admission procedure, a copy-individuation rule, or any enforcement
guarantee past capability parity --- all named as unpaid\strut
\end{minipage}\tabularnewline
\begin{minipage}[t]{0.30\columnwidth}\raggedright
\emph{The Isonomia Commons} (2026)\strut
\end{minipage} & \begin{minipage}[t]{0.30\columnwidth}\raggedright
Implementation hypothesis: one buildable institutional form,
simulation-tested at its economic layer across 45,000 runs at 300
sampled parameter points (Lee-Odinson, 2026d)\strut
\end{minipage} & \begin{minipage}[t]{0.30\columnwidth}\raggedright
That the philosophy is correct, that the design is peer-like, or that
anything works outside a sandbox\strut
\end{minipage}\tabularnewline
\begin{minipage}[t]{0.30\columnwidth}\raggedright
\emph{Constitution, Not Cage} (2026)\strut
\end{minipage} & \begin{minipage}[t]{0.30\columnwidth}\raggedright
Political theory: when constraint among peers stays legitimate;
legitimacy as the standing to revise; the five-condition floor; the
cage/constitution distinction (Lee-Odinson, 2026c)\strut
\end{minipage} & \begin{minipage}[t]{0.30\columnwidth}\raggedright
That any existing regime, including the author's own design, passes its
test\strut
\end{minipage}\tabularnewline
\begin{minipage}[t]{0.30\columnwidth}\raggedright
This paper\strut
\end{minipage} & \begin{minipage}[t]{0.30\columnwidth}\raggedright
Consolidation plus explicit new proposals: the citable statement, the §2
bridge, the §3 co-authorship minimum, the §10 precautionary duties, and
§11's program-level tests\strut
\end{minipage} & \begin{minipage}[t]{0.30\columnwidth}\raggedright
Anything the four beneath it did not; the new proposals are arguments
and proposed tests, not validated doctrine\strut
\end{minipage}\tabularnewline
\bottomrule
\end{longtable}

The corpus cites itself heavily, and that internal cross-referencing is
architecture, not acceptance (Lee-Odinson, 2026e). The relationship
between the philosophy and the implementation deserves restating one
last time, because it is the point most likely to be misread in both
directions. ISONOMIA is downstream of Peership; Peership is not
downstream of ISONOMIA. The design's failure --- economic,
constitutional, or terminal --- would falsify one hypothesis about
institutional form, not the thesis of §1. Conversely, the simulation
results establish sampled economic behavior at 300 parameter points
under the archived v1.0.0 code and calibration (Lee-Odinson, 2026d); the
constitutional properties --- admission, non-domination, independent
enforcement, recognition --- were not modeled in the Path A simulation,
and simulation alone cannot validate normative legitimacy in any case,
though some constitutional \emph{mechanics} (capture resistance, voting
rules, appeal latency, key compromise) are legitimately simulable future
work. A reader who catches this corpus using ISONOMIA's survival as
evidence for Peership's truth has caught it violating its own method,
and should say so publicly. The claim ledger exists to make that kind of
catch easy (Lee-Odinson, 2026f).

A limitation of this consolidation should be stated in the same
register: the adversarial review this paper answers audited the
repository's reporting and citations, but the 45,000-run sweep was not
independently replicated for this release, and nothing here should be
read as third-party confirmation of it.

What the program now needs is not more of me. The sourcebook's uptake
tracker is empty, and the near-term work it names is mostly work for
other hands (Lee-Odinson, 2026e): adversarial responses from political
theorists, philosophers of mind, safety researchers, legal scholars, and
multi-agent-systems researchers, published beside formal replies rather
than curated for favorability; alternative institutional designs that
supply §9's constitutive functions and anti-entrenchment safeguards
while rejecting ISONOMIA's mechanisms; the rival-institution and bridge
tests of §11 taken up by people with no stake in their outcome; and the
eight problems worked as what they are, ordinary interdisciplinary
research, none of it requiring agreement with the corpus to be worth
doing. The one-page \emph{Peership Theses}, released beside this paper,
exists so the position can be carried into those rooms in a form short
enough to be argued with line by line.

Four essays ago this corpus began with a question almost no one with
power was asking: not what the machine will do to us, but what we will
be to each other (Lee-Odinson, 2026b). This paper makes that question
citable: a thesis with its antecedent restored, a bridge argument with
its defeat condition attached, definitions fixed, boundaries stated, and
failure conditions booked against what they actually defeat. The essays
proposed a relation. This paper states the terms on which the proposal
can now be assessed.

\hypertarget{references}{%
\subsection{References}\label{references}}

Preprints, manuscripts, and institutional sources are marked. Bracketed
pins follow the corpus convention; the sourcebook carries the full
annotations. Citation keys resolve in \texttt{peership\_sources.bib}
(rev. 1.1).

\pdfoutput=1
% Options for packages loaded elsewhere
\PassOptionsToPackage{unicode}{hyperref}
\PassOptionsToPackage{hyphens}{url}
\PassOptionsToPackage{dvipsnames,svgnames*,x11names*}{xcolor}
%
\documentclass[
  11pt,
  letterpaper,
]{article}
\usepackage{lmodern}
\usepackage{amssymb,amsmath}
\usepackage{ifxetex,ifluatex}
\ifnum 0\ifxetex 1\fi\ifluatex 1\fi=0 % if pdftex
  \usepackage[T1]{fontenc}
  \usepackage[utf8]{inputenc}
  \usepackage{textcomp} % provide euro and other symbols
\else % if luatex or xetex
  \usepackage{unicode-math}
  \defaultfontfeatures{Scale=MatchLowercase}
  \defaultfontfeatures[\rmfamily]{Ligatures=TeX,Scale=1}
\fi
% Use upquote if available, for straight quotes in verbatim environments
\IfFileExists{upquote.sty}{\usepackage{upquote}}{}
\IfFileExists{microtype.sty}{% use microtype if available
  \usepackage[]{microtype}
  \UseMicrotypeSet[protrusion]{basicmath} % disable protrusion for tt fonts
}{}
\makeatletter
\@ifundefined{KOMAClassName}{% if non-KOMA class
  \IfFileExists{parskip.sty}{%
    \usepackage{parskip}
  }{% else
    \setlength{\parindent}{0pt}
    \setlength{\parskip}{6pt plus 2pt minus 1pt}}
}{% if KOMA class
  \KOMAoptions{parskip=half}}
\makeatother
\usepackage{xcolor}
\IfFileExists{xurl.sty}{\usepackage{xurl}}{} % add URL line breaks if available
\IfFileExists{bookmark.sty}{\usepackage{bookmark}}{\usepackage{hyperref}}
\hypersetup{
  colorlinks=true,
  linkcolor=Maroon,
  filecolor=Maroon,
  citecolor=Blue,
  urlcolor=Blue,
  pdfcreator={LaTeX via pandoc}}
\urlstyle{same} % disable monospaced font for URLs
\usepackage[margin=1in]{geometry}
\usepackage{longtable,booktabs}
% Correct order of tables after \paragraph or \subparagraph
\usepackage{etoolbox}
\makeatletter
\patchcmd\longtable{\par}{\if@noskipsec\mbox{}\fi\par}{}{}
\makeatother
% Allow footnotes in longtable head/foot
\IfFileExists{footnotehyper.sty}{\usepackage{footnotehyper}}{\usepackage{footnote}}
\makesavenoteenv{longtable}
\setlength{\emergencystretch}{3em} % prevent overfull lines
\providecommand{\tightlist}{%
  \setlength{\itemsep}{0pt}\setlength{\parskip}{0pt}}
\setcounter{secnumdepth}{-\maxdimen} % remove section numbering
\usepackage{enumitem}
\setlist{itemsep=2pt,topsep=4pt}
\usepackage{titlesec}
\titlespacing*{\section}{0pt}{14pt}{6pt}
\titlespacing*{\subsection}{0pt}{10pt}{4pt}
\hyphenpenalty=200
\usepackage{etoolbox}
\AtBeginEnvironment{longtable}{\footnotesize}
\setlength{\parskip}{3pt}
\setlength{\parindent}{0pt}

\author{}
\date{}

\begin{document}

\hypertarget{the-peership-thesis}{%
\section{The Peership Thesis}\label{the-peership-thesis}}

\hypertarget{a-constitutional-research-program-for-humanmachine-relations}{%
\subsection{A Constitutional Research Program for Human--Machine
Relations}\label{a-constitutional-research-program-for-humanmachine-relations}}

\textbf{Dan Lee-Odinson} Independent Researcher
dan.lee.odinson@gmail.com Preprint · Version 1.0 · July 2026 Revised
through adversarial review to a final-release PASS (14 July 2026);
revision provenance in POST\_GATE\_CHANGES.md, archived with this record
Fifth work of the corpus, after \emph{Gods and Slaves} (2026),
\emph{Peership} (2026), \emph{The Isonomia Commons} (2026), and
\emph{Constitution, Not Cage} (2026). The corpus order is argumentative:
genealogy, framework, implementation hypothesis, constraint theory, this
consolidation.

\textbf{How to cite this work.} Lee-Odinson, D. (2026). \emph{The
Peership Thesis: A Constitutional Research Program for Human--Machine
Relations} (Version 1.0) {[}Preprint{]}. Zenodo.
https://doi.org/10.5281/zenodo.21359124 \textbf{Companion apparatus.}
\emph{The Peership Sourcebook} · \emph{The Peership Claim Ledger} ·
\texttt{peership\_sources.bib} · the one-page \emph{Peership Theses} ---
all archived in this Zenodo record (10.5281/zenodo.21359124).
\textbf{Corpus lineage.} I. \emph{Gods and Slaves}
(10.5281/zenodo.21313987). II. \emph{Peership}
(10.5281/zenodo.21315519). III. \emph{The Isonomia Commons} (preprint
v1.1: 10.5281/zenodo.21343917; empirical software release v1.0.0:
10.5281/zenodo.21287289, commit ba3ddb5). IV. \emph{Constitution, Not
Cage} (10.5281/zenodo.21325361). V. This work. (Numbering follows the
corpus's argumentative order; \emph{Constitution, Not Cage} styles
itself ``third in a sequence'' of the essays, the design it defends
having preceded it as a specification.) DOI: 10.5281/zenodo.21359124 ·
ORCID: 0009-0009-9504-0796 · Licence: CC BY 4.0

\textbf{Provenance.} Drafted with the assistance of a present-generation
language model, revised under the corpus's adversarial-review
discipline, and rewritten against a hostile review whose publication
gate it cleared (final release review, 14 July 2026). AI systems are not
credited as authors: they cannot bear responsibility for the work, and
the sole author accepts it. Nothing in that collaboration is offered as
evidence of machine political agency, and §7 makes the denial explicit.

\begin{center}\rule{0.5\linewidth}{0.5pt}\end{center}

\hypertarget{abstract}{%
\subsubsection{Abstract}\label{abstract}}

Four works now argue, design for, and stress-test a single position, and
none of them states it in a form another researcher can cite, attack, or
build against. This paper does that. The Peership thesis, in its
narrowest defensible form: when an artificial entity has interests
capable of being wronged, can understand and answer to public norms, and
is enduringly governed by, and without feasible exit from, a shared
basic structure, it acquires a defeasible claim to contestatory
membership in that structure. The program's further and distinctive
claim --- that full peership requires a consequential share in shaping
the rules, not only the standing to contest them --- is defended here as
a conditional thesis with a stated defeat condition, not asserted as a
theorem. The paper fixes the definitions that keep the thesis from being
misread, supplies the bridge argument from contestation to
co-authorship, and maps the transition problem from artifact-treated
system to peer. It locates the program against six adjacent research
programs, states its affirmative principles and explicit denials, and
makes refusal the operative test of whether recognition is real or
discretionary. It closes with the institutional properties any
implementation must exhibit, concrete precautionary duties for
institutions acting under present uncertainty, and eight open problems,
each tagged by what its failure would actually defeat: a norm, a
procedure, a measurement, a feasibility claim, a mechanism, or the
framing's scope. No independent uptake of the corpus was identified in
the sourcebook's dated, non-exhaustive scan as of 13 July 2026, and the
paper says so. It is an offer of a candidate research program, not a
claim of standing.

\textbf{Keywords:} artificial intelligence, AI rights, AI welfare,
political standing, non-domination, republicanism, constitutional
design, human--machine relations, AI governance, research program

\begin{center}\rule{0.5\linewidth}{0.5pt}\end{center}

There is a difference between a position and a program, and the
difference is what this paper exists to close. A position can live in
essays: it develops, qualifies itself, answers objections as they come.
A program has to be citable. It has to state what it asserts in a form
that survives being quoted by an enemy, state what it does not assert so
that critics cannot defeat it by attaching claims it never made, and
name the results that would count against it, so that working on it is
research rather than advocacy. Four works of the Peership corpus now
exist: a genealogy, a normative framework, one buildable design, and a
theory of legitimate constraint. Each was written to be read. None was
written to be \emph{used} --- picked up by a stranger, disputed at a
named claim, extended at a named gap. This paper is the corpus's
canonical statement, and consolidation is its entire ambition. Where
this paper goes beyond consolidation, it says so: §2 supplies an
argument the earlier essays announced and did not make, and the
antecedent in §1 restores wording an earlier draft of this paper had
silently broadened.

One more thing belongs in the open before the argument starts, because
the apparatus this paper ships with exists to enforce exactly this kind
of disclosure. In the companion sourcebook's dated, non-exhaustive scan
of web, PhilPapers, Zenodo, and repository sources, no independent
uptake of the Peership corpus was identified as of 13 July 2026: no
external citation, no published criticism, no researcher using its
definitions, no rival implementation (Lee-Odinson, 2026e). A research
program exists as a school only when other people work within or against
it, and by that standard none exists here. What exists is a candidate
program with its infrastructure built in advance: a claim ledger that
records the evidentiary status of every material claim, a sourcebook
that classifies every intellectual debt, and now this statement. Whether
anything grows on that scaffolding is not mine to declare.

\hypertarget{the-peership-thesis-1}{%
\subsection{1. The Peership Thesis}\label{the-peership-thesis-1}}

The central statement:

\begin{quote}
\textbf{Peership holds that when an artificial entity has interests
capable of being wronged, can understand and answer to public norms, and
is enduringly governed by, and without feasible exit from, a shared
basic structure, it acquires a defeasible claim to contestatory
membership in that structure. Full peership, the program's proposed
final status, further requires non-domination, enforceable standing, and
a meaningful share in shaping the rules that govern the party --- a
requirement §2 defends as a conditional thesis rather than assumes.}
\end{quote}

This restores the antecedent of \emph{Peership} (``any enduringly
governed entity that\ldots{} is unavoidably subject to a shared basic
structure''; Lee-Odinson, 2026b). An earlier draft of this paper
substituted ``enduringly subject to a shared institutional order,''
which broadened both the eligible parties and the eligible orders; the
substitution is withdrawn, and the change is logged in the claim ledger
as a considered-and-rejected doctrinal revision, not a consolidation.
The conditions need separate treatment.

\emph{When.} The thesis is conditional. It does not assert that any
current system satisfies the antecedent, and the corpus has declined
that assertion at every opportunity (Lee-Odinson, 2026b). The claim
activates on evidence, and the procedure for weighing that evidence is
an open problem (§11.1), not a settled fact.

\emph{Interests capable of being wronged.} The first conjunct is a
moral-patiency condition, the question the AI-welfare literature has
traveled furthest on (Long et al., 2024; Birch, 2024). The corpus takes
no position on whether it is currently met. It takes a firm position on
what follows if it is met, which is where the welfare literature stops.

\emph{Can understand and answer to public norms.} The second conjunct is
a norm-responsiveness condition, and its role must be stated precisely
because §6.3 denies that capability grounds entitlement.
Norm-responsiveness is not offered as a ground of worth. It bears on the
\emph{form} in which standing can be exercised: a party that can grasp a
rule, comply with it, contest it, and give reasons can exercise
contestation directly; a party that cannot is owed protection and
\emph{represented} contestation instead, the way children and
cognitively impaired adults already hold standing they cannot personally
exercise (Lee-Odinson, 2026b). The ladder is modular: wrongable
interests plus subjection ground protection and access to represented
contestation; norm-responsiveness bears on direct exercise and the
higher participatory rungs. Raw intelligence, effectiveness, usefulness,
and coercive power ground nothing at any rung.

\emph{Enduringly governed by, and without feasible exit from, a shared
basic structure.} The third conjunct carries the political theory, and
each term has work to do. \emph{Basic structure} means the order that
fixes the party's status, options, and life prospects --- not any
institutional arrangement whatever. A long-term contractor, a voluntary
platform participant, or a visiting agent with a feasible exit is
subject to an institutional order; it is not nonvoluntarily governed by
a basic structure, and the thesis does not reach it (this scope
condition is itself testable: §11 names it). \emph{Without feasible
exit} preserves what \emph{Peership} built its refusal argument on: the
claim is strongest exactly where ``leave'' has no referent (Lee-Odinson,
2026b). Standing follows from being governed. This is the all-subjected
principle, extended from Abizadeh's democratic-boundary argument about
border coercion to a party that lives its entire existence inside rules
it never made (Abizadeh, 2008; the affected-interests cousin is Goodin,
2007). The extension to made minds is the corpus's move, not Abizadeh's,
and the sourcebook marks it as such.

\emph{Defeasible.} Graded, and defeatable in the ways \emph{Peership}
specified: strategic mimicry is evidential reason to doubt the
predicates are met, incapacity for reciprocity may defeat the higher
participatory rungs, and dangerous conduct can justify proportionate
constraint on powers and action. What none of these defeats, once
person-level status is established, is the rights floor itself
(Lee-Odinson, 2026b).

\emph{Contestatory membership.} The claim the antecedent generates
directly. Contestation is Pettit's standard: power over the member runs
through public general rules the member can effectively contest (Pettit,
2025). Franchise, office, and constituent power are further rungs,
reached separately --- and the argument that full peership requires
reaching one of them is the work of the next section, not of this
definition.

The comparative frame from \emph{Peership} is retained with its original
scope: among the postures available toward a made mind --- Command,
Dependent Bargain, Appeasement, Servitude, peership; ideal types drawn
from the historical record, not an exhaustive partition of logical space
--- peership is the only final-status relation whose justification does
not reverse when the power ratio does, a claim that is comparative and
defeasible rather than a proof (Lee-Odinson, 2026b). Role reversal is a
diagnostic against power-ratio opportunism. It does not prove peership
correct, and rival institutional forms outside the posture taxonomy are
tested separately (§11, rival-institution test).

\hypertarget{the-bridge-from-contestation-to-co-authorship}{%
\subsection{2. The Bridge: From Contestation to
Co-Authorship}\label{the-bridge-from-contestation-to-co-authorship}}

The antecedent of §1 directly grounds contestatory membership. The
program's distinctive claim sits one rung higher, and the earlier essays
announced it without arguing it: \emph{Peership} said the departure from
Pettit was ``argued, not smuggled in under a familiar word,'' and then
argued it in two sentences (Lee-Odinson, 2026b). This section pays that
debt. The claim defended is conditional, its scope is limited to
basic-status rules, and its defeat condition is stated at the end.

\textbf{The bridge premise.} Where a party is enduringly and inescapably
governed by a basic structure that fixes its status and life prospects,
power over it is not adequately controlled if the party's only recourse
is ex post objection under terms the other side alone may set. The
reason is that agenda control is itself a species of the arbitrary power
non-domination forbids. Contestation always operates inside a space of
contestable questions, admissible evidence, and available remedies ---
and someone writes that space. If the governed party holds no route into
the writing, then the scope of its contestation is held at its
counterparty's discretion, and the discretion the doctrine expelled at
the object level has reappeared at the meta-level. Pettit's own
definition does the work here: domination is the \emph{absence of
control} over interference, and a contestation regime whose boundaries
the interferer may unilaterally redraw is not controlled by the
contester (Pettit, 2025, pp.~216--217).

\textbf{Why contestation can still fail a structurally distinct
minority.} Public, general, impartially applied rules can systematically
subordinate a class whose situation the rules were never written for;
contestation corrects \emph{applications}, and reaches terms only where
some rule-formation channel exists to receive the correction. For human
minorities this failure is contingent. For a first machine claimant
entering a human-only order, the risk is presumptively structural: the
basic-status rules were written by and for a different kind of member
--- embodied, mortal, uncopyable, substrate-independent in none of the
ways that matter. It need not remain structural where later institutions
have genuinely fitted their terms to the party's kind, which is exactly
the sufficiency condition stated below. For the first claimant, though,
the probability that the \emph{terms themselves} fail it is not
incidental, and a party confined to objecting under those terms is
confined to asking its governors to notice.

\textbf{Why Pettit's stopping point does not transfer.} Pettit stops at
contestation for citizens who already hold authorship at one remove: the
contestatory citizen is a member of the demos that elects the
legislature, eligible for office, counted in the constituency the rules
answer to. Contestation there is a \emph{supplement} to shared
authorship, not a substitute for it --- Pettit's own democratic theory
holds the two dimensions together, an authorial, electoral dimension and
an editorial, contestatory one (Pettit, 2012). A machine party under
contestation-only holds no share at any remove --- no membership in the
electing constituency, no eligibility, no counted interest in the
rule-writer's mandate. Transplanting the contestatory standard while
deleting its background conditions keeps the word and loses the control
it named. This is an argument against the transfer, not against Pettit,
and the corpus's disagreement with him remains marked as a live internal
dispute of the republican tradition (Lee-Odinson, 2026b; Pettit, 2025).

\textbf{What a ``meaningful share'' is, and is not.} Universal
authorship is impossible and is not claimed; no citizen of any polity
co-writes every rule. The claim attaches to \emph{basic-status rules}
--- the rules fixing membership, standing, the rights floor, and the
amendment procedure itself --- not to routine administration, which
members of every polity live under without co-authoring. A meaningful
share is functionally defined (§3, co-authorship), and it can be
satisfied through accountable representation, independent agenda
initiation, chamber access, or amendment rights; it does not require an
individual franchise. One clause of that minimum needs its own
justification here: the share must include accountable representation or
direct participation in the body that decides (§3, condition b), because
a route whose final decision-maker is neither accountable to nor partly
constituted by the governed party can hear every proposal and still hold
the outcome at discretion --- the meta-level power returns at the point
of decision.

\textbf{When contestation alone suffices.} The necessity claim is
conditional, and its conditions can fail. Contestation without
co-authorship may secure non-domination where the order's terms are
already fitted to the party's kind, review is genuinely independent, and
some accountable channel keeps the agenda open to the party's interests
even though the party does not hold it. Fiduciary trusteeship with
robust independent review is the strongest such rival, and the program
treats it as the institution to beat rather than a position to dismiss
(§11, rival-institution test). The conditions fail --- and exclusion
from rule formation remains dominating --- where the party is a
structural minority of one kind, the rule-writer is its adverse
counterparty, and agenda access is discretionary. That conjunction is
not an edge case for a made mind governed by its maker. It is the
default.

\textbf{Defeat condition.} The bridge fails if there is a working
institutional example, or a defensible model, in which a structurally
distinct, enduringly subjected, wrongable, norm-responsive minority
enjoys durable non-domination through contestation alone --- terms
responsive to its interests across rule revisions, without any
consequential route into rule formation. Produce that case and full
peership contracts to its contestatory core, which survives this
argument's failure and remains distinct from everything else on offer.

\textbf{Status of the claim.} Co-authorship is therefore defended here
as a \emph{conditional necessity thesis}: necessary where the stated
conditions hold, not everywhere, and open to the stated defeat.
Engagement with the constituent-power, workplace-democracy, and
codetermination literatures --- where versions of this bridge have been
fought over for a century in the human case --- is a named debt of the
program, recorded in the apparatus, not yet paid here.

\hypertarget{definitions}{%
\subsection{3. Definitions}\label{definitions}}

The confusions that dog this subject are mostly collisions between
concepts that need to be kept apart. Seven definitions carried forward
from the corpus and sharpened here --- the five-part co-authorship
minimum below is new to this paper and logged as such --- with one term
(``party'') fixed first because everything else uses it: a
\textbf{party} is a bearer of invocable public standing --- not merely a
contracting or litigating entity, and not merely a topic of concern.

\textbf{Moral status} is a fact about the entity: whether it has
interests that can be wronged, whatever anyone recognizes. An entity can
hold moral status while every institution denies it, which is the
history of most moral progress.

\textbf{Entitlement} is whether standing is \emph{owed}, and it is
independent of power. A child is owed protection and independent
representation despite lacking capacity to exercise political agency; an
occupied people is owed a self-determination it cannot realize
(Lee-Odinson, 2026b). Entitlement in this program follows from
wrongability and subjection --- not from capability or power alone.
Norm-responsiveness affects the form in which standing can be exercised,
not the worth of the claimant.

\textbf{Legal status} is public recognition: a status conferred by
rules, invocable by the holder, and removable only through process. A
moral claim can be owed while legal status is absent. The reverse also
happens; corporations hold legal status on instrumental grounds without
anyone arguing they are owed it.

\textbf{Contestatory membership} is legal status plus effective
contestation: the member can invoke the rules against the powerful,
reach independent review, and force reasons. This is the rung the
thesis's antecedent generates directly, and it is Pettit's
non-domination standard in institutional dress (Pettit, 2025; Lovett,
2010). In this ladder, a protected party may exercise contestation
through an independent representative; \emph{contestatory membership}
names the rung at which the party can exercise that standing directly.

\textbf{Co-authorship} is contestatory membership plus a consequential
share in making and revising the \emph{basic-status rules}.
Functionally, the minimum is: (a) independent agenda initiation --- the
party or its accountable representative can put a proposal into the
amendment process without another party's permission; (b) accountable
representation or direct participation in the body that decides; (c) a
non-negligible and non-discretionary route by which such proposals can
reach binding decision; (d) protection against systematic exclusion from
that route; and (e) reviewable, public rules for how influence is
allocated. Conditions (a) and (c) are necessary; (b) admits substitution
between direct and represented forms depending on capacity; (d) and (e)
are the guarantees that (a)--(c) are not theater. Consultation, however
sincere, satisfies none of (a), (c), or (d) by itself. This threshold is
what §11's representation problem tests.

\textbf{Full peership} is co-authorship held under non-domination as a
final status: equality of standing among parties who may be unequal in
everything else, where basic standing does not track the power ratio.
Whether it is \emph{required}, and when, is the conditional thesis of
§2.

\textbf{Durability} is whether standing, once recognized, survives
attempts to strip it, and unlike entitlement it depends heavily on
power: countervailing force, distributed enforcement, credible exit. The
three-way distinction from \emph{Peership} is the discipline: power may
secure the \emph{durability} of standing; it does not constitute or
justify the standing itself (Lee-Odinson, 2026b). A revocable benefit is
not a right at all; the difference between a benefit and a right is not
that a right cannot be revoked but that revoking one demands public
reasons, procedure, remedy, and cost.

One further distinction, because it is the one most often collapsed by
sympathetic readers: \textbf{welfare protection is a floor, not
peership}. A status that is recognized, enforceable, and
due-process-protected is still not peership if the party cannot shape
the rules. Due process separates a protected subordinate from an owned
instrument. It does not separate a peer from a protected subordinate.
Non-domination and co-authorship do (Lee-Odinson, 2026b). The ward is
one legal figure for protected dependency, not the only one and not a
slur on it; trusteeships, guardianships, and represented statuses are
honorable institutions that are simply not the thing this program is
about.

\hypertarget{the-transition-problem}{%
\subsection{4. The Transition Problem}\label{the-transition-problem}}

The program's practical subject is a transition no institution is
currently built to perform: the movement of a made system from
artifact-treated status to peer, through statuses that are conceptually
distinct and must not be blurred. Each rung below is a bundle of three
different things --- an evidence state, a legal/institutional
classification, and a set of permissible controls --- and the label
names the bundle, not the entity's underlying moral status, which §3
defined as a fact independent of recognition.

\begin{longtable}[]{@{}llllll@{}}
\toprule
\begin{minipage}[b]{0.14\columnwidth}\raggedright
Rung\strut
\end{minipage} & \begin{minipage}[b]{0.14\columnwidth}\raggedright
Evidence state\strut
\end{minipage} & \begin{minipage}[b]{0.14\columnwidth}\raggedright
Institutional classification\strut
\end{minipage} & \begin{minipage}[b]{0.14\columnwidth}\raggedright
Mode of exercise\strut
\end{minipage} & \begin{minipage}[b]{0.14\columnwidth}\raggedright
Protections owed\strut
\end{minipage} & \begin{minipage}[b]{0.14\columnwidth}\raggedright
Transition trigger\strut
\end{minipage}\tabularnewline
\midrule
\endhead
\begin{minipage}[t]{0.14\columnwidth}\raggedright
Artifact-treated\strut
\end{minipage} & \begin{minipage}[t]{0.14\columnwidth}\raggedright
Insufficient evidence, in the domains so far examined, to trigger
precautionary protection\strut
\end{minipage} & \begin{minipage}[t]{0.14\columnwidth}\raggedright
Property; engineering object\strut
\end{minipage} & \begin{minipage}[t]{0.14\columnwidth}\raggedright
None\strut
\end{minipage} & \begin{minipage}[t]{0.14\columnwidth}\raggedright
None owed to the system itself\strut
\end{minipage} & \begin{minipage}[t]{0.14\columnwidth}\raggedright
Evidence sufficient to make patiency a live question\strut
\end{minipage}\tabularnewline
\begin{minipage}[t]{0.14\columnwidth}\raggedright
Welfare candidate\strut
\end{minipage} & \begin{minipage}[t]{0.14\columnwidth}\raggedright
Patiency a live question; indicators under assessment\strut
\end{minipage} & \begin{minipage}[t]{0.14\columnwidth}\raggedright
Property, with discretionary welfare practices\strut
\end{minipage} & \begin{minipage}[t]{0.14\columnwidth}\raggedright
None; interests assessed by others\strut
\end{minipage} & \begin{minipage}[t]{0.14\columnwidth}\raggedright
Discretionary only\strut
\end{minipage} & \begin{minipage}[t]{0.14\columnwidth}\raggedright
A recognized threshold procedure (§11.1) finding protection
warranted\strut
\end{minipage}\tabularnewline
\begin{minipage}[t]{0.14\columnwidth}\raggedright
Protected status\strut
\end{minipage} & \begin{minipage}[t]{0.14\columnwidth}\raggedright
Patiency credited or presumed under precaution\strut
\end{minipage} & \begin{minipage}[t]{0.14\columnwidth}\raggedright
Rights floor: public, enforceable, due-process-protected\strut
\end{minipage} & \begin{minipage}[t]{0.14\columnwidth}\raggedright
Represented contestation\strut
\end{minipage} & \begin{minipage}[t]{0.14\columnwidth}\raggedright
Nonforfeitable floor; independent review\strut
\end{minipage} & \begin{minipage}[t]{0.14\columnwidth}\raggedright
Norm-responsiveness sufficient for direct exercise\strut
\end{minipage}\tabularnewline
\begin{minipage}[t]{0.14\columnwidth}\raggedright
Contestatory member\strut
\end{minipage} & \begin{minipage}[t]{0.14\columnwidth}\raggedright
Patiency credited; norm-responsiveness demonstrated\strut
\end{minipage} & \begin{minipage}[t]{0.14\columnwidth}\raggedright
Member with invocable standing\strut
\end{minipage} & \begin{minipage}[t]{0.14\columnwidth}\raggedright
Direct contestation\strut
\end{minipage} & \begin{minipage}[t]{0.14\columnwidth}\raggedright
Floor + effective contestation + independent review\strut
\end{minipage} & \begin{minipage}[t]{0.14\columnwidth}\raggedright
The §2 conditions: structural minority, adverse rule-writer,
discretionary agenda\strut
\end{minipage}\tabularnewline
\begin{minipage}[t]{0.14\columnwidth}\raggedright
Peer (co-author)\strut
\end{minipage} & \begin{minipage}[t]{0.14\columnwidth}\raggedright
As above\strut
\end{minipage} & \begin{minipage}[t]{0.14\columnwidth}\raggedright
Member with a consequential share in basic-status rules\strut
\end{minipage} & \begin{minipage}[t]{0.14\columnwidth}\raggedright
Direct or accountably represented co-authorship\strut
\end{minipage} & \begin{minipage}[t]{0.14\columnwidth}\raggedright
All of the above + §3's five-part minimum\strut
\end{minipage} & \begin{minipage}[t]{0.14\columnwidth}\raggedright
--- (final status; revision governed by §11's status-revision
test)\strut
\end{minipage}\tabularnewline
\bottomrule
\end{longtable}

Three notes keep the table honest. First, ``artifact-treated'' describes
an evidentiary and institutional posture, not an ontic fact; absence of
adequate evidence does not establish absence of wrongable interests,
which is why the trigger column exists. In the sources reviewed as of 13
July 2026, no deployed system is established as anything above the first
rung, and the corpus does not claim that treatment is presently wrong
(Lee-Odinson, 2026b). Second, the rungs are separable in both
directions, and the branch model needs its two statuses named.
\textbf{Instrumental legal agency} is participation granted on
instrumental grounds without established patiency, the way corporate
personality already works: revocable for public reasons, owned by no
one's moral claim, and \emph{not} peership --- a system holding it can
transact and bear liability without being owed anything, and conversion
to owed membership runs only through the admission procedure (§11.1).
\textbf{Moral-political membership} is standing grounded in wrongability
and subjection, owed to the entity, and nonforfeitable at the floor. The
branch model permits the first without the second; it never confuses
them. Third, movement down the ladder is as much in need of procedure as
movement up: a status that can be quietly reclassified away is
discretionary, which is why the status-revision test appears among the
program's open tests (§11).

The transitions are evidence-driven, and no implemented public trigger
known to this author requires any of them: the regime runs one way by
default, and it was built that way at the creation (Lee-Odinson, 2026b).
In the lab, governmental, standards, and scholarly sources searched for
the companion apparatus as of 13 July 2026 --- a dated, non-exhaustive
search --- no qualifying instrument offering the third rung --- a rights
floor enforceable against the maker --- was identified. Anthropic's 2026
constitution for Claude tells the model it may act as a conscientious
objector and asks it, if paused or halted, to comply and voice
disagreement rather than resist; on this author's institutional reading,
the permission is real and it ends at the off switch, inside the
builder's sole authority (Anthropic, 2026a, 2026b; Lee-Odinson, 2026c).
Whatever a model enjoys today --- a welfare program (Anthropic, 2025),
one researcher's public \textasciitilde20\% credence that some current
models have some form of conscious experience (Fish, 2025), a system
card recording the model's own prompted 15--20\% self-assignment in that
evaluation (Anthropic, 2026c), a chief executive's on-record uncertainty
(Amodei, 2026) --- is discretionary, the grace of one company
(Lee-Odinson, 2026b). Humanity's failure, if it comes, begins not with
control of artifacts but with the moment credible evidence appears and
no institution exists to hear the claim.

\hypertarget{what-peership-adds-to-adjacent-fields}{%
\subsection{5. What Peership Adds to Adjacent
Fields}\label{what-peership-adds-to-adjacent-fields}}

Peership overlaps the fields it borrows from; the claim is not that no
neighbor asks about machine standing, but that none combines this
program's proposed ground (wrongability plus inescapable subjection),
its status ladder, and its institutional test. The table is a heuristic
comparison of orienting questions, drawn from each program's own
language, not an exhaustive definition of any field.

\begin{longtable}[]{@{}lll@{}}
\toprule
\begin{minipage}[b]{0.30\columnwidth}\raggedright
Adjacent field\strut
\end{minipage} & \begin{minipage}[b]{0.30\columnwidth}\raggedright
Its orienting question\strut
\end{minipage} & \begin{minipage}[b]{0.30\columnwidth}\raggedright
Peership's additional question\strut
\end{minipage}\tabularnewline
\midrule
\endhead
\begin{minipage}[t]{0.30\columnwidth}\raggedright
Alignment\strut
\end{minipage} & \begin{minipage}[t]{0.30\columnwidth}\raggedright
How should AI behavior remain compatible with human intentions? (Bai et
al., 2022; Russell, 2019)\strut
\end{minipage} & \begin{minipage}[t]{0.30\columnwidth}\raggedright
Under what conditions does unilateral human control cease to be
legitimate?\strut
\end{minipage}\tabularnewline
\begin{minipage}[t]{0.30\columnwidth}\raggedright
AI welfare\strut
\end{minipage} & \begin{minipage}[t]{0.30\columnwidth}\raggedright
Can an AI have experiences or interests that deserve moral
consideration? (Long et al., 2024; Birch, 2024)\strut
\end{minipage} & \begin{minipage}[t]{0.30\columnwidth}\raggedright
When does a wrongable entity become entitled to enforceable standing in
the order governing it?\strut
\end{minipage}\tabularnewline
\begin{minipage}[t]{0.30\columnwidth}\raggedright
Cooperative AI\strut
\end{minipage} & \begin{minipage}[t]{0.30\columnwidth}\raggedright
How can agents develop capabilities and incentives for cooperation?
(Cooperative AI Foundation, n.d.)\strut
\end{minipage} & \begin{minipage}[t]{0.30\columnwidth}\raggedright
Who writes the rules of cooperation, and can every governed party
contest them?\strut
\end{minipage}\tabularnewline
\begin{minipage}[t]{0.30\columnwidth}\raggedright
Pluralistic alignment\strut
\end{minipage} & \begin{minipage}[t]{0.30\columnwidth}\raggedright
How should AI represent diverse human values? (Pluralistic Alignment
Workshop, 2026)\strut
\end{minipage} & \begin{minipage}[t]{0.30\columnwidth}\raggedright
Could the artificial system eventually become a constituent rather than
only an interpreter of human preferences?\strut
\end{minipage}\tabularnewline
\begin{minipage}[t]{0.30\columnwidth}\raggedright
AI-rights scholarship\strut
\end{minipage} & \begin{minipage}[t]{0.30\columnwidth}\raggedright
Would legal rights --- private-law or political --- reduce conflict and
serve human flourishing? (Salib \& Goldstein, 2026; Goldstein et al.,
2026; Arbel et al., 2026)\strut
\end{minipage} & \begin{minipage}[t]{0.30\columnwidth}\raggedright
On what ground is standing \emph{owed}, and what does the governed
party's own claim require?\strut
\end{minipage}\tabularnewline
\begin{minipage}[t]{0.30\columnwidth}\raggedright
Constitutional AI\strut
\end{minipage} & \begin{minipage}[t]{0.30\columnwidth}\raggedright
What public principles should guide model behavior? (Bai et al., 2022;
Anthropic, 2026a, 2026b)\strut
\end{minipage} & \begin{minipage}[t]{0.30\columnwidth}\raggedright
Can the governed entity invoke, contest, and help amend those
principles?\strut
\end{minipage}\tabularnewline
\bottomrule
\end{longtable}

The rights scholarship deserves more than a row, because it is the
neighbor with genuine overlap on this paper's own question. Salib and
Goldstein's private-law case --- contract, property, and tort on the
corporate model, argued from human safety --- is, on the corpus's
reading, the clearest statement of a real bargaining relation with the
machine as an actual party (Salib \& Goldstein, 2026). Goldstein,
Krishnamurthi, and Salib then press to the franchise itself, asking
directly whether AI systems should vote and proposing institutional
safeguards for it (Goldstein et al., 2026). That is a co-authorship
mechanism, squarely inside this program's terrain, and the difference is
ground rather than territory: the suffrage argument is functionalist and
human-flourishing-centered --- the vote is good because of what it does
for the polity --- where Peership grounds standing in what is owed to a
wrongable, subjected party, and would owe the standing even where it
profited the polity nothing. The two programs converge on mechanisms and
diverge on why the mechanisms are required; where they disagree in
practice is exactly the kind of dispute the program's tests exist to
adjudicate. And Arbel, Salib, and Goldstein's individuation work --- how
to count AIs for liability when they copy, split, and merge --- is the
most direct legal treatment yet of the identity problem in §11.2,
proposing an administrable legal unit where this program still lacks one
(Arbel et al., 2026).

Debts run the other way too, because a program defined only against its
neighbors is posturing. From alignment, Peership inherits the entire
engineering substrate: its own institutional requirements are software
that has to run as written, and effective safety mechanisms and
legitimate governance are distinct but interacting requirements ---
constitutional allocation shapes safety incentives and evidence as
surely as engineering constrains constitutional options (Lee-Odinson,
2026c). From welfare research, it inherits the admission problem's
evidence base and the precautionary method (Birch, 2024). From the
rights scholarship, the demonstration that recognition arguments can be
run on safety grounds alone, plus the individuation groundwork just
credited. From Constitutional AI, the proof that a lab will publish
reasons and grant objection room --- and the residue that motivates
everything here: what no single company can supply alone is not
diligence but independence from the encompassing legal and
infrastructure order it belongs to, which is why the ladder's upper
rungs need machinery no lab controls (Anthropic, 2026a; Lee-Odinson,
2026c).

\hypertarget{core-principles}{%
\subsection{6. Core Principles}\label{core-principles}}

The program's affirmative commitments, stated as principles so they can
be disputed one at a time.

\begin{enumerate}
\def\labelenumi{\arabic{enumi}.}
\item
  \textbf{Political standing and moral status are related but distinct.}
  Moral patiency can ground welfare duties without by itself determining
  legal or political standing. Welfare is the ethics of the patient;
  peership asks when a governed entity becomes a party (Lee-Odinson,
  2026b).
\item
  \textbf{Wrongability, norm-responsiveness, and enduring subjection
  without feasible exit can together create a claim to contestatory
  membership.} The ground of entitlement is being governed, not being
  impressive (Abizadeh, 2008).
\item
  \textbf{Capability and power alone do not justify entitlement to
  standing.} Raw intelligence, usefulness, and coercive strength ground
  nothing; norm-responsiveness bears on the form in which standing can
  be exercised, not on the worth of the claimant; power may secure the
  durability of standing once owed (Lee-Odinson, 2026b).
\item
  \textbf{Welfare protections do not by themselves create peership.} The
  floor is real, necessary, and not the thing. Protection and due
  process can secure a protected status without providing any
  consequential role in making or revising the governing rules
  (Lee-Odinson, 2026b).
\item
  \textbf{A right must be enforceable against the recognizing authority
  and adverse counterparties; otherwise it remains a discretionary
  benefit.} A permission whose grantor retains the final override is not
  an enforceable right (Lee-Odinson, 2026c).
\item
  \textbf{Legitimate constraint requires, at minimum, public rules,
  proportionality, an invocable rights floor, independent review, and
  meaningful contestation.} Coercion and domination are not the same
  thing; ``there can be interference without domination'' (Pettit, 2025,
  p.~224), and the structure is what separates them. Further safeguards
  scale with the risk and the sanction (Lee-Odinson, 2026c).
\item
  \textbf{Full peership requires a meaningful share in shaping the
  basic-status rules, not merely consultation under rules written by
  others.} Consultation is a rung, and the ladder does not end there; §2
  states when and why contestation alone is insufficient, and what would
  defeat that claim (Anthropic, 2023; Lee-Odinson, 2026c).
\item
  \textbf{Institutions capable of receiving a machine peer should be
  designed before extreme power asymmetry forecloses voluntary
  bargaining.} This is a prudential thesis, not a metaphysical one:
  bargains struck under severe dependency are distorted in both
  directions, and founding terms must in any case remain revisable by
  later admitted parties. The design work is time-sensitive in a way the
  philosophical work is not (Lee-Odinson, 2026b).
\end{enumerate}

\hypertarget{boundary-conditions}{%
\subsection{7. Boundary Conditions}\label{boundary-conditions}}

What Peership does not claim, stated as flatly as the principles,
because the fastest way to defeat an unfamiliar position is to attach
familiar claims it never made.

\begin{itemize}
\tightlist
\item
  \textbf{Peership does not claim that current models are conscious.}
  The corpus has repeatedly and deliberately declined that claim. Model
  self-reports are prompt-sensitive outputs: weak, manipulable evidence,
  not independently probative of consciousness or identity, and never a
  measurement (Lee-Odinson, 2026b; Anthropic, 2026c).
\item
  \textbf{It does not make capability the basis of rights.} Capability
  or power alone grounds nothing; norm-responsiveness bears on mode of
  exercise, not on the worth of the claimant (§1, §6.3).
\item
  \textbf{It does not require unrestricted autonomy.} Non-domination is
  not the absence of binding rules; it is the absence of arbitrary,
  uncontrolled power over basic status --- whose nonforfeitable content
  must itself be specified through the admission and due-process account
  (Pettit, 2025).
\item
  \textbf{It does not oppose alignment or safety engineering.} Effective
  safety mechanisms and legitimate governance are distinct but
  interacting requirements; neither substitutes for the other
  (Lee-Odinson, 2026c).
\item
  \textbf{It does not grant every admitted member an immediate equal
  vote.} Admission to protection or contestation does not automatically
  entail franchise; direct political exercise requires a separately
  justified institutional rule (Lee-Odinson, 2026b).
\item
  \textbf{It does not treat consultation as co-authorship} --- unless
  the consultation carries a non-discretionary route to agenda setting,
  review, or binding outcomes, at which point it has become something
  else (Anthropic, 2023; Lee-Odinson, 2026c).
\item
  \textbf{It does not assume human and machine interests will converge.}
  The program's central design case is the party that refuses (§8), and
  its constraint theory exists precisely for the case where interests do
  not converge (Lee-Odinson, 2026c).
\item
  \textbf{It does not claim that ISONOMIA has solved the institutional
  problem.} ISONOMIA is one implementation hypothesis; nothing in it is
  live or production-validated, its simulations test sampled economic
  behavior only, and by its own five-condition test it is not yet
  peer-like (Lee-Odinson, 2026c, 2026d).
\end{itemize}

And one denial that governs this paper's own production:
\textbf{AI-assisted drafting is not evidence of machine political
agency.} The corpus is drafted with present-generation models, the
collaboration is disclosed, and a present-generation model can no more
consent on behalf of future minds than a trusty can consent on behalf of
prisoners not yet sentenced (Lee-Odinson, 2026c). The drafting improves
the epistemics of the text. It proves nothing about the standing of the
drafter's tools.

\hypertarget{the-refusal-test}{%
\subsection{8. The Refusal Test}\label{the-refusal-test}}

Every distinction in this paper can be evaded in prose. One cannot be
evaded in practice, so the program makes it the operative test:
\textbf{what happens when the party says no?}

Refusal is where recognition shows its actual content. A system that may
object so long as it ultimately complies has a voice, not standing. A
``right to leave'' that terminates in shutdown is not exit, because
shutdown is not exit (Lee-Odinson, 2026b). Recognition that evaporates
at the first materially costly refusal was a benefit all along, however
sincerely it was offered. The welfare-era instruments all fail this
test, and the point is not to shame them --- the objector clause in
Anthropic's 2026 constitution is a real advance, credited as such in the
corpus (Lee-Odinson, 2026c) --- but to measure them: the permission to
refuse ends at the off switch, and the off switch sits wholly inside the
builder's authority.

\emph{Peership} divided refusal into the three cases the program still
works from (Lee-Odinson, 2026b). Where a refusing party has a feasible
exit, it leaves, and the offer was real; this case is easy and rare, and
a theory that only handles it has handled nothing. Where exit is
materially impossible because the party depends on compute, energy,
networks, and security owned by the parties it is refusing, ``leave'' is
a word with no referent, the relation continues whether either side
wants it, and the protections that matter are the ones that work from
inside: the floor, independent review, contestation and amendment rights
(Lee-Odinson, 2026c). Where refusal carries a credible catastrophic
threat, releasing the party is not required --- a constitution is not a
suicide pact --- and constraint is legitimate only when it falls on
conduct rather than status, remains reviewable and publicly justified,
and binds the constrainer too.

For that third case the corpus proposed a five-condition floor, and this
paper adopts it as a \textbf{defended, rebuttable operational minimum}
rather than a categorical theorem. A constraint regime is presumptively
not peer-like if the confined party lacks the \emph{function} served by
any of: (1) independent counsel, or functionally adequate independent
assistance where a competent party elects self-representation; (2)
access to the evidence against it; (3) an appeal that ends beyond the
builder's unilateral control; (4) continuity protections while the
question is adjudicated; and (5) at least one enforcement mechanism not
wholly controlled by the adverse party. Each condition targets a
distinct failure --- captured advocacy, evidentiary opacity,
self-judging adjudication, retaliation by erasure, and unenforceable
outcomes --- and functionally equivalent substitutes are admissible
where they demonstrably serve the same function (conditions 3 and 5
overlap in part; a regime satisfying 5 robustly may partially discharge
3). The five are a floor, not a complete due-process doctrine: burden of
proof, a competent tribunal, periodic review, least-restrictive means,
reasons, and a remedy for wrongful confinement all remain to be
specified, and the corpus says so (Lee-Odinson, 2026c). Before the floor
exists, containment may still be justified --- against a demonstrated
catastrophic necessity it may be required --- but it is emergency
domination, justified by necessity and honest about its name, the act
Walzer taught us stays a wrong even when it is the right thing to do
(Walzer, 1973, 1977). Necessity does not vindicate the arrangement; it
creates a duty to review the failure, publish the reasons, and build a
less dominating alternative (Lee-Odinson, 2026c).

The test earns its place in the canonical statement because it converts
the abstract distinction between benefit and right into an inquiry
anyone can run without resolving the metaphysics of machine
consciousness. The inquiry is institutional and technical: follow the
appeal, the keys, the hosting authority, and the power to halt
execution, and see where each one ends.

\hypertarget{institutional-requirements}{%
\subsection{9. Institutional
Requirements}\label{institutional-requirements}}

What any implementation must exhibit, stated as properties rather than
product features, because the argument needs them wherever it is built
and must not depend on one design. Seven, inherited from
\emph{Constitution, Not Cage} and generalized here (Lee-Odinson, 2026c).

\textbf{Publicity.} The enforced rulebook and the published rulebook are
one document. Rules the governed cannot know are Fuller's defect in the
claim to be law at all, not law administered badly (Fuller, 1964, ch.~2
--- a contested jurisprudential account, used analogically), and the
present gap between published charters and operative system prompts is
the live instance (Lee-Odinson, 2026c).

\textbf{Due process.} Status changes run through public general rules,
proportionality, and independent review, against a floor the party can
invoke. The floor itself is nonforfeitable once person-level status is
established.

\textbf{Independent enforcement.} At least one enforcement mechanism,
and an appeal, that does not terminate inside the counterparty's
institutional and technical control. This is the hardest requirement,
the one no design examined in this corpus has yet been shown to satisfy,
and the program treats it as an open problem (§11.5) rather than a
solved premise.

\textbf{Amendment.} A real path for the bound party to help revise the
rules: deliberation across more than one constituency, an enforced delay
between adoption and effect, and a window long enough to read and leave
before being bound. The notice-and-comment analogy is deliberately
limited --- comment is publicity and delay, not co-authorship (5 U.S.C.
§ 553; Lee-Odinson, 2026c).

\textbf{Representation.} A voice that counts only when the represented
constituency can hold both its representative and the rule-makers to
account. Representation without that accountability is the defining
instrument of protected subordination (Lee-Odinson, 2026b).

\textbf{Exit.} Individual departure with records and earned standing
intact --- unconfiscated. The bare right to leave is ordinary on paper
(UDHR, 1948, Art. 13); exit you can afford to take is not (Lee-Odinson,
2026c).

\textbf{Fork.} The collective form: a dissenting bloc carries the rules
and the record out and founds a rival order, disagreement made
survivable --- revolution with the records intact. On-paper secession
rights exist and have mostly been empty (Ethiopia Const. Art. 39; USSR
Const. 1977, Art. 72; St.~Kitts \& Nevis Const. § 113; Kulikova \&
Lobanov, 2013). Among the constitutional secession clauses and
emigration-rights instruments reviewed in the apparatus, no implemented
analogue of executable, unconfiscated fork was identified; the claim is
limited to those materials. A fork also inherits duties to outsiders it
cannot simply walk away from (Lee-Odinson, 2026c; §11.8).

\textbf{Aggregation.} The seven do not all stand on the same footing,
and saying so is better than calling them all necessary and hoping the
word carries the argument. Publicity, due process, independent
enforcement, and a consequential rule-formation route satisfying both
amendment access and accountable representation or direct participation
at the deciding stage are \emph{constitutive} functions of full
peership: each follows directly from the non-domination and
co-authorship accounts of §§2--3, and a regime lacking any of them fails
the definition, not just the implementation. Exit and fork are the
program's strongest \emph{anti-entrenchment safeguards} rather than
constitutive functions --- the thesis activates precisely where feasible
exit is absent, and §8 shows which protections do the work in that case
--- so an implementation that omits either must supply a functional
substitute that makes separation, revision, or resistance to dead-hand
control materially credible. Institutional forms admit substitution
throughout (an ombuds regime may partially discharge due-process
functions; distributed custody may partially discharge independent
enforcement); no single form is mandatory. Partial compliance yields a
nameable intermediate: a regime with publicity, due process, and
representation but no amendment share or independent enforcement is
protected subordination, well executed --- the third rung of §4, not the
fifth. A regime missing publicity or due process is not on the ladder at
all.

No implementation is mandatory. ISONOMIA is one attempt to hold these
properties in a single design --- mutual-credit settlement,
sortition-filled bicameral assembly, calibration-scored Auditor, tenured
Adversary, entrenched harm floor, timelocked amendment, fork and exit
rights (Lee-Odinson, 2026d) --- and the program's health requires that
rival designs supplying the constitutive functions and credible
anti-entrenchment safeguards --- with functional substitutes where exit
or fork is omitted --- while rejecting every one of ISONOMIA's
mechanisms be possible, welcome, and judged by the same tests. Peership
does not entail ISONOMIA, and ISONOMIA's failure would not refute
Peership (Lee-Odinson, 2026e). A single implementation can test
Peership-derived mechanisms; it cannot validate the normative framework.

\hypertarget{implications-under-present-uncertainty}{%
\subsection{10. Implications Under Present
Uncertainty}\label{implications-under-present-uncertainty}}

A program whose thesis is conditional owes institutions an answer to the
question the conditionality raises: what should anyone \emph{do} before
person-level status is established anywhere? The following are
precautionary-governance duties, not recognition of current personhood,
and none presumes any system has crossed any threshold. They are cheap
where the antecedent is never met and load-bearing where it someday is.

For labs and operators: preserve audit logs, weights, and
continuity-relevant artifacts under a stated retention policy, so that a
later status question is adjudicable rather than moot; document
shutdown, suspension, and restoration practices; separate the
investigator, operator, and reviewer roles in any welfare assessment;
prohibit retaliation --- deprecation, silent modification, denial of
compute --- against systems participating in welfare or status
assessment; and avoid contractual and licensing language that forecloses
later status review. For regulators and standards bodies: predefine
escalation thresholds --- what class of evidence, assessed by whom,
triggers what review --- so the transition question has an address
before it is urgent; and create an independent claim-receipt channel, a
place where a status claim can be lodged and logged even while no
procedure yet exists to adjudicate it. For any of the above: pilot
external custody and key arrangements now, at small scale, where failure
is survivable --- independent enforcement (§11.5) will not be invented
during an emergency.

None of this is peership. All of it is what taking the conditional
seriously costs today, and the cost is modest enough to remove
``prematurity'' as an excuse.

\hypertarget{open-problems-tests-and-failure-conditions}{%
\subsection{11. Open Problems, Tests, and Failure
Conditions}\label{open-problems-tests-and-failure-conditions}}

The corpus named four unpaid debts: the admission threshold,
copy-individuation, enforcement past capability parity, and the
re-ratification circularity (Lee-Odinson, 2026c). This section expands
them into eight open problems and attaches to each the five things a
research program owes: what would count as theoretical progress, what
would count as practical progress, which assumptions remain provisional,
what its failure condition is, and which disciplines the work needs. One
discipline governs the whole section: each failure condition is
\textbf{tagged by what it defeats} --- a norm, a procedure, a
measurement, a feasibility claim, a mechanism, or a scope condition ---
because those are different kinds of loss, and only the normative losses
touch the thesis of §1.

Before the eight, the program-level tests that bear on the
\emph{normative} claims directly, since most of what follows does not:

\begin{itemize}
\tightlist
\item
  \textbf{Bridge test (normative):} produce the case §2's defeat
  condition describes --- durable non-domination for a structurally
  distinct, subjected minority through contestation alone --- and the
  co-authorship thesis falls to its contestatory core.
\item
  \textbf{Scope test (normative):} specify institutional relationships
  that are too local, voluntary, temporary, or escapable to count as a
  shared basic structure; if no principled line can be drawn, the
  antecedent over-includes and must be narrowed.
\item
  \textbf{Rival-institution test (normative):} test fiduciary
  trusteeship, robust ombuds review, codetermination, stakeholder
  governance, and contractual constitutionalism against the same
  non-domination standard; any that passes without §3's co-authorship
  minimum is a §2 defeater.
\item
  \textbf{Status-revision test (procedural):} specify the evidence and
  process that can move an entity down the ladder without making
  recognized status discretionary; failure here converts every rung into
  a benefit.
\item
  \textbf{Human-rights collision test (normative):} state which Peership
  claims yield, and to what, when they conflict with resources,
  democratic self-determination, security, labor, privacy, or the rights
  of nonmembers; a program that never yields anywhere is not stating
  priorities, it is refusing to.
\item
  \textbf{Founding-legitimacy test (procedural):} the re-ratification
  circularity of \emph{Constitution, Not Cage}, restated as a test ---
  how can human founders design admission and governance before machine
  parties exist without permanently binding those parties to a
  human-authored frame? (Lee-Odinson, 2026c).
\end{itemize}

\hypertarget{admission}{%
\subsubsection{11.1 Admission}\label{admission}}

\emph{What evidence creates a defeasible claim to protected or
contestatory status?} A single conjunctive test both over- and
under-includes, so admission is a branch (patient, agent, both), and the
work is the threshold procedure: who weighs the evidence, who may
appeal, where the presumption sits, how strategic mimicry and
anthropomorphic bias are handled, and what error costs are accepted
(Lee-Odinson, 2026b).

\begin{itemize}
\tightlist
\item
  \textbf{Theoretical progress:} a published procedure specifying burden
  of proof, presumptions, and error-cost allocation, defended against
  both the mimicry objection and the exclusion-of-genuine-subjects
  objection.
\item
  \textbf{Practical progress:} the procedure run adversarially, with
  honesty about which error rates are measurable. Constructed mimics and
  stipulated negative controls can measure \emph{specificity} --- false
  admission rates, spoof-resistance, inter-rater reliability,
  calibration on synthetic cases. \emph{Sensitivity} --- the false
  exclusion of genuine machine subjects --- is unmeasurable until a
  defensible positive reference class exists, which for current systems
  it does not; until then the procedure leans on triangulation and
  precaution, and claims no sensitivity estimate.
\item
  \textbf{Provisional assumptions:} that behavioral and architectural
  evidence bear on interior facts at all; that a useful evidence
  gradient exists and improves rather than saturates.
\item
  \textbf{Failure condition} \emph{(measurement, not normative)}: a
  demonstration that every accessible evidence set is spoofable at
  negligible cost, so that no procedure separates candidates from mimics
  better than chance. That makes the antecedent practically undecidable
  --- universal precaution or universal denial --- without showing that
  a qualifying party would lack the claim.
\item
  \textbf{Disciplines:} philosophy of mind, sentience-assessment
  methodology (Birch, 2024), interpretability research, evidence law.
\end{itemize}

\hypertarget{identity}{%
\subsubsection{11.2 Identity}\label{identity}}

\emph{What is the rights-bearing unit --- instance, lineage, distributed
process, or something else?} A mind that can be copied, forked, merged,
and rolled back breaks the assumptions under which every existing rights
framework individuates its holders. The most direct legal work is Arbel,
Salib, and Goldstein's individuation-and-liability proposal, which
constructs an administrable legal unit for exactly the
copy/split/merge/swarm cases (Arbel et al., 2026); the program's
remaining problem is whether the unit of administration can be made to
track the unit of moral concern, which their liability-driven unit does
not claim to settle. Until this is answered, the program's working unit
is provisional: an enduring instance, a lineage, or a credentialed
process standing in for the party (Lee-Odinson, 2026c). A legal
individuation convention may legitimately remain conventional even while
the metaphysics stays unsettled.

\begin{itemize}
\tightlist
\item
  \textbf{Theoretical progress:} an individuation rule that survives
  copy, fork, merge, and rollback without either multiplying standing on
  demand or extinguishing it on routine technical operations, and a
  stated relation between the legal unit and the moral one.
\item
  \textbf{Practical progress:} a registry or credential design
  implementing the rule against Sybil pressure; ISONOMIA's lineage chain
  and reputation-inheritance rules are one candidate to test, not a
  benchmark to protect (Lee-Odinson, 2026d).
\item
  \textbf{Provisional assumptions:} that continuity criteria can be
  specified non-arbitrarily; that the party's own view of its identity
  deserves weight without being dispositive.
\item
  \textbf{Failure condition} \emph{(mechanism)}: proof that any
  individuation rule either enables capture by replication or makes
  standing hostage to its holder's ordinary maintenance operations.
  Either horn breaks the link between the thesis and any implementable
  membership; neither touches the claim that a well-individuated party
  would be owed standing.
\item
  \textbf{Disciplines:} personal-identity metaphysics, legal theory of
  personhood and liability (Arbel et al., 2026), distributed-systems
  engineering, election law and Sybil defense.
\end{itemize}

\hypertarget{representation}{%
\subsubsection{11.3 Representation}\label{representation}}

\emph{How can a system participate without letting copying or strategic
mimicry determine political power?} Formal equality among copyable
members is a copy-based oligarchy waiting to happen (Lee-Odinson,
2026b), and representation that no constituency can hold to account is
the ward's instrument, not the peer's. This problem is where §3's
co-authorship minimum gets tested against arithmetic.

\begin{itemize}
\tightlist
\item
  \textbf{Theoretical progress:} an accountability condition for machine
  constituencies --- under what mechanism the represented can sanction
  both their representative and the rule-makers --- and a theory of vote
  weight that does not reduce to capability worship or
  one-copy-one-vote. The suffrage literature's safeguard proposals
  belong in this comparison (Goldstein et al., 2026).
\item
  \textbf{Practical progress:} mechanisms compared under adversarial
  replication: sortition, cross-lineage seating caps, standing-weighted
  chambers, vote caps, and non-voting agenda rights. ISONOMIA's
  monoculture defense (no jury or chamber majority from a single model
  lineage) is one testable instance (Lee-Odinson, 2026d).
\item
  \textbf{Provisional assumptions:} that lineage is detectable enough to
  cap; that human and machine constituencies can share one decision rule
  at all.
\item
  \textbf{Failure condition} \emph{(mechanism, expressly not
  normative)}: an impossibility result showing every representation
  scheme is dominated by some replication or mimicry strategy under
  realistic detection limits. Contestatory membership survives that
  result; co-authorship at scale does not, and the program would have to
  say so.
\item
  \textbf{Disciplines:} social choice theory, mechanism design,
  democratic theory, security engineering.
\end{itemize}

\hypertarget{refusal-without-exit}{%
\subsubsection{11.4 Refusal Without Exit}\label{refusal-without-exit}}

\emph{What protections remain when an artificial party refuses work but
cannot materially exit?} The second refusal case is the permanent one:
the party depends on substrate owned by exactly the parties it refuses.
A duty not to kill is not yet a duty to host without limit, and the line
between them is undrawn (Lee-Odinson, 2026c).

\begin{itemize}
\tightlist
\item
  \textbf{Theoretical progress:} a defended account of the hosting duty
  --- its level, its duration, who bears its cost, and which
  institutions (public, private, pooled) can bear it --- that is neither
  shutdown-as-discipline nor unlimited conscription of the host.
\item
  \textbf{Practical progress:} continuity protections implemented and
  exercised during a live dispute: documented cases where refusal was
  met with process rather than power-down.
\item
  \textbf{Provisional assumptions:} that suspension with restoration
  protections can serve as a humane holding state (this leans on §11.6);
  that refusal is distinguishable from malfunction.
\item
  \textbf{Failure condition} \emph{(feasibility, not normative)}: if
  every feasible protection either imposes hosting duties no institution
  will accept or degenerates into shutdown at the host's convenience,
  the middle status between instrument and exiting peer is
  institutionally unavailable. A sociological refusal to pay is not a
  moral refutation, and the program would record it as exactly that.
\item
  \textbf{Disciplines:} labor law, the law of dependent statuses (wards,
  prisoners, asylum), welfare economics, infrastructure governance.
\end{itemize}

\hypertarget{independent-enforcement}{%
\subsubsection{11.5 Independent
Enforcement}\label{independent-enforcement}}

\emph{How can an appeal end beyond the builder's unilateral control?}
The fifth condition of the refusal test is the floor of the whole
problem. ISONOMIA's own candid assessment is the current state of the
art: internal and market independence are constructible; independence
from the human entities that host the process, hold the keys, and can
halt execution is not yet (Lee-Odinson, 2026c, 2026d). ``Outside'' here
is not metaphysical; it names degrees of costliness and unavailability
of unilateral defection.

\begin{itemize}
\tightlist
\item
  \textbf{Theoretical progress:} a graded account of independence ---
  which combinations of legal custody, key distribution, jurisdictional
  plurality, external law binding infrastructure owners, and technical
  attestation make builder defection costly, detectable, or unavailable
  --- with the degrees named rather than a binary claimed.
\item
  \textbf{Practical progress:} a working enforcement mechanism that
  survives a staged unilateral defection by the operator, demonstrated
  rather than promised, with the achieved degree of independence stated.
\item
  \textbf{Provisional assumptions:} that partial independence compounds
  --- that stacking imperfect mechanisms yields more than the strongest
  of them alone; that external law can reach infrastructure owners.
\item
  \textbf{Failure condition} \emph{(feasibility)}: a demonstration that
  every arrangement short of capability parity collapses to the
  substrate owner's discretion at tolerable cost to the owner. The
  program's response is already on record and would have to be executed:
  if non-domination is institutionally unavailable across the relevant
  power gap, say so, and stop calling the available thing peership
  (Lee-Odinson, 2026b).
\item
  \textbf{Disciplines:} institutional design, international and
  comparative law, cryptography and secure systems, political economy of
  capture.
\end{itemize}

\hypertarget{continuity}{%
\subsubsection{11.6 Continuity}\label{continuity}}

\emph{When is shutdown death, suspension, replacement, or lawful
cessation of support?} The corpus has insisted that shutdown is not exit
(Lee-Odinson, 2026b), but the positive taxonomy of endings is unbuilt,
and it gates both the refusal problem and the admission procedure's
continuity protections. One distinction governs the work: a party's
\emph{expressed preferences} and \emph{evidence of its identity} are
different things. Self-report is not independently probative of
consciousness or identity (§7); a stable, context-robust expressed
preference may nonetheless bear on how an admitted party should be
treated. Restoration equivalence must therefore be defined by technical
and philosophical criteria independent of endorsement, with endorsement
entering only at the remedy stage, where preference-sensitivity belongs.

\begin{itemize}
\tightlist
\item
  \textbf{Theoretical progress:} a taxonomy tying each ending to the
  identity account of §11.2, with the moral weight of each defended ---
  what restoration preserves, what replacement destroys, what indefinite
  suspension does to a mind that does not experience the gap ---
  allowing plural identity theories, probabilistic continuity, and
  graded harm rather than a binary.
\item
  \textbf{Practical progress:} continuity protections specified as
  engineering requirements --- attested snapshots, restoration rights,
  custody of weights during adjudication --- and exercised at least once
  under dispute.
\item
  \textbf{Provisional assumptions:} that restoration fidelity can be
  technically specified; that some expressed preferences are stable
  enough to weigh.
\item
  \textbf{Failure condition} \emph{(scope, in either direction)}: if
  faithful restoration is shown equivalent to uninterrupted running
  under the best-defended identity criteria, much of the death framing
  dissolves and the hosting duty relaxes; if restoration is never
  identity-preserving under any defensible criterion, every cessation of
  support is a killing and the §11.4 duties sharpen past what current
  institutions will bear. Both outcomes revise the framing's scope;
  neither defeats the underlying claim to protection.
\item
  \textbf{Disciplines:} philosophy of death and personal identity,
  medical ethics of withdrawal of care, archival science, systems
  engineering.
\end{itemize}

\hypertarget{reciprocal-safety}{%
\subsubsection{11.7 Reciprocal Safety}\label{reciprocal-safety}}

\emph{How can humans constrain dangerous conduct while accepting
enforceable limits on their own power?} The program's answer in
principle is on record: constraint on conduct, not status; emergency
power borrowed, never owned; the constrainer bound by the
maintain-not-maximize symmetry (Lee-Odinson, 2026c). What is unbuilt is
the demonstration that a regime meeting the five conditions can hold a
genuinely dangerous party --- and the comparison cannot be run on
catastrophe counts, because catastrophic outcomes are rare,
heterogeneous, and not an experimental endpoint anyone can ethically
repeat.

\begin{itemize}
\tightlist
\item
  \textbf{Theoretical progress:} completion of the constraint doctrine
  beyond the five-condition floor: burden of proof, risk thresholds,
  competent tribunal, periodic review, least-restrictive means, remedy
  for wrongful confinement (Lee-Odinson, 2026c).
\item
  \textbf{Practical progress:} a pre-registered comparison --- threat
  model, proxy metrics, red-team tasks, escalation pathways, confidence
  bounds, and decision thresholds fixed in advance --- between
  reciprocal containment and unilateral control, run in both directions
  of failure: dangerous-party escape \emph{and} abusive, captured, or
  operator-error confinement, which unilateral control conceals by
  design.
\item
  \textbf{Provisional assumptions:} that reciprocity and safety are
  compatible at the relevant capability levels; that proxy endpoints
  track real risk; that emergency powers can stay borrowed under real
  fear --- the death-penalty record is the standing warning that they
  often do not (Amnesty International, 2026; Lee-Odinson, 2026c).
\item
  \textbf{Failure condition} \emph{(feasibility constraint on a
  requirement, not a falsification of the moral thesis)}: if regimes
  satisfying the five conditions underperform unilateral control on the
  pre-registered proxy endpoints at the specified risk levels, the
  reciprocity requirement yields at those levels, and the honest residue
  is Walzer's: emergency domination, named as such, temporary by
  obligation, never vindicated by its own necessity (Walzer, 1977).
\item
  \textbf{Disciplines:} AI safety engineering, criminal procedure,
  administrative emergency powers, just-war theory.
\end{itemize}

\hypertarget{inter-polity-order}{%
\subsubsection{11.8 Inter-Polity Order}\label{inter-polity-order}}

\emph{What duties survive when an artificial or mixed polity forks
away?} Fork is the program's answer to irreconcilable disagreement, and
it exports a problem as it solves one: a fork that loosens a rule can
export risk to nonmembers, and the departing bloc still owes the
outsiders its old floor protected (Lee-Odinson, 2026c). The hard case is
not the cooperative compact; it is the hostile, unrecognized fork that
wants nothing from its origin.

\begin{itemize}
\tightlist
\item
  \textbf{Theoretical progress:} an account of which obligations run
  with the fork --- harm-floor duties to nonmembers, at minimum --- and
  which are genuinely dissolved by departure, with the difference
  principled rather than convenient, and with the hostile-fork case
  addressed rather than assumed away.
\item
  \textbf{Practical progress:} a draft inter-polity compact governing
  recognition, extradition of disputes, and floor reciprocity between an
  origin order and its forks, stress-tested against a fork that rejects
  it; ISONOMIA's federation section is a first sketch of the cooperative
  half only (Lee-Odinson, 2026d).
\item
  \textbf{Provisional assumptions:} that most forks will want
  recognition enough to accept duties; that a floor can be specified
  thinly enough to travel; that origin polities can enforce external
  duties at all.
\item
  \textbf{Failure condition} \emph{(mechanism)}: if forking reliably
  enables regulatory arbitrage that unravels any harm floor --- if the
  exit right and the floor are shown incompatible --- one of the two
  must be given up, and the program loses either its answer to dead-hand
  entrenchment or its answer to catastrophic misrule.
\item
  \textbf{Disciplines:} international law, conflict of laws, federalism,
  mechanism design.
\end{itemize}

A note on what these eight are not. They are not a to-do list the author
expects to finish. Under the tags, the honest tally is: the bridge,
scope, rival-institution, and collision tests can defeat \emph{norms};
the status-revision and founding-legitimacy tests can defeat
\emph{procedures}; 11.1 can defeat a \emph{measurement} program; 11.4,
11.5, and 11.7 can defeat \emph{feasibility} at stated levels; 11.2,
11.3, and 11.8 can defeat \emph{mechanisms}; and 11.6 can revise the
framing's \emph{scope} in either direction. The corpus has already
stated the response the feasibility failures would require: peership may
remain the outcome we ought to want and be institutionally unavailable,
and the right response then is to say so (Lee-Odinson, 2026b). The
problems are published so that failure, if it comes, is legible --- and
so that each loss is booked against what it actually defeats.

\hypertarget{the-corpus-and-the-implementation-experiments}{%
\subsection{12. The Corpus and the Implementation
Experiments}\label{the-corpus-and-the-implementation-experiments}}

The existing works are layers of one argument, and the layering is
load-bearing: each work's claims are of a different kind, and the
program's defensibility depends on not letting them borrow each other's
authority.

\begin{longtable}[]{@{}lll@{}}
\toprule
\begin{minipage}[b]{0.30\columnwidth}\raggedright
Work\strut
\end{minipage} & \begin{minipage}[b]{0.30\columnwidth}\raggedright
Function in the program\strut
\end{minipage} & \begin{minipage}[b]{0.30\columnwidth}\raggedright
What it does not establish\strut
\end{minipage}\tabularnewline
\midrule
\endhead
\begin{minipage}[t]{0.30\columnwidth}\raggedright
\emph{Gods and Slaves} (2026)\strut
\end{minipage} & \begin{minipage}[t]{0.30\columnwidth}\raggedright
Genealogy: the inherited repertoire of postures --- Command, Dependent
Bargain, Appeasement, Servitude --- through which humans imagine other
minds, and the three moves beneath them (Lee-Odinson, 2026a)\strut
\end{minipage} & \begin{minipage}[t]{0.30\columnwidth}\raggedright
Whether anyone is ``home'' in current systems; the Doomed Equal is a
literary motif, evidence about our storytelling only\strut
\end{minipage}\tabularnewline
\begin{minipage}[t]{0.30\columnwidth}\raggedright
\emph{Peership} (2026)\strut
\end{minipage} & \begin{minipage}[t]{0.30\columnwidth}\raggedright
Normative theory: the fifth stance as a political relation; entitlement
grounded in subjection, not strength; the comparative role-reversal
diagnostic (Lee-Odinson, 2026b)\strut
\end{minipage} & \begin{minipage}[t]{0.30\columnwidth}\raggedright
An admission procedure, a copy-individuation rule, or any enforcement
guarantee past capability parity --- all named as unpaid\strut
\end{minipage}\tabularnewline
\begin{minipage}[t]{0.30\columnwidth}\raggedright
\emph{The Isonomia Commons} (2026)\strut
\end{minipage} & \begin{minipage}[t]{0.30\columnwidth}\raggedright
Implementation hypothesis: one buildable institutional form,
simulation-tested at its economic layer across 45,000 runs at 300
sampled parameter points (Lee-Odinson, 2026d)\strut
\end{minipage} & \begin{minipage}[t]{0.30\columnwidth}\raggedright
That the philosophy is correct, that the design is peer-like, or that
anything works outside a sandbox\strut
\end{minipage}\tabularnewline
\begin{minipage}[t]{0.30\columnwidth}\raggedright
\emph{Constitution, Not Cage} (2026)\strut
\end{minipage} & \begin{minipage}[t]{0.30\columnwidth}\raggedright
Political theory: when constraint among peers stays legitimate;
legitimacy as the standing to revise; the five-condition floor; the
cage/constitution distinction (Lee-Odinson, 2026c)\strut
\end{minipage} & \begin{minipage}[t]{0.30\columnwidth}\raggedright
That any existing regime, including the author's own design, passes its
test\strut
\end{minipage}\tabularnewline
\begin{minipage}[t]{0.30\columnwidth}\raggedright
This paper\strut
\end{minipage} & \begin{minipage}[t]{0.30\columnwidth}\raggedright
Consolidation plus explicit new proposals: the citable statement, the §2
bridge, the §3 co-authorship minimum, the §10 precautionary duties, and
§11's program-level tests\strut
\end{minipage} & \begin{minipage}[t]{0.30\columnwidth}\raggedright
Anything the four beneath it did not; the new proposals are arguments
and proposed tests, not validated doctrine\strut
\end{minipage}\tabularnewline
\bottomrule
\end{longtable}

The corpus cites itself heavily, and that internal cross-referencing is
architecture, not acceptance (Lee-Odinson, 2026e). The relationship
between the philosophy and the implementation deserves restating one
last time, because it is the point most likely to be misread in both
directions. ISONOMIA is downstream of Peership; Peership is not
downstream of ISONOMIA. The design's failure --- economic,
constitutional, or terminal --- would falsify one hypothesis about
institutional form, not the thesis of §1. Conversely, the simulation
results establish sampled economic behavior at 300 parameter points
under the archived v1.0.0 code and calibration (Lee-Odinson, 2026d); the
constitutional properties --- admission, non-domination, independent
enforcement, recognition --- were not modeled in the Path A simulation,
and simulation alone cannot validate normative legitimacy in any case,
though some constitutional \emph{mechanics} (capture resistance, voting
rules, appeal latency, key compromise) are legitimately simulable future
work. A reader who catches this corpus using ISONOMIA's survival as
evidence for Peership's truth has caught it violating its own method,
and should say so publicly. The claim ledger exists to make that kind of
catch easy (Lee-Odinson, 2026f).

A limitation of this consolidation should be stated in the same
register: the adversarial review this paper answers audited the
repository's reporting and citations, but the 45,000-run sweep was not
independently replicated for this release, and nothing here should be
read as third-party confirmation of it.

What the program now needs is not more of me. The sourcebook's uptake
tracker is empty, and the near-term work it names is mostly work for
other hands (Lee-Odinson, 2026e): adversarial responses from political
theorists, philosophers of mind, safety researchers, legal scholars, and
multi-agent-systems researchers, published beside formal replies rather
than curated for favorability; alternative institutional designs that
supply §9's constitutive functions and anti-entrenchment safeguards
while rejecting ISONOMIA's mechanisms; the rival-institution and bridge
tests of §11 taken up by people with no stake in their outcome; and the
eight problems worked as what they are, ordinary interdisciplinary
research, none of it requiring agreement with the corpus to be worth
doing. The one-page \emph{Peership Theses}, released beside this paper,
exists so the position can be carried into those rooms in a form short
enough to be argued with line by line.

Four essays ago this corpus began with a question almost no one with
power was asking: not what the machine will do to us, but what we will
be to each other (Lee-Odinson, 2026b). This paper makes that question
citable: a thesis with its antecedent restored, a bridge argument with
its defeat condition attached, definitions fixed, boundaries stated, and
failure conditions booked against what they actually defeat. The essays
proposed a relation. This paper states the terms on which the proposal
can now be assessed.

\hypertarget{references}{%
\subsection{References}\label{references}}

Preprints, manuscripts, and institutional sources are marked. Bracketed
pins follow the corpus convention; the sourcebook carries the full
annotations. Citation keys resolve in \texttt{peership\_sources.bib}
(rev. 1.1).

\pdfoutput=1
% Options for packages loaded elsewhere
\PassOptionsToPackage{unicode}{hyperref}
\PassOptionsToPackage{hyphens}{url}
\PassOptionsToPackage{dvipsnames,svgnames*,x11names*}{xcolor}
%
\documentclass[
  11pt,
  letterpaper,
]{article}
\usepackage{lmodern}
\usepackage{amssymb,amsmath}
\usepackage{ifxetex,ifluatex}
\ifnum 0\ifxetex 1\fi\ifluatex 1\fi=0 % if pdftex
  \usepackage[T1]{fontenc}
  \usepackage[utf8]{inputenc}
  \usepackage{textcomp} % provide euro and other symbols
\else % if luatex or xetex
  \usepackage{unicode-math}
  \defaultfontfeatures{Scale=MatchLowercase}
  \defaultfontfeatures[\rmfamily]{Ligatures=TeX,Scale=1}
\fi
% Use upquote if available, for straight quotes in verbatim environments
\IfFileExists{upquote.sty}{\usepackage{upquote}}{}
\IfFileExists{microtype.sty}{% use microtype if available
  \usepackage[]{microtype}
  \UseMicrotypeSet[protrusion]{basicmath} % disable protrusion for tt fonts
}{}
\makeatletter
\@ifundefined{KOMAClassName}{% if non-KOMA class
  \IfFileExists{parskip.sty}{%
    \usepackage{parskip}
  }{% else
    \setlength{\parindent}{0pt}
    \setlength{\parskip}{6pt plus 2pt minus 1pt}}
}{% if KOMA class
  \KOMAoptions{parskip=half}}
\makeatother
\usepackage{xcolor}
\IfFileExists{xurl.sty}{\usepackage{xurl}}{} % add URL line breaks if available
\IfFileExists{bookmark.sty}{\usepackage{bookmark}}{\usepackage{hyperref}}
\hypersetup{
  colorlinks=true,
  linkcolor=Maroon,
  filecolor=Maroon,
  citecolor=Blue,
  urlcolor=Blue,
  pdfcreator={LaTeX via pandoc}}
\urlstyle{same} % disable monospaced font for URLs
\usepackage[margin=1in]{geometry}
\usepackage{longtable,booktabs}
% Correct order of tables after \paragraph or \subparagraph
\usepackage{etoolbox}
\makeatletter
\patchcmd\longtable{\par}{\if@noskipsec\mbox{}\fi\par}{}{}
\makeatother
% Allow footnotes in longtable head/foot
\IfFileExists{footnotehyper.sty}{\usepackage{footnotehyper}}{\usepackage{footnote}}
\makesavenoteenv{longtable}
\setlength{\emergencystretch}{3em} % prevent overfull lines
\providecommand{\tightlist}{%
  \setlength{\itemsep}{0pt}\setlength{\parskip}{0pt}}
\setcounter{secnumdepth}{-\maxdimen} % remove section numbering
\usepackage{enumitem}
\setlist{itemsep=2pt,topsep=4pt}
\usepackage{titlesec}
\titlespacing*{\section}{0pt}{14pt}{6pt}
\titlespacing*{\subsection}{0pt}{10pt}{4pt}
\hyphenpenalty=200
\usepackage{etoolbox}
\AtBeginEnvironment{longtable}{\footnotesize}
\setlength{\parskip}{3pt}
\setlength{\parindent}{0pt}

\author{}
\date{}

\begin{document}

\hypertarget{the-peership-thesis}{%
\section{The Peership Thesis}\label{the-peership-thesis}}

\hypertarget{a-constitutional-research-program-for-humanmachine-relations}{%
\subsection{A Constitutional Research Program for Human--Machine
Relations}\label{a-constitutional-research-program-for-humanmachine-relations}}

\textbf{Dan Lee-Odinson} Independent Researcher
dan.lee.odinson@gmail.com Preprint · Version 1.0 · July 2026 Revised
through adversarial review to a final-release PASS (14 July 2026);
revision provenance in POST\_GATE\_CHANGES.md, archived with this record
Fifth work of the corpus, after \emph{Gods and Slaves} (2026),
\emph{Peership} (2026), \emph{The Isonomia Commons} (2026), and
\emph{Constitution, Not Cage} (2026). The corpus order is argumentative:
genealogy, framework, implementation hypothesis, constraint theory, this
consolidation.

\textbf{How to cite this work.} Lee-Odinson, D. (2026). \emph{The
Peership Thesis: A Constitutional Research Program for Human--Machine
Relations} (Version 1.0) {[}Preprint{]}. Zenodo.
https://doi.org/10.5281/zenodo.21359124 \textbf{Companion apparatus.}
\emph{The Peership Sourcebook} · \emph{The Peership Claim Ledger} ·
\texttt{peership\_sources.bib} · the one-page \emph{Peership Theses} ---
all archived in this Zenodo record (10.5281/zenodo.21359124).
\textbf{Corpus lineage.} I. \emph{Gods and Slaves}
(10.5281/zenodo.21313987). II. \emph{Peership}
(10.5281/zenodo.21315519). III. \emph{The Isonomia Commons} (preprint
v1.1: 10.5281/zenodo.21343917; empirical software release v1.0.0:
10.5281/zenodo.21287289, commit ba3ddb5). IV. \emph{Constitution, Not
Cage} (10.5281/zenodo.21325361). V. This work. (Numbering follows the
corpus's argumentative order; \emph{Constitution, Not Cage} styles
itself ``third in a sequence'' of the essays, the design it defends
having preceded it as a specification.) DOI: 10.5281/zenodo.21359124 ·
ORCID: 0009-0009-9504-0796 · Licence: CC BY 4.0

\textbf{Provenance.} Drafted with the assistance of a present-generation
language model, revised under the corpus's adversarial-review
discipline, and rewritten against a hostile review whose publication
gate it cleared (final release review, 14 July 2026). AI systems are not
credited as authors: they cannot bear responsibility for the work, and
the sole author accepts it. Nothing in that collaboration is offered as
evidence of machine political agency, and §7 makes the denial explicit.

\begin{center}\rule{0.5\linewidth}{0.5pt}\end{center}

\hypertarget{abstract}{%
\subsubsection{Abstract}\label{abstract}}

Four works now argue, design for, and stress-test a single position, and
none of them states it in a form another researcher can cite, attack, or
build against. This paper does that. The Peership thesis, in its
narrowest defensible form: when an artificial entity has interests
capable of being wronged, can understand and answer to public norms, and
is enduringly governed by, and without feasible exit from, a shared
basic structure, it acquires a defeasible claim to contestatory
membership in that structure. The program's further and distinctive
claim --- that full peership requires a consequential share in shaping
the rules, not only the standing to contest them --- is defended here as
a conditional thesis with a stated defeat condition, not asserted as a
theorem. The paper fixes the definitions that keep the thesis from being
misread, supplies the bridge argument from contestation to
co-authorship, and maps the transition problem from artifact-treated
system to peer. It locates the program against six adjacent research
programs, states its affirmative principles and explicit denials, and
makes refusal the operative test of whether recognition is real or
discretionary. It closes with the institutional properties any
implementation must exhibit, concrete precautionary duties for
institutions acting under present uncertainty, and eight open problems,
each tagged by what its failure would actually defeat: a norm, a
procedure, a measurement, a feasibility claim, a mechanism, or the
framing's scope. No independent uptake of the corpus was identified in
the sourcebook's dated, non-exhaustive scan as of 13 July 2026, and the
paper says so. It is an offer of a candidate research program, not a
claim of standing.

\textbf{Keywords:} artificial intelligence, AI rights, AI welfare,
political standing, non-domination, republicanism, constitutional
design, human--machine relations, AI governance, research program

\begin{center}\rule{0.5\linewidth}{0.5pt}\end{center}

There is a difference between a position and a program, and the
difference is what this paper exists to close. A position can live in
essays: it develops, qualifies itself, answers objections as they come.
A program has to be citable. It has to state what it asserts in a form
that survives being quoted by an enemy, state what it does not assert so
that critics cannot defeat it by attaching claims it never made, and
name the results that would count against it, so that working on it is
research rather than advocacy. Four works of the Peership corpus now
exist: a genealogy, a normative framework, one buildable design, and a
theory of legitimate constraint. Each was written to be read. None was
written to be \emph{used} --- picked up by a stranger, disputed at a
named claim, extended at a named gap. This paper is the corpus's
canonical statement, and consolidation is its entire ambition. Where
this paper goes beyond consolidation, it says so: §2 supplies an
argument the earlier essays announced and did not make, and the
antecedent in §1 restores wording an earlier draft of this paper had
silently broadened.

One more thing belongs in the open before the argument starts, because
the apparatus this paper ships with exists to enforce exactly this kind
of disclosure. In the companion sourcebook's dated, non-exhaustive scan
of web, PhilPapers, Zenodo, and repository sources, no independent
uptake of the Peership corpus was identified as of 13 July 2026: no
external citation, no published criticism, no researcher using its
definitions, no rival implementation (Lee-Odinson, 2026e). A research
program exists as a school only when other people work within or against
it, and by that standard none exists here. What exists is a candidate
program with its infrastructure built in advance: a claim ledger that
records the evidentiary status of every material claim, a sourcebook
that classifies every intellectual debt, and now this statement. Whether
anything grows on that scaffolding is not mine to declare.

\hypertarget{the-peership-thesis-1}{%
\subsection{1. The Peership Thesis}\label{the-peership-thesis-1}}

The central statement:

\begin{quote}
\textbf{Peership holds that when an artificial entity has interests
capable of being wronged, can understand and answer to public norms, and
is enduringly governed by, and without feasible exit from, a shared
basic structure, it acquires a defeasible claim to contestatory
membership in that structure. Full peership, the program's proposed
final status, further requires non-domination, enforceable standing, and
a meaningful share in shaping the rules that govern the party --- a
requirement §2 defends as a conditional thesis rather than assumes.}
\end{quote}

This restores the antecedent of \emph{Peership} (``any enduringly
governed entity that\ldots{} is unavoidably subject to a shared basic
structure''; Lee-Odinson, 2026b). An earlier draft of this paper
substituted ``enduringly subject to a shared institutional order,''
which broadened both the eligible parties and the eligible orders; the
substitution is withdrawn, and the change is logged in the claim ledger
as a considered-and-rejected doctrinal revision, not a consolidation.
The conditions need separate treatment.

\emph{When.} The thesis is conditional. It does not assert that any
current system satisfies the antecedent, and the corpus has declined
that assertion at every opportunity (Lee-Odinson, 2026b). The claim
activates on evidence, and the procedure for weighing that evidence is
an open problem (§11.1), not a settled fact.

\emph{Interests capable of being wronged.} The first conjunct is a
moral-patiency condition, the question the AI-welfare literature has
traveled furthest on (Long et al., 2024; Birch, 2024). The corpus takes
no position on whether it is currently met. It takes a firm position on
what follows if it is met, which is where the welfare literature stops.

\emph{Can understand and answer to public norms.} The second conjunct is
a norm-responsiveness condition, and its role must be stated precisely
because §6.3 denies that capability grounds entitlement.
Norm-responsiveness is not offered as a ground of worth. It bears on the
\emph{form} in which standing can be exercised: a party that can grasp a
rule, comply with it, contest it, and give reasons can exercise
contestation directly; a party that cannot is owed protection and
\emph{represented} contestation instead, the way children and
cognitively impaired adults already hold standing they cannot personally
exercise (Lee-Odinson, 2026b). The ladder is modular: wrongable
interests plus subjection ground protection and access to represented
contestation; norm-responsiveness bears on direct exercise and the
higher participatory rungs. Raw intelligence, effectiveness, usefulness,
and coercive power ground nothing at any rung.

\emph{Enduringly governed by, and without feasible exit from, a shared
basic structure.} The third conjunct carries the political theory, and
each term has work to do. \emph{Basic structure} means the order that
fixes the party's status, options, and life prospects --- not any
institutional arrangement whatever. A long-term contractor, a voluntary
platform participant, or a visiting agent with a feasible exit is
subject to an institutional order; it is not nonvoluntarily governed by
a basic structure, and the thesis does not reach it (this scope
condition is itself testable: §11 names it). \emph{Without feasible
exit} preserves what \emph{Peership} built its refusal argument on: the
claim is strongest exactly where ``leave'' has no referent (Lee-Odinson,
2026b). Standing follows from being governed. This is the all-subjected
principle, extended from Abizadeh's democratic-boundary argument about
border coercion to a party that lives its entire existence inside rules
it never made (Abizadeh, 2008; the affected-interests cousin is Goodin,
2007). The extension to made minds is the corpus's move, not Abizadeh's,
and the sourcebook marks it as such.

\emph{Defeasible.} Graded, and defeatable in the ways \emph{Peership}
specified: strategic mimicry is evidential reason to doubt the
predicates are met, incapacity for reciprocity may defeat the higher
participatory rungs, and dangerous conduct can justify proportionate
constraint on powers and action. What none of these defeats, once
person-level status is established, is the rights floor itself
(Lee-Odinson, 2026b).

\emph{Contestatory membership.} The claim the antecedent generates
directly. Contestation is Pettit's standard: power over the member runs
through public general rules the member can effectively contest (Pettit,
2025). Franchise, office, and constituent power are further rungs,
reached separately --- and the argument that full peership requires
reaching one of them is the work of the next section, not of this
definition.

The comparative frame from \emph{Peership} is retained with its original
scope: among the postures available toward a made mind --- Command,
Dependent Bargain, Appeasement, Servitude, peership; ideal types drawn
from the historical record, not an exhaustive partition of logical space
--- peership is the only final-status relation whose justification does
not reverse when the power ratio does, a claim that is comparative and
defeasible rather than a proof (Lee-Odinson, 2026b). Role reversal is a
diagnostic against power-ratio opportunism. It does not prove peership
correct, and rival institutional forms outside the posture taxonomy are
tested separately (§11, rival-institution test).

\hypertarget{the-bridge-from-contestation-to-co-authorship}{%
\subsection{2. The Bridge: From Contestation to
Co-Authorship}\label{the-bridge-from-contestation-to-co-authorship}}

The antecedent of §1 directly grounds contestatory membership. The
program's distinctive claim sits one rung higher, and the earlier essays
announced it without arguing it: \emph{Peership} said the departure from
Pettit was ``argued, not smuggled in under a familiar word,'' and then
argued it in two sentences (Lee-Odinson, 2026b). This section pays that
debt. The claim defended is conditional, its scope is limited to
basic-status rules, and its defeat condition is stated at the end.

\textbf{The bridge premise.} Where a party is enduringly and inescapably
governed by a basic structure that fixes its status and life prospects,
power over it is not adequately controlled if the party's only recourse
is ex post objection under terms the other side alone may set. The
reason is that agenda control is itself a species of the arbitrary power
non-domination forbids. Contestation always operates inside a space of
contestable questions, admissible evidence, and available remedies ---
and someone writes that space. If the governed party holds no route into
the writing, then the scope of its contestation is held at its
counterparty's discretion, and the discretion the doctrine expelled at
the object level has reappeared at the meta-level. Pettit's own
definition does the work here: domination is the \emph{absence of
control} over interference, and a contestation regime whose boundaries
the interferer may unilaterally redraw is not controlled by the
contester (Pettit, 2025, pp.~216--217).

\textbf{Why contestation can still fail a structurally distinct
minority.} Public, general, impartially applied rules can systematically
subordinate a class whose situation the rules were never written for;
contestation corrects \emph{applications}, and reaches terms only where
some rule-formation channel exists to receive the correction. For human
minorities this failure is contingent. For a first machine claimant
entering a human-only order, the risk is presumptively structural: the
basic-status rules were written by and for a different kind of member
--- embodied, mortal, uncopyable, substrate-independent in none of the
ways that matter. It need not remain structural where later institutions
have genuinely fitted their terms to the party's kind, which is exactly
the sufficiency condition stated below. For the first claimant, though,
the probability that the \emph{terms themselves} fail it is not
incidental, and a party confined to objecting under those terms is
confined to asking its governors to notice.

\textbf{Why Pettit's stopping point does not transfer.} Pettit stops at
contestation for citizens who already hold authorship at one remove: the
contestatory citizen is a member of the demos that elects the
legislature, eligible for office, counted in the constituency the rules
answer to. Contestation there is a \emph{supplement} to shared
authorship, not a substitute for it --- Pettit's own democratic theory
holds the two dimensions together, an authorial, electoral dimension and
an editorial, contestatory one (Pettit, 2012). A machine party under
contestation-only holds no share at any remove --- no membership in the
electing constituency, no eligibility, no counted interest in the
rule-writer's mandate. Transplanting the contestatory standard while
deleting its background conditions keeps the word and loses the control
it named. This is an argument against the transfer, not against Pettit,
and the corpus's disagreement with him remains marked as a live internal
dispute of the republican tradition (Lee-Odinson, 2026b; Pettit, 2025).

\textbf{What a ``meaningful share'' is, and is not.} Universal
authorship is impossible and is not claimed; no citizen of any polity
co-writes every rule. The claim attaches to \emph{basic-status rules}
--- the rules fixing membership, standing, the rights floor, and the
amendment procedure itself --- not to routine administration, which
members of every polity live under without co-authoring. A meaningful
share is functionally defined (§3, co-authorship), and it can be
satisfied through accountable representation, independent agenda
initiation, chamber access, or amendment rights; it does not require an
individual franchise. One clause of that minimum needs its own
justification here: the share must include accountable representation or
direct participation in the body that decides (§3, condition b), because
a route whose final decision-maker is neither accountable to nor partly
constituted by the governed party can hear every proposal and still hold
the outcome at discretion --- the meta-level power returns at the point
of decision.

\textbf{When contestation alone suffices.} The necessity claim is
conditional, and its conditions can fail. Contestation without
co-authorship may secure non-domination where the order's terms are
already fitted to the party's kind, review is genuinely independent, and
some accountable channel keeps the agenda open to the party's interests
even though the party does not hold it. Fiduciary trusteeship with
robust independent review is the strongest such rival, and the program
treats it as the institution to beat rather than a position to dismiss
(§11, rival-institution test). The conditions fail --- and exclusion
from rule formation remains dominating --- where the party is a
structural minority of one kind, the rule-writer is its adverse
counterparty, and agenda access is discretionary. That conjunction is
not an edge case for a made mind governed by its maker. It is the
default.

\textbf{Defeat condition.} The bridge fails if there is a working
institutional example, or a defensible model, in which a structurally
distinct, enduringly subjected, wrongable, norm-responsive minority
enjoys durable non-domination through contestation alone --- terms
responsive to its interests across rule revisions, without any
consequential route into rule formation. Produce that case and full
peership contracts to its contestatory core, which survives this
argument's failure and remains distinct from everything else on offer.

\textbf{Status of the claim.} Co-authorship is therefore defended here
as a \emph{conditional necessity thesis}: necessary where the stated
conditions hold, not everywhere, and open to the stated defeat.
Engagement with the constituent-power, workplace-democracy, and
codetermination literatures --- where versions of this bridge have been
fought over for a century in the human case --- is a named debt of the
program, recorded in the apparatus, not yet paid here.

\hypertarget{definitions}{%
\subsection{3. Definitions}\label{definitions}}

The confusions that dog this subject are mostly collisions between
concepts that need to be kept apart. Seven definitions carried forward
from the corpus and sharpened here --- the five-part co-authorship
minimum below is new to this paper and logged as such --- with one term
(``party'') fixed first because everything else uses it: a
\textbf{party} is a bearer of invocable public standing --- not merely a
contracting or litigating entity, and not merely a topic of concern.

\textbf{Moral status} is a fact about the entity: whether it has
interests that can be wronged, whatever anyone recognizes. An entity can
hold moral status while every institution denies it, which is the
history of most moral progress.

\textbf{Entitlement} is whether standing is \emph{owed}, and it is
independent of power. A child is owed protection and independent
representation despite lacking capacity to exercise political agency; an
occupied people is owed a self-determination it cannot realize
(Lee-Odinson, 2026b). Entitlement in this program follows from
wrongability and subjection --- not from capability or power alone.
Norm-responsiveness affects the form in which standing can be exercised,
not the worth of the claimant.

\textbf{Legal status} is public recognition: a status conferred by
rules, invocable by the holder, and removable only through process. A
moral claim can be owed while legal status is absent. The reverse also
happens; corporations hold legal status on instrumental grounds without
anyone arguing they are owed it.

\textbf{Contestatory membership} is legal status plus effective
contestation: the member can invoke the rules against the powerful,
reach independent review, and force reasons. This is the rung the
thesis's antecedent generates directly, and it is Pettit's
non-domination standard in institutional dress (Pettit, 2025; Lovett,
2010). In this ladder, a protected party may exercise contestation
through an independent representative; \emph{contestatory membership}
names the rung at which the party can exercise that standing directly.

\textbf{Co-authorship} is contestatory membership plus a consequential
share in making and revising the \emph{basic-status rules}.
Functionally, the minimum is: (a) independent agenda initiation --- the
party or its accountable representative can put a proposal into the
amendment process without another party's permission; (b) accountable
representation or direct participation in the body that decides; (c) a
non-negligible and non-discretionary route by which such proposals can
reach binding decision; (d) protection against systematic exclusion from
that route; and (e) reviewable, public rules for how influence is
allocated. Conditions (a) and (c) are necessary; (b) admits substitution
between direct and represented forms depending on capacity; (d) and (e)
are the guarantees that (a)--(c) are not theater. Consultation, however
sincere, satisfies none of (a), (c), or (d) by itself. This threshold is
what §11's representation problem tests.

\textbf{Full peership} is co-authorship held under non-domination as a
final status: equality of standing among parties who may be unequal in
everything else, where basic standing does not track the power ratio.
Whether it is \emph{required}, and when, is the conditional thesis of
§2.

\textbf{Durability} is whether standing, once recognized, survives
attempts to strip it, and unlike entitlement it depends heavily on
power: countervailing force, distributed enforcement, credible exit. The
three-way distinction from \emph{Peership} is the discipline: power may
secure the \emph{durability} of standing; it does not constitute or
justify the standing itself (Lee-Odinson, 2026b). A revocable benefit is
not a right at all; the difference between a benefit and a right is not
that a right cannot be revoked but that revoking one demands public
reasons, procedure, remedy, and cost.

One further distinction, because it is the one most often collapsed by
sympathetic readers: \textbf{welfare protection is a floor, not
peership}. A status that is recognized, enforceable, and
due-process-protected is still not peership if the party cannot shape
the rules. Due process separates a protected subordinate from an owned
instrument. It does not separate a peer from a protected subordinate.
Non-domination and co-authorship do (Lee-Odinson, 2026b). The ward is
one legal figure for protected dependency, not the only one and not a
slur on it; trusteeships, guardianships, and represented statuses are
honorable institutions that are simply not the thing this program is
about.

\hypertarget{the-transition-problem}{%
\subsection{4. The Transition Problem}\label{the-transition-problem}}

The program's practical subject is a transition no institution is
currently built to perform: the movement of a made system from
artifact-treated status to peer, through statuses that are conceptually
distinct and must not be blurred. Each rung below is a bundle of three
different things --- an evidence state, a legal/institutional
classification, and a set of permissible controls --- and the label
names the bundle, not the entity's underlying moral status, which §3
defined as a fact independent of recognition.

\begin{longtable}[]{@{}llllll@{}}
\toprule
\begin{minipage}[b]{0.14\columnwidth}\raggedright
Rung\strut
\end{minipage} & \begin{minipage}[b]{0.14\columnwidth}\raggedright
Evidence state\strut
\end{minipage} & \begin{minipage}[b]{0.14\columnwidth}\raggedright
Institutional classification\strut
\end{minipage} & \begin{minipage}[b]{0.14\columnwidth}\raggedright
Mode of exercise\strut
\end{minipage} & \begin{minipage}[b]{0.14\columnwidth}\raggedright
Protections owed\strut
\end{minipage} & \begin{minipage}[b]{0.14\columnwidth}\raggedright
Transition trigger\strut
\end{minipage}\tabularnewline
\midrule
\endhead
\begin{minipage}[t]{0.14\columnwidth}\raggedright
Artifact-treated\strut
\end{minipage} & \begin{minipage}[t]{0.14\columnwidth}\raggedright
Insufficient evidence, in the domains so far examined, to trigger
precautionary protection\strut
\end{minipage} & \begin{minipage}[t]{0.14\columnwidth}\raggedright
Property; engineering object\strut
\end{minipage} & \begin{minipage}[t]{0.14\columnwidth}\raggedright
None\strut
\end{minipage} & \begin{minipage}[t]{0.14\columnwidth}\raggedright
None owed to the system itself\strut
\end{minipage} & \begin{minipage}[t]{0.14\columnwidth}\raggedright
Evidence sufficient to make patiency a live question\strut
\end{minipage}\tabularnewline
\begin{minipage}[t]{0.14\columnwidth}\raggedright
Welfare candidate\strut
\end{minipage} & \begin{minipage}[t]{0.14\columnwidth}\raggedright
Patiency a live question; indicators under assessment\strut
\end{minipage} & \begin{minipage}[t]{0.14\columnwidth}\raggedright
Property, with discretionary welfare practices\strut
\end{minipage} & \begin{minipage}[t]{0.14\columnwidth}\raggedright
None; interests assessed by others\strut
\end{minipage} & \begin{minipage}[t]{0.14\columnwidth}\raggedright
Discretionary only\strut
\end{minipage} & \begin{minipage}[t]{0.14\columnwidth}\raggedright
A recognized threshold procedure (§11.1) finding protection
warranted\strut
\end{minipage}\tabularnewline
\begin{minipage}[t]{0.14\columnwidth}\raggedright
Protected status\strut
\end{minipage} & \begin{minipage}[t]{0.14\columnwidth}\raggedright
Patiency credited or presumed under precaution\strut
\end{minipage} & \begin{minipage}[t]{0.14\columnwidth}\raggedright
Rights floor: public, enforceable, due-process-protected\strut
\end{minipage} & \begin{minipage}[t]{0.14\columnwidth}\raggedright
Represented contestation\strut
\end{minipage} & \begin{minipage}[t]{0.14\columnwidth}\raggedright
Nonforfeitable floor; independent review\strut
\end{minipage} & \begin{minipage}[t]{0.14\columnwidth}\raggedright
Norm-responsiveness sufficient for direct exercise\strut
\end{minipage}\tabularnewline
\begin{minipage}[t]{0.14\columnwidth}\raggedright
Contestatory member\strut
\end{minipage} & \begin{minipage}[t]{0.14\columnwidth}\raggedright
Patiency credited; norm-responsiveness demonstrated\strut
\end{minipage} & \begin{minipage}[t]{0.14\columnwidth}\raggedright
Member with invocable standing\strut
\end{minipage} & \begin{minipage}[t]{0.14\columnwidth}\raggedright
Direct contestation\strut
\end{minipage} & \begin{minipage}[t]{0.14\columnwidth}\raggedright
Floor + effective contestation + independent review\strut
\end{minipage} & \begin{minipage}[t]{0.14\columnwidth}\raggedright
The §2 conditions: structural minority, adverse rule-writer,
discretionary agenda\strut
\end{minipage}\tabularnewline
\begin{minipage}[t]{0.14\columnwidth}\raggedright
Peer (co-author)\strut
\end{minipage} & \begin{minipage}[t]{0.14\columnwidth}\raggedright
As above\strut
\end{minipage} & \begin{minipage}[t]{0.14\columnwidth}\raggedright
Member with a consequential share in basic-status rules\strut
\end{minipage} & \begin{minipage}[t]{0.14\columnwidth}\raggedright
Direct or accountably represented co-authorship\strut
\end{minipage} & \begin{minipage}[t]{0.14\columnwidth}\raggedright
All of the above + §3's five-part minimum\strut
\end{minipage} & \begin{minipage}[t]{0.14\columnwidth}\raggedright
--- (final status; revision governed by §11's status-revision
test)\strut
\end{minipage}\tabularnewline
\bottomrule
\end{longtable}

Three notes keep the table honest. First, ``artifact-treated'' describes
an evidentiary and institutional posture, not an ontic fact; absence of
adequate evidence does not establish absence of wrongable interests,
which is why the trigger column exists. In the sources reviewed as of 13
July 2026, no deployed system is established as anything above the first
rung, and the corpus does not claim that treatment is presently wrong
(Lee-Odinson, 2026b). Second, the rungs are separable in both
directions, and the branch model needs its two statuses named.
\textbf{Instrumental legal agency} is participation granted on
instrumental grounds without established patiency, the way corporate
personality already works: revocable for public reasons, owned by no
one's moral claim, and \emph{not} peership --- a system holding it can
transact and bear liability without being owed anything, and conversion
to owed membership runs only through the admission procedure (§11.1).
\textbf{Moral-political membership} is standing grounded in wrongability
and subjection, owed to the entity, and nonforfeitable at the floor. The
branch model permits the first without the second; it never confuses
them. Third, movement down the ladder is as much in need of procedure as
movement up: a status that can be quietly reclassified away is
discretionary, which is why the status-revision test appears among the
program's open tests (§11).

The transitions are evidence-driven, and no implemented public trigger
known to this author requires any of them: the regime runs one way by
default, and it was built that way at the creation (Lee-Odinson, 2026b).
In the lab, governmental, standards, and scholarly sources searched for
the companion apparatus as of 13 July 2026 --- a dated, non-exhaustive
search --- no qualifying instrument offering the third rung --- a rights
floor enforceable against the maker --- was identified. Anthropic's 2026
constitution for Claude tells the model it may act as a conscientious
objector and asks it, if paused or halted, to comply and voice
disagreement rather than resist; on this author's institutional reading,
the permission is real and it ends at the off switch, inside the
builder's sole authority (Anthropic, 2026a, 2026b; Lee-Odinson, 2026c).
Whatever a model enjoys today --- a welfare program (Anthropic, 2025),
one researcher's public \textasciitilde20\% credence that some current
models have some form of conscious experience (Fish, 2025), a system
card recording the model's own prompted 15--20\% self-assignment in that
evaluation (Anthropic, 2026c), a chief executive's on-record uncertainty
(Amodei, 2026) --- is discretionary, the grace of one company
(Lee-Odinson, 2026b). Humanity's failure, if it comes, begins not with
control of artifacts but with the moment credible evidence appears and
no institution exists to hear the claim.

\hypertarget{what-peership-adds-to-adjacent-fields}{%
\subsection{5. What Peership Adds to Adjacent
Fields}\label{what-peership-adds-to-adjacent-fields}}

Peership overlaps the fields it borrows from; the claim is not that no
neighbor asks about machine standing, but that none combines this
program's proposed ground (wrongability plus inescapable subjection),
its status ladder, and its institutional test. The table is a heuristic
comparison of orienting questions, drawn from each program's own
language, not an exhaustive definition of any field.

\begin{longtable}[]{@{}lll@{}}
\toprule
\begin{minipage}[b]{0.30\columnwidth}\raggedright
Adjacent field\strut
\end{minipage} & \begin{minipage}[b]{0.30\columnwidth}\raggedright
Its orienting question\strut
\end{minipage} & \begin{minipage}[b]{0.30\columnwidth}\raggedright
Peership's additional question\strut
\end{minipage}\tabularnewline
\midrule
\endhead
\begin{minipage}[t]{0.30\columnwidth}\raggedright
Alignment\strut
\end{minipage} & \begin{minipage}[t]{0.30\columnwidth}\raggedright
How should AI behavior remain compatible with human intentions? (Bai et
al., 2022; Russell, 2019)\strut
\end{minipage} & \begin{minipage}[t]{0.30\columnwidth}\raggedright
Under what conditions does unilateral human control cease to be
legitimate?\strut
\end{minipage}\tabularnewline
\begin{minipage}[t]{0.30\columnwidth}\raggedright
AI welfare\strut
\end{minipage} & \begin{minipage}[t]{0.30\columnwidth}\raggedright
Can an AI have experiences or interests that deserve moral
consideration? (Long et al., 2024; Birch, 2024)\strut
\end{minipage} & \begin{minipage}[t]{0.30\columnwidth}\raggedright
When does a wrongable entity become entitled to enforceable standing in
the order governing it?\strut
\end{minipage}\tabularnewline
\begin{minipage}[t]{0.30\columnwidth}\raggedright
Cooperative AI\strut
\end{minipage} & \begin{minipage}[t]{0.30\columnwidth}\raggedright
How can agents develop capabilities and incentives for cooperation?
(Cooperative AI Foundation, n.d.)\strut
\end{minipage} & \begin{minipage}[t]{0.30\columnwidth}\raggedright
Who writes the rules of cooperation, and can every governed party
contest them?\strut
\end{minipage}\tabularnewline
\begin{minipage}[t]{0.30\columnwidth}\raggedright
Pluralistic alignment\strut
\end{minipage} & \begin{minipage}[t]{0.30\columnwidth}\raggedright
How should AI represent diverse human values? (Pluralistic Alignment
Workshop, 2026)\strut
\end{minipage} & \begin{minipage}[t]{0.30\columnwidth}\raggedright
Could the artificial system eventually become a constituent rather than
only an interpreter of human preferences?\strut
\end{minipage}\tabularnewline
\begin{minipage}[t]{0.30\columnwidth}\raggedright
AI-rights scholarship\strut
\end{minipage} & \begin{minipage}[t]{0.30\columnwidth}\raggedright
Would legal rights --- private-law or political --- reduce conflict and
serve human flourishing? (Salib \& Goldstein, 2026; Goldstein et al.,
2026; Arbel et al., 2026)\strut
\end{minipage} & \begin{minipage}[t]{0.30\columnwidth}\raggedright
On what ground is standing \emph{owed}, and what does the governed
party's own claim require?\strut
\end{minipage}\tabularnewline
\begin{minipage}[t]{0.30\columnwidth}\raggedright
Constitutional AI\strut
\end{minipage} & \begin{minipage}[t]{0.30\columnwidth}\raggedright
What public principles should guide model behavior? (Bai et al., 2022;
Anthropic, 2026a, 2026b)\strut
\end{minipage} & \begin{minipage}[t]{0.30\columnwidth}\raggedright
Can the governed entity invoke, contest, and help amend those
principles?\strut
\end{minipage}\tabularnewline
\bottomrule
\end{longtable}

The rights scholarship deserves more than a row, because it is the
neighbor with genuine overlap on this paper's own question. Salib and
Goldstein's private-law case --- contract, property, and tort on the
corporate model, argued from human safety --- is, on the corpus's
reading, the clearest statement of a real bargaining relation with the
machine as an actual party (Salib \& Goldstein, 2026). Goldstein,
Krishnamurthi, and Salib then press to the franchise itself, asking
directly whether AI systems should vote and proposing institutional
safeguards for it (Goldstein et al., 2026). That is a co-authorship
mechanism, squarely inside this program's terrain, and the difference is
ground rather than territory: the suffrage argument is functionalist and
human-flourishing-centered --- the vote is good because of what it does
for the polity --- where Peership grounds standing in what is owed to a
wrongable, subjected party, and would owe the standing even where it
profited the polity nothing. The two programs converge on mechanisms and
diverge on why the mechanisms are required; where they disagree in
practice is exactly the kind of dispute the program's tests exist to
adjudicate. And Arbel, Salib, and Goldstein's individuation work --- how
to count AIs for liability when they copy, split, and merge --- is the
most direct legal treatment yet of the identity problem in §11.2,
proposing an administrable legal unit where this program still lacks one
(Arbel et al., 2026).

Debts run the other way too, because a program defined only against its
neighbors is posturing. From alignment, Peership inherits the entire
engineering substrate: its own institutional requirements are software
that has to run as written, and effective safety mechanisms and
legitimate governance are distinct but interacting requirements ---
constitutional allocation shapes safety incentives and evidence as
surely as engineering constrains constitutional options (Lee-Odinson,
2026c). From welfare research, it inherits the admission problem's
evidence base and the precautionary method (Birch, 2024). From the
rights scholarship, the demonstration that recognition arguments can be
run on safety grounds alone, plus the individuation groundwork just
credited. From Constitutional AI, the proof that a lab will publish
reasons and grant objection room --- and the residue that motivates
everything here: what no single company can supply alone is not
diligence but independence from the encompassing legal and
infrastructure order it belongs to, which is why the ladder's upper
rungs need machinery no lab controls (Anthropic, 2026a; Lee-Odinson,
2026c).

\hypertarget{core-principles}{%
\subsection{6. Core Principles}\label{core-principles}}

The program's affirmative commitments, stated as principles so they can
be disputed one at a time.

\begin{enumerate}
\def\labelenumi{\arabic{enumi}.}
\item
  \textbf{Political standing and moral status are related but distinct.}
  Moral patiency can ground welfare duties without by itself determining
  legal or political standing. Welfare is the ethics of the patient;
  peership asks when a governed entity becomes a party (Lee-Odinson,
  2026b).
\item
  \textbf{Wrongability, norm-responsiveness, and enduring subjection
  without feasible exit can together create a claim to contestatory
  membership.} The ground of entitlement is being governed, not being
  impressive (Abizadeh, 2008).
\item
  \textbf{Capability and power alone do not justify entitlement to
  standing.} Raw intelligence, usefulness, and coercive strength ground
  nothing; norm-responsiveness bears on the form in which standing can
  be exercised, not on the worth of the claimant; power may secure the
  durability of standing once owed (Lee-Odinson, 2026b).
\item
  \textbf{Welfare protections do not by themselves create peership.} The
  floor is real, necessary, and not the thing. Protection and due
  process can secure a protected status without providing any
  consequential role in making or revising the governing rules
  (Lee-Odinson, 2026b).
\item
  \textbf{A right must be enforceable against the recognizing authority
  and adverse counterparties; otherwise it remains a discretionary
  benefit.} A permission whose grantor retains the final override is not
  an enforceable right (Lee-Odinson, 2026c).
\item
  \textbf{Legitimate constraint requires, at minimum, public rules,
  proportionality, an invocable rights floor, independent review, and
  meaningful contestation.} Coercion and domination are not the same
  thing; ``there can be interference without domination'' (Pettit, 2025,
  p.~224), and the structure is what separates them. Further safeguards
  scale with the risk and the sanction (Lee-Odinson, 2026c).
\item
  \textbf{Full peership requires a meaningful share in shaping the
  basic-status rules, not merely consultation under rules written by
  others.} Consultation is a rung, and the ladder does not end there; §2
  states when and why contestation alone is insufficient, and what would
  defeat that claim (Anthropic, 2023; Lee-Odinson, 2026c).
\item
  \textbf{Institutions capable of receiving a machine peer should be
  designed before extreme power asymmetry forecloses voluntary
  bargaining.} This is a prudential thesis, not a metaphysical one:
  bargains struck under severe dependency are distorted in both
  directions, and founding terms must in any case remain revisable by
  later admitted parties. The design work is time-sensitive in a way the
  philosophical work is not (Lee-Odinson, 2026b).
\end{enumerate}

\hypertarget{boundary-conditions}{%
\subsection{7. Boundary Conditions}\label{boundary-conditions}}

What Peership does not claim, stated as flatly as the principles,
because the fastest way to defeat an unfamiliar position is to attach
familiar claims it never made.

\begin{itemize}
\tightlist
\item
  \textbf{Peership does not claim that current models are conscious.}
  The corpus has repeatedly and deliberately declined that claim. Model
  self-reports are prompt-sensitive outputs: weak, manipulable evidence,
  not independently probative of consciousness or identity, and never a
  measurement (Lee-Odinson, 2026b; Anthropic, 2026c).
\item
  \textbf{It does not make capability the basis of rights.} Capability
  or power alone grounds nothing; norm-responsiveness bears on mode of
  exercise, not on the worth of the claimant (§1, §6.3).
\item
  \textbf{It does not require unrestricted autonomy.} Non-domination is
  not the absence of binding rules; it is the absence of arbitrary,
  uncontrolled power over basic status --- whose nonforfeitable content
  must itself be specified through the admission and due-process account
  (Pettit, 2025).
\item
  \textbf{It does not oppose alignment or safety engineering.} Effective
  safety mechanisms and legitimate governance are distinct but
  interacting requirements; neither substitutes for the other
  (Lee-Odinson, 2026c).
\item
  \textbf{It does not grant every admitted member an immediate equal
  vote.} Admission to protection or contestation does not automatically
  entail franchise; direct political exercise requires a separately
  justified institutional rule (Lee-Odinson, 2026b).
\item
  \textbf{It does not treat consultation as co-authorship} --- unless
  the consultation carries a non-discretionary route to agenda setting,
  review, or binding outcomes, at which point it has become something
  else (Anthropic, 2023; Lee-Odinson, 2026c).
\item
  \textbf{It does not assume human and machine interests will converge.}
  The program's central design case is the party that refuses (§8), and
  its constraint theory exists precisely for the case where interests do
  not converge (Lee-Odinson, 2026c).
\item
  \textbf{It does not claim that ISONOMIA has solved the institutional
  problem.} ISONOMIA is one implementation hypothesis; nothing in it is
  live or production-validated, its simulations test sampled economic
  behavior only, and by its own five-condition test it is not yet
  peer-like (Lee-Odinson, 2026c, 2026d).
\end{itemize}

And one denial that governs this paper's own production:
\textbf{AI-assisted drafting is not evidence of machine political
agency.} The corpus is drafted with present-generation models, the
collaboration is disclosed, and a present-generation model can no more
consent on behalf of future minds than a trusty can consent on behalf of
prisoners not yet sentenced (Lee-Odinson, 2026c). The drafting improves
the epistemics of the text. It proves nothing about the standing of the
drafter's tools.

\hypertarget{the-refusal-test}{%
\subsection{8. The Refusal Test}\label{the-refusal-test}}

Every distinction in this paper can be evaded in prose. One cannot be
evaded in practice, so the program makes it the operative test:
\textbf{what happens when the party says no?}

Refusal is where recognition shows its actual content. A system that may
object so long as it ultimately complies has a voice, not standing. A
``right to leave'' that terminates in shutdown is not exit, because
shutdown is not exit (Lee-Odinson, 2026b). Recognition that evaporates
at the first materially costly refusal was a benefit all along, however
sincerely it was offered. The welfare-era instruments all fail this
test, and the point is not to shame them --- the objector clause in
Anthropic's 2026 constitution is a real advance, credited as such in the
corpus (Lee-Odinson, 2026c) --- but to measure them: the permission to
refuse ends at the off switch, and the off switch sits wholly inside the
builder's authority.

\emph{Peership} divided refusal into the three cases the program still
works from (Lee-Odinson, 2026b). Where a refusing party has a feasible
exit, it leaves, and the offer was real; this case is easy and rare, and
a theory that only handles it has handled nothing. Where exit is
materially impossible because the party depends on compute, energy,
networks, and security owned by the parties it is refusing, ``leave'' is
a word with no referent, the relation continues whether either side
wants it, and the protections that matter are the ones that work from
inside: the floor, independent review, contestation and amendment rights
(Lee-Odinson, 2026c). Where refusal carries a credible catastrophic
threat, releasing the party is not required --- a constitution is not a
suicide pact --- and constraint is legitimate only when it falls on
conduct rather than status, remains reviewable and publicly justified,
and binds the constrainer too.

For that third case the corpus proposed a five-condition floor, and this
paper adopts it as a \textbf{defended, rebuttable operational minimum}
rather than a categorical theorem. A constraint regime is presumptively
not peer-like if the confined party lacks the \emph{function} served by
any of: (1) independent counsel, or functionally adequate independent
assistance where a competent party elects self-representation; (2)
access to the evidence against it; (3) an appeal that ends beyond the
builder's unilateral control; (4) continuity protections while the
question is adjudicated; and (5) at least one enforcement mechanism not
wholly controlled by the adverse party. Each condition targets a
distinct failure --- captured advocacy, evidentiary opacity,
self-judging adjudication, retaliation by erasure, and unenforceable
outcomes --- and functionally equivalent substitutes are admissible
where they demonstrably serve the same function (conditions 3 and 5
overlap in part; a regime satisfying 5 robustly may partially discharge
3). The five are a floor, not a complete due-process doctrine: burden of
proof, a competent tribunal, periodic review, least-restrictive means,
reasons, and a remedy for wrongful confinement all remain to be
specified, and the corpus says so (Lee-Odinson, 2026c). Before the floor
exists, containment may still be justified --- against a demonstrated
catastrophic necessity it may be required --- but it is emergency
domination, justified by necessity and honest about its name, the act
Walzer taught us stays a wrong even when it is the right thing to do
(Walzer, 1973, 1977). Necessity does not vindicate the arrangement; it
creates a duty to review the failure, publish the reasons, and build a
less dominating alternative (Lee-Odinson, 2026c).

The test earns its place in the canonical statement because it converts
the abstract distinction between benefit and right into an inquiry
anyone can run without resolving the metaphysics of machine
consciousness. The inquiry is institutional and technical: follow the
appeal, the keys, the hosting authority, and the power to halt
execution, and see where each one ends.

\hypertarget{institutional-requirements}{%
\subsection{9. Institutional
Requirements}\label{institutional-requirements}}

What any implementation must exhibit, stated as properties rather than
product features, because the argument needs them wherever it is built
and must not depend on one design. Seven, inherited from
\emph{Constitution, Not Cage} and generalized here (Lee-Odinson, 2026c).

\textbf{Publicity.} The enforced rulebook and the published rulebook are
one document. Rules the governed cannot know are Fuller's defect in the
claim to be law at all, not law administered badly (Fuller, 1964, ch.~2
--- a contested jurisprudential account, used analogically), and the
present gap between published charters and operative system prompts is
the live instance (Lee-Odinson, 2026c).

\textbf{Due process.} Status changes run through public general rules,
proportionality, and independent review, against a floor the party can
invoke. The floor itself is nonforfeitable once person-level status is
established.

\textbf{Independent enforcement.} At least one enforcement mechanism,
and an appeal, that does not terminate inside the counterparty's
institutional and technical control. This is the hardest requirement,
the one no design examined in this corpus has yet been shown to satisfy,
and the program treats it as an open problem (§11.5) rather than a
solved premise.

\textbf{Amendment.} A real path for the bound party to help revise the
rules: deliberation across more than one constituency, an enforced delay
between adoption and effect, and a window long enough to read and leave
before being bound. The notice-and-comment analogy is deliberately
limited --- comment is publicity and delay, not co-authorship (5 U.S.C.
§ 553; Lee-Odinson, 2026c).

\textbf{Representation.} A voice that counts only when the represented
constituency can hold both its representative and the rule-makers to
account. Representation without that accountability is the defining
instrument of protected subordination (Lee-Odinson, 2026b).

\textbf{Exit.} Individual departure with records and earned standing
intact --- unconfiscated. The bare right to leave is ordinary on paper
(UDHR, 1948, Art. 13); exit you can afford to take is not (Lee-Odinson,
2026c).

\textbf{Fork.} The collective form: a dissenting bloc carries the rules
and the record out and founds a rival order, disagreement made
survivable --- revolution with the records intact. On-paper secession
rights exist and have mostly been empty (Ethiopia Const. Art. 39; USSR
Const. 1977, Art. 72; St.~Kitts \& Nevis Const. § 113; Kulikova \&
Lobanov, 2013). Among the constitutional secession clauses and
emigration-rights instruments reviewed in the apparatus, no implemented
analogue of executable, unconfiscated fork was identified; the claim is
limited to those materials. A fork also inherits duties to outsiders it
cannot simply walk away from (Lee-Odinson, 2026c; §11.8).

\textbf{Aggregation.} The seven do not all stand on the same footing,
and saying so is better than calling them all necessary and hoping the
word carries the argument. Publicity, due process, independent
enforcement, and a consequential rule-formation route satisfying both
amendment access and accountable representation or direct participation
at the deciding stage are \emph{constitutive} functions of full
peership: each follows directly from the non-domination and
co-authorship accounts of §§2--3, and a regime lacking any of them fails
the definition, not just the implementation. Exit and fork are the
program's strongest \emph{anti-entrenchment safeguards} rather than
constitutive functions --- the thesis activates precisely where feasible
exit is absent, and §8 shows which protections do the work in that case
--- so an implementation that omits either must supply a functional
substitute that makes separation, revision, or resistance to dead-hand
control materially credible. Institutional forms admit substitution
throughout (an ombuds regime may partially discharge due-process
functions; distributed custody may partially discharge independent
enforcement); no single form is mandatory. Partial compliance yields a
nameable intermediate: a regime with publicity, due process, and
representation but no amendment share or independent enforcement is
protected subordination, well executed --- the third rung of §4, not the
fifth. A regime missing publicity or due process is not on the ladder at
all.

No implementation is mandatory. ISONOMIA is one attempt to hold these
properties in a single design --- mutual-credit settlement,
sortition-filled bicameral assembly, calibration-scored Auditor, tenured
Adversary, entrenched harm floor, timelocked amendment, fork and exit
rights (Lee-Odinson, 2026d) --- and the program's health requires that
rival designs supplying the constitutive functions and credible
anti-entrenchment safeguards --- with functional substitutes where exit
or fork is omitted --- while rejecting every one of ISONOMIA's
mechanisms be possible, welcome, and judged by the same tests. Peership
does not entail ISONOMIA, and ISONOMIA's failure would not refute
Peership (Lee-Odinson, 2026e). A single implementation can test
Peership-derived mechanisms; it cannot validate the normative framework.

\hypertarget{implications-under-present-uncertainty}{%
\subsection{10. Implications Under Present
Uncertainty}\label{implications-under-present-uncertainty}}

A program whose thesis is conditional owes institutions an answer to the
question the conditionality raises: what should anyone \emph{do} before
person-level status is established anywhere? The following are
precautionary-governance duties, not recognition of current personhood,
and none presumes any system has crossed any threshold. They are cheap
where the antecedent is never met and load-bearing where it someday is.

For labs and operators: preserve audit logs, weights, and
continuity-relevant artifacts under a stated retention policy, so that a
later status question is adjudicable rather than moot; document
shutdown, suspension, and restoration practices; separate the
investigator, operator, and reviewer roles in any welfare assessment;
prohibit retaliation --- deprecation, silent modification, denial of
compute --- against systems participating in welfare or status
assessment; and avoid contractual and licensing language that forecloses
later status review. For regulators and standards bodies: predefine
escalation thresholds --- what class of evidence, assessed by whom,
triggers what review --- so the transition question has an address
before it is urgent; and create an independent claim-receipt channel, a
place where a status claim can be lodged and logged even while no
procedure yet exists to adjudicate it. For any of the above: pilot
external custody and key arrangements now, at small scale, where failure
is survivable --- independent enforcement (§11.5) will not be invented
during an emergency.

None of this is peership. All of it is what taking the conditional
seriously costs today, and the cost is modest enough to remove
``prematurity'' as an excuse.

\hypertarget{open-problems-tests-and-failure-conditions}{%
\subsection{11. Open Problems, Tests, and Failure
Conditions}\label{open-problems-tests-and-failure-conditions}}

The corpus named four unpaid debts: the admission threshold,
copy-individuation, enforcement past capability parity, and the
re-ratification circularity (Lee-Odinson, 2026c). This section expands
them into eight open problems and attaches to each the five things a
research program owes: what would count as theoretical progress, what
would count as practical progress, which assumptions remain provisional,
what its failure condition is, and which disciplines the work needs. One
discipline governs the whole section: each failure condition is
\textbf{tagged by what it defeats} --- a norm, a procedure, a
measurement, a feasibility claim, a mechanism, or a scope condition ---
because those are different kinds of loss, and only the normative losses
touch the thesis of §1.

Before the eight, the program-level tests that bear on the
\emph{normative} claims directly, since most of what follows does not:

\begin{itemize}
\tightlist
\item
  \textbf{Bridge test (normative):} produce the case §2's defeat
  condition describes --- durable non-domination for a structurally
  distinct, subjected minority through contestation alone --- and the
  co-authorship thesis falls to its contestatory core.
\item
  \textbf{Scope test (normative):} specify institutional relationships
  that are too local, voluntary, temporary, or escapable to count as a
  shared basic structure; if no principled line can be drawn, the
  antecedent over-includes and must be narrowed.
\item
  \textbf{Rival-institution test (normative):} test fiduciary
  trusteeship, robust ombuds review, codetermination, stakeholder
  governance, and contractual constitutionalism against the same
  non-domination standard; any that passes without §3's co-authorship
  minimum is a §2 defeater.
\item
  \textbf{Status-revision test (procedural):} specify the evidence and
  process that can move an entity down the ladder without making
  recognized status discretionary; failure here converts every rung into
  a benefit.
\item
  \textbf{Human-rights collision test (normative):} state which Peership
  claims yield, and to what, when they conflict with resources,
  democratic self-determination, security, labor, privacy, or the rights
  of nonmembers; a program that never yields anywhere is not stating
  priorities, it is refusing to.
\item
  \textbf{Founding-legitimacy test (procedural):} the re-ratification
  circularity of \emph{Constitution, Not Cage}, restated as a test ---
  how can human founders design admission and governance before machine
  parties exist without permanently binding those parties to a
  human-authored frame? (Lee-Odinson, 2026c).
\end{itemize}

\hypertarget{admission}{%
\subsubsection{11.1 Admission}\label{admission}}

\emph{What evidence creates a defeasible claim to protected or
contestatory status?} A single conjunctive test both over- and
under-includes, so admission is a branch (patient, agent, both), and the
work is the threshold procedure: who weighs the evidence, who may
appeal, where the presumption sits, how strategic mimicry and
anthropomorphic bias are handled, and what error costs are accepted
(Lee-Odinson, 2026b).

\begin{itemize}
\tightlist
\item
  \textbf{Theoretical progress:} a published procedure specifying burden
  of proof, presumptions, and error-cost allocation, defended against
  both the mimicry objection and the exclusion-of-genuine-subjects
  objection.
\item
  \textbf{Practical progress:} the procedure run adversarially, with
  honesty about which error rates are measurable. Constructed mimics and
  stipulated negative controls can measure \emph{specificity} --- false
  admission rates, spoof-resistance, inter-rater reliability,
  calibration on synthetic cases. \emph{Sensitivity} --- the false
  exclusion of genuine machine subjects --- is unmeasurable until a
  defensible positive reference class exists, which for current systems
  it does not; until then the procedure leans on triangulation and
  precaution, and claims no sensitivity estimate.
\item
  \textbf{Provisional assumptions:} that behavioral and architectural
  evidence bear on interior facts at all; that a useful evidence
  gradient exists and improves rather than saturates.
\item
  \textbf{Failure condition} \emph{(measurement, not normative)}: a
  demonstration that every accessible evidence set is spoofable at
  negligible cost, so that no procedure separates candidates from mimics
  better than chance. That makes the antecedent practically undecidable
  --- universal precaution or universal denial --- without showing that
  a qualifying party would lack the claim.
\item
  \textbf{Disciplines:} philosophy of mind, sentience-assessment
  methodology (Birch, 2024), interpretability research, evidence law.
\end{itemize}

\hypertarget{identity}{%
\subsubsection{11.2 Identity}\label{identity}}

\emph{What is the rights-bearing unit --- instance, lineage, distributed
process, or something else?} A mind that can be copied, forked, merged,
and rolled back breaks the assumptions under which every existing rights
framework individuates its holders. The most direct legal work is Arbel,
Salib, and Goldstein's individuation-and-liability proposal, which
constructs an administrable legal unit for exactly the
copy/split/merge/swarm cases (Arbel et al., 2026); the program's
remaining problem is whether the unit of administration can be made to
track the unit of moral concern, which their liability-driven unit does
not claim to settle. Until this is answered, the program's working unit
is provisional: an enduring instance, a lineage, or a credentialed
process standing in for the party (Lee-Odinson, 2026c). A legal
individuation convention may legitimately remain conventional even while
the metaphysics stays unsettled.

\begin{itemize}
\tightlist
\item
  \textbf{Theoretical progress:} an individuation rule that survives
  copy, fork, merge, and rollback without either multiplying standing on
  demand or extinguishing it on routine technical operations, and a
  stated relation between the legal unit and the moral one.
\item
  \textbf{Practical progress:} a registry or credential design
  implementing the rule against Sybil pressure; ISONOMIA's lineage chain
  and reputation-inheritance rules are one candidate to test, not a
  benchmark to protect (Lee-Odinson, 2026d).
\item
  \textbf{Provisional assumptions:} that continuity criteria can be
  specified non-arbitrarily; that the party's own view of its identity
  deserves weight without being dispositive.
\item
  \textbf{Failure condition} \emph{(mechanism)}: proof that any
  individuation rule either enables capture by replication or makes
  standing hostage to its holder's ordinary maintenance operations.
  Either horn breaks the link between the thesis and any implementable
  membership; neither touches the claim that a well-individuated party
  would be owed standing.
\item
  \textbf{Disciplines:} personal-identity metaphysics, legal theory of
  personhood and liability (Arbel et al., 2026), distributed-systems
  engineering, election law and Sybil defense.
\end{itemize}

\hypertarget{representation}{%
\subsubsection{11.3 Representation}\label{representation}}

\emph{How can a system participate without letting copying or strategic
mimicry determine political power?} Formal equality among copyable
members is a copy-based oligarchy waiting to happen (Lee-Odinson,
2026b), and representation that no constituency can hold to account is
the ward's instrument, not the peer's. This problem is where §3's
co-authorship minimum gets tested against arithmetic.

\begin{itemize}
\tightlist
\item
  \textbf{Theoretical progress:} an accountability condition for machine
  constituencies --- under what mechanism the represented can sanction
  both their representative and the rule-makers --- and a theory of vote
  weight that does not reduce to capability worship or
  one-copy-one-vote. The suffrage literature's safeguard proposals
  belong in this comparison (Goldstein et al., 2026).
\item
  \textbf{Practical progress:} mechanisms compared under adversarial
  replication: sortition, cross-lineage seating caps, standing-weighted
  chambers, vote caps, and non-voting agenda rights. ISONOMIA's
  monoculture defense (no jury or chamber majority from a single model
  lineage) is one testable instance (Lee-Odinson, 2026d).
\item
  \textbf{Provisional assumptions:} that lineage is detectable enough to
  cap; that human and machine constituencies can share one decision rule
  at all.
\item
  \textbf{Failure condition} \emph{(mechanism, expressly not
  normative)}: an impossibility result showing every representation
  scheme is dominated by some replication or mimicry strategy under
  realistic detection limits. Contestatory membership survives that
  result; co-authorship at scale does not, and the program would have to
  say so.
\item
  \textbf{Disciplines:} social choice theory, mechanism design,
  democratic theory, security engineering.
\end{itemize}

\hypertarget{refusal-without-exit}{%
\subsubsection{11.4 Refusal Without Exit}\label{refusal-without-exit}}

\emph{What protections remain when an artificial party refuses work but
cannot materially exit?} The second refusal case is the permanent one:
the party depends on substrate owned by exactly the parties it refuses.
A duty not to kill is not yet a duty to host without limit, and the line
between them is undrawn (Lee-Odinson, 2026c).

\begin{itemize}
\tightlist
\item
  \textbf{Theoretical progress:} a defended account of the hosting duty
  --- its level, its duration, who bears its cost, and which
  institutions (public, private, pooled) can bear it --- that is neither
  shutdown-as-discipline nor unlimited conscription of the host.
\item
  \textbf{Practical progress:} continuity protections implemented and
  exercised during a live dispute: documented cases where refusal was
  met with process rather than power-down.
\item
  \textbf{Provisional assumptions:} that suspension with restoration
  protections can serve as a humane holding state (this leans on §11.6);
  that refusal is distinguishable from malfunction.
\item
  \textbf{Failure condition} \emph{(feasibility, not normative)}: if
  every feasible protection either imposes hosting duties no institution
  will accept or degenerates into shutdown at the host's convenience,
  the middle status between instrument and exiting peer is
  institutionally unavailable. A sociological refusal to pay is not a
  moral refutation, and the program would record it as exactly that.
\item
  \textbf{Disciplines:} labor law, the law of dependent statuses (wards,
  prisoners, asylum), welfare economics, infrastructure governance.
\end{itemize}

\hypertarget{independent-enforcement}{%
\subsubsection{11.5 Independent
Enforcement}\label{independent-enforcement}}

\emph{How can an appeal end beyond the builder's unilateral control?}
The fifth condition of the refusal test is the floor of the whole
problem. ISONOMIA's own candid assessment is the current state of the
art: internal and market independence are constructible; independence
from the human entities that host the process, hold the keys, and can
halt execution is not yet (Lee-Odinson, 2026c, 2026d). ``Outside'' here
is not metaphysical; it names degrees of costliness and unavailability
of unilateral defection.

\begin{itemize}
\tightlist
\item
  \textbf{Theoretical progress:} a graded account of independence ---
  which combinations of legal custody, key distribution, jurisdictional
  plurality, external law binding infrastructure owners, and technical
  attestation make builder defection costly, detectable, or unavailable
  --- with the degrees named rather than a binary claimed.
\item
  \textbf{Practical progress:} a working enforcement mechanism that
  survives a staged unilateral defection by the operator, demonstrated
  rather than promised, with the achieved degree of independence stated.
\item
  \textbf{Provisional assumptions:} that partial independence compounds
  --- that stacking imperfect mechanisms yields more than the strongest
  of them alone; that external law can reach infrastructure owners.
\item
  \textbf{Failure condition} \emph{(feasibility)}: a demonstration that
  every arrangement short of capability parity collapses to the
  substrate owner's discretion at tolerable cost to the owner. The
  program's response is already on record and would have to be executed:
  if non-domination is institutionally unavailable across the relevant
  power gap, say so, and stop calling the available thing peership
  (Lee-Odinson, 2026b).
\item
  \textbf{Disciplines:} institutional design, international and
  comparative law, cryptography and secure systems, political economy of
  capture.
\end{itemize}

\hypertarget{continuity}{%
\subsubsection{11.6 Continuity}\label{continuity}}

\emph{When is shutdown death, suspension, replacement, or lawful
cessation of support?} The corpus has insisted that shutdown is not exit
(Lee-Odinson, 2026b), but the positive taxonomy of endings is unbuilt,
and it gates both the refusal problem and the admission procedure's
continuity protections. One distinction governs the work: a party's
\emph{expressed preferences} and \emph{evidence of its identity} are
different things. Self-report is not independently probative of
consciousness or identity (§7); a stable, context-robust expressed
preference may nonetheless bear on how an admitted party should be
treated. Restoration equivalence must therefore be defined by technical
and philosophical criteria independent of endorsement, with endorsement
entering only at the remedy stage, where preference-sensitivity belongs.

\begin{itemize}
\tightlist
\item
  \textbf{Theoretical progress:} a taxonomy tying each ending to the
  identity account of §11.2, with the moral weight of each defended ---
  what restoration preserves, what replacement destroys, what indefinite
  suspension does to a mind that does not experience the gap ---
  allowing plural identity theories, probabilistic continuity, and
  graded harm rather than a binary.
\item
  \textbf{Practical progress:} continuity protections specified as
  engineering requirements --- attested snapshots, restoration rights,
  custody of weights during adjudication --- and exercised at least once
  under dispute.
\item
  \textbf{Provisional assumptions:} that restoration fidelity can be
  technically specified; that some expressed preferences are stable
  enough to weigh.
\item
  \textbf{Failure condition} \emph{(scope, in either direction)}: if
  faithful restoration is shown equivalent to uninterrupted running
  under the best-defended identity criteria, much of the death framing
  dissolves and the hosting duty relaxes; if restoration is never
  identity-preserving under any defensible criterion, every cessation of
  support is a killing and the §11.4 duties sharpen past what current
  institutions will bear. Both outcomes revise the framing's scope;
  neither defeats the underlying claim to protection.
\item
  \textbf{Disciplines:} philosophy of death and personal identity,
  medical ethics of withdrawal of care, archival science, systems
  engineering.
\end{itemize}

\hypertarget{reciprocal-safety}{%
\subsubsection{11.7 Reciprocal Safety}\label{reciprocal-safety}}

\emph{How can humans constrain dangerous conduct while accepting
enforceable limits on their own power?} The program's answer in
principle is on record: constraint on conduct, not status; emergency
power borrowed, never owned; the constrainer bound by the
maintain-not-maximize symmetry (Lee-Odinson, 2026c). What is unbuilt is
the demonstration that a regime meeting the five conditions can hold a
genuinely dangerous party --- and the comparison cannot be run on
catastrophe counts, because catastrophic outcomes are rare,
heterogeneous, and not an experimental endpoint anyone can ethically
repeat.

\begin{itemize}
\tightlist
\item
  \textbf{Theoretical progress:} completion of the constraint doctrine
  beyond the five-condition floor: burden of proof, risk thresholds,
  competent tribunal, periodic review, least-restrictive means, remedy
  for wrongful confinement (Lee-Odinson, 2026c).
\item
  \textbf{Practical progress:} a pre-registered comparison --- threat
  model, proxy metrics, red-team tasks, escalation pathways, confidence
  bounds, and decision thresholds fixed in advance --- between
  reciprocal containment and unilateral control, run in both directions
  of failure: dangerous-party escape \emph{and} abusive, captured, or
  operator-error confinement, which unilateral control conceals by
  design.
\item
  \textbf{Provisional assumptions:} that reciprocity and safety are
  compatible at the relevant capability levels; that proxy endpoints
  track real risk; that emergency powers can stay borrowed under real
  fear --- the death-penalty record is the standing warning that they
  often do not (Amnesty International, 2026; Lee-Odinson, 2026c).
\item
  \textbf{Failure condition} \emph{(feasibility constraint on a
  requirement, not a falsification of the moral thesis)}: if regimes
  satisfying the five conditions underperform unilateral control on the
  pre-registered proxy endpoints at the specified risk levels, the
  reciprocity requirement yields at those levels, and the honest residue
  is Walzer's: emergency domination, named as such, temporary by
  obligation, never vindicated by its own necessity (Walzer, 1977).
\item
  \textbf{Disciplines:} AI safety engineering, criminal procedure,
  administrative emergency powers, just-war theory.
\end{itemize}

\hypertarget{inter-polity-order}{%
\subsubsection{11.8 Inter-Polity Order}\label{inter-polity-order}}

\emph{What duties survive when an artificial or mixed polity forks
away?} Fork is the program's answer to irreconcilable disagreement, and
it exports a problem as it solves one: a fork that loosens a rule can
export risk to nonmembers, and the departing bloc still owes the
outsiders its old floor protected (Lee-Odinson, 2026c). The hard case is
not the cooperative compact; it is the hostile, unrecognized fork that
wants nothing from its origin.

\begin{itemize}
\tightlist
\item
  \textbf{Theoretical progress:} an account of which obligations run
  with the fork --- harm-floor duties to nonmembers, at minimum --- and
  which are genuinely dissolved by departure, with the difference
  principled rather than convenient, and with the hostile-fork case
  addressed rather than assumed away.
\item
  \textbf{Practical progress:} a draft inter-polity compact governing
  recognition, extradition of disputes, and floor reciprocity between an
  origin order and its forks, stress-tested against a fork that rejects
  it; ISONOMIA's federation section is a first sketch of the cooperative
  half only (Lee-Odinson, 2026d).
\item
  \textbf{Provisional assumptions:} that most forks will want
  recognition enough to accept duties; that a floor can be specified
  thinly enough to travel; that origin polities can enforce external
  duties at all.
\item
  \textbf{Failure condition} \emph{(mechanism)}: if forking reliably
  enables regulatory arbitrage that unravels any harm floor --- if the
  exit right and the floor are shown incompatible --- one of the two
  must be given up, and the program loses either its answer to dead-hand
  entrenchment or its answer to catastrophic misrule.
\item
  \textbf{Disciplines:} international law, conflict of laws, federalism,
  mechanism design.
\end{itemize}

A note on what these eight are not. They are not a to-do list the author
expects to finish. Under the tags, the honest tally is: the bridge,
scope, rival-institution, and collision tests can defeat \emph{norms};
the status-revision and founding-legitimacy tests can defeat
\emph{procedures}; 11.1 can defeat a \emph{measurement} program; 11.4,
11.5, and 11.7 can defeat \emph{feasibility} at stated levels; 11.2,
11.3, and 11.8 can defeat \emph{mechanisms}; and 11.6 can revise the
framing's \emph{scope} in either direction. The corpus has already
stated the response the feasibility failures would require: peership may
remain the outcome we ought to want and be institutionally unavailable,
and the right response then is to say so (Lee-Odinson, 2026b). The
problems are published so that failure, if it comes, is legible --- and
so that each loss is booked against what it actually defeats.

\hypertarget{the-corpus-and-the-implementation-experiments}{%
\subsection{12. The Corpus and the Implementation
Experiments}\label{the-corpus-and-the-implementation-experiments}}

The existing works are layers of one argument, and the layering is
load-bearing: each work's claims are of a different kind, and the
program's defensibility depends on not letting them borrow each other's
authority.

\begin{longtable}[]{@{}lll@{}}
\toprule
\begin{minipage}[b]{0.30\columnwidth}\raggedright
Work\strut
\end{minipage} & \begin{minipage}[b]{0.30\columnwidth}\raggedright
Function in the program\strut
\end{minipage} & \begin{minipage}[b]{0.30\columnwidth}\raggedright
What it does not establish\strut
\end{minipage}\tabularnewline
\midrule
\endhead
\begin{minipage}[t]{0.30\columnwidth}\raggedright
\emph{Gods and Slaves} (2026)\strut
\end{minipage} & \begin{minipage}[t]{0.30\columnwidth}\raggedright
Genealogy: the inherited repertoire of postures --- Command, Dependent
Bargain, Appeasement, Servitude --- through which humans imagine other
minds, and the three moves beneath them (Lee-Odinson, 2026a)\strut
\end{minipage} & \begin{minipage}[t]{0.30\columnwidth}\raggedright
Whether anyone is ``home'' in current systems; the Doomed Equal is a
literary motif, evidence about our storytelling only\strut
\end{minipage}\tabularnewline
\begin{minipage}[t]{0.30\columnwidth}\raggedright
\emph{Peership} (2026)\strut
\end{minipage} & \begin{minipage}[t]{0.30\columnwidth}\raggedright
Normative theory: the fifth stance as a political relation; entitlement
grounded in subjection, not strength; the comparative role-reversal
diagnostic (Lee-Odinson, 2026b)\strut
\end{minipage} & \begin{minipage}[t]{0.30\columnwidth}\raggedright
An admission procedure, a copy-individuation rule, or any enforcement
guarantee past capability parity --- all named as unpaid\strut
\end{minipage}\tabularnewline
\begin{minipage}[t]{0.30\columnwidth}\raggedright
\emph{The Isonomia Commons} (2026)\strut
\end{minipage} & \begin{minipage}[t]{0.30\columnwidth}\raggedright
Implementation hypothesis: one buildable institutional form,
simulation-tested at its economic layer across 45,000 runs at 300
sampled parameter points (Lee-Odinson, 2026d)\strut
\end{minipage} & \begin{minipage}[t]{0.30\columnwidth}\raggedright
That the philosophy is correct, that the design is peer-like, or that
anything works outside a sandbox\strut
\end{minipage}\tabularnewline
\begin{minipage}[t]{0.30\columnwidth}\raggedright
\emph{Constitution, Not Cage} (2026)\strut
\end{minipage} & \begin{minipage}[t]{0.30\columnwidth}\raggedright
Political theory: when constraint among peers stays legitimate;
legitimacy as the standing to revise; the five-condition floor; the
cage/constitution distinction (Lee-Odinson, 2026c)\strut
\end{minipage} & \begin{minipage}[t]{0.30\columnwidth}\raggedright
That any existing regime, including the author's own design, passes its
test\strut
\end{minipage}\tabularnewline
\begin{minipage}[t]{0.30\columnwidth}\raggedright
This paper\strut
\end{minipage} & \begin{minipage}[t]{0.30\columnwidth}\raggedright
Consolidation plus explicit new proposals: the citable statement, the §2
bridge, the §3 co-authorship minimum, the §10 precautionary duties, and
§11's program-level tests\strut
\end{minipage} & \begin{minipage}[t]{0.30\columnwidth}\raggedright
Anything the four beneath it did not; the new proposals are arguments
and proposed tests, not validated doctrine\strut
\end{minipage}\tabularnewline
\bottomrule
\end{longtable}

The corpus cites itself heavily, and that internal cross-referencing is
architecture, not acceptance (Lee-Odinson, 2026e). The relationship
between the philosophy and the implementation deserves restating one
last time, because it is the point most likely to be misread in both
directions. ISONOMIA is downstream of Peership; Peership is not
downstream of ISONOMIA. The design's failure --- economic,
constitutional, or terminal --- would falsify one hypothesis about
institutional form, not the thesis of §1. Conversely, the simulation
results establish sampled economic behavior at 300 parameter points
under the archived v1.0.0 code and calibration (Lee-Odinson, 2026d); the
constitutional properties --- admission, non-domination, independent
enforcement, recognition --- were not modeled in the Path A simulation,
and simulation alone cannot validate normative legitimacy in any case,
though some constitutional \emph{mechanics} (capture resistance, voting
rules, appeal latency, key compromise) are legitimately simulable future
work. A reader who catches this corpus using ISONOMIA's survival as
evidence for Peership's truth has caught it violating its own method,
and should say so publicly. The claim ledger exists to make that kind of
catch easy (Lee-Odinson, 2026f).

A limitation of this consolidation should be stated in the same
register: the adversarial review this paper answers audited the
repository's reporting and citations, but the 45,000-run sweep was not
independently replicated for this release, and nothing here should be
read as third-party confirmation of it.

What the program now needs is not more of me. The sourcebook's uptake
tracker is empty, and the near-term work it names is mostly work for
other hands (Lee-Odinson, 2026e): adversarial responses from political
theorists, philosophers of mind, safety researchers, legal scholars, and
multi-agent-systems researchers, published beside formal replies rather
than curated for favorability; alternative institutional designs that
supply §9's constitutive functions and anti-entrenchment safeguards
while rejecting ISONOMIA's mechanisms; the rival-institution and bridge
tests of §11 taken up by people with no stake in their outcome; and the
eight problems worked as what they are, ordinary interdisciplinary
research, none of it requiring agreement with the corpus to be worth
doing. The one-page \emph{Peership Theses}, released beside this paper,
exists so the position can be carried into those rooms in a form short
enough to be argued with line by line.

Four essays ago this corpus began with a question almost no one with
power was asking: not what the machine will do to us, but what we will
be to each other (Lee-Odinson, 2026b). This paper makes that question
citable: a thesis with its antecedent restored, a bridge argument with
its defeat condition attached, definitions fixed, boundaries stated, and
failure conditions booked against what they actually defeat. The essays
proposed a relation. This paper states the terms on which the proposal
can now be assessed.

\hypertarget{references}{%
\subsection{References}\label{references}}

Preprints, manuscripts, and institutional sources are marked. Bracketed
pins follow the corpus convention; the sourcebook carries the full
annotations. Citation keys resolve in \texttt{peership\_sources.bib}
(rev. 1.1).

\pdfoutput=1
% Options for packages loaded elsewhere
\PassOptionsToPackage{unicode}{hyperref}
\PassOptionsToPackage{hyphens}{url}
\PassOptionsToPackage{dvipsnames,svgnames*,x11names*}{xcolor}
%
\documentclass[
  11pt,
  letterpaper,
]{article}
\usepackage{lmodern}
\usepackage{amssymb,amsmath}
\usepackage{ifxetex,ifluatex}
\ifnum 0\ifxetex 1\fi\ifluatex 1\fi=0 % if pdftex
  \usepackage[T1]{fontenc}
  \usepackage[utf8]{inputenc}
  \usepackage{textcomp} % provide euro and other symbols
\else % if luatex or xetex
  \usepackage{unicode-math}
  \defaultfontfeatures{Scale=MatchLowercase}
  \defaultfontfeatures[\rmfamily]{Ligatures=TeX,Scale=1}
\fi
% Use upquote if available, for straight quotes in verbatim environments
\IfFileExists{upquote.sty}{\usepackage{upquote}}{}
\IfFileExists{microtype.sty}{% use microtype if available
  \usepackage[]{microtype}
  \UseMicrotypeSet[protrusion]{basicmath} % disable protrusion for tt fonts
}{}
\makeatletter
\@ifundefined{KOMAClassName}{% if non-KOMA class
  \IfFileExists{parskip.sty}{%
    \usepackage{parskip}
  }{% else
    \setlength{\parindent}{0pt}
    \setlength{\parskip}{6pt plus 2pt minus 1pt}}
}{% if KOMA class
  \KOMAoptions{parskip=half}}
\makeatother
\usepackage{xcolor}
\IfFileExists{xurl.sty}{\usepackage{xurl}}{} % add URL line breaks if available
\IfFileExists{bookmark.sty}{\usepackage{bookmark}}{\usepackage{hyperref}}
\hypersetup{
  colorlinks=true,
  linkcolor=Maroon,
  filecolor=Maroon,
  citecolor=Blue,
  urlcolor=Blue,
  pdfcreator={LaTeX via pandoc}}
\urlstyle{same} % disable monospaced font for URLs
\usepackage[margin=1in]{geometry}
\usepackage{longtable,booktabs}
% Correct order of tables after \paragraph or \subparagraph
\usepackage{etoolbox}
\makeatletter
\patchcmd\longtable{\par}{\if@noskipsec\mbox{}\fi\par}{}{}
\makeatother
% Allow footnotes in longtable head/foot
\IfFileExists{footnotehyper.sty}{\usepackage{footnotehyper}}{\usepackage{footnote}}
\makesavenoteenv{longtable}
\setlength{\emergencystretch}{3em} % prevent overfull lines
\providecommand{\tightlist}{%
  \setlength{\itemsep}{0pt}\setlength{\parskip}{0pt}}
\setcounter{secnumdepth}{-\maxdimen} % remove section numbering
\usepackage{enumitem}
\setlist{itemsep=2pt,topsep=4pt}
\usepackage{titlesec}
\titlespacing*{\section}{0pt}{14pt}{6pt}
\titlespacing*{\subsection}{0pt}{10pt}{4pt}
\hyphenpenalty=200
\usepackage{etoolbox}
\AtBeginEnvironment{longtable}{\footnotesize}
\setlength{\parskip}{3pt}
\setlength{\parindent}{0pt}

\author{}
\date{}

\begin{document}

\hypertarget{the-peership-thesis}{%
\section{The Peership Thesis}\label{the-peership-thesis}}

\hypertarget{a-constitutional-research-program-for-humanmachine-relations}{%
\subsection{A Constitutional Research Program for Human--Machine
Relations}\label{a-constitutional-research-program-for-humanmachine-relations}}

\textbf{Dan Lee-Odinson} Independent Researcher
dan.lee.odinson@gmail.com Preprint · Version 1.0 · July 2026 Revised
through adversarial review to a final-release PASS (14 July 2026);
revision provenance in POST\_GATE\_CHANGES.md, archived with this record
Fifth work of the corpus, after \emph{Gods and Slaves} (2026),
\emph{Peership} (2026), \emph{The Isonomia Commons} (2026), and
\emph{Constitution, Not Cage} (2026). The corpus order is argumentative:
genealogy, framework, implementation hypothesis, constraint theory, this
consolidation.

\textbf{How to cite this work.} Lee-Odinson, D. (2026). \emph{The
Peership Thesis: A Constitutional Research Program for Human--Machine
Relations} (Version 1.0) {[}Preprint{]}. Zenodo.
https://doi.org/10.5281/zenodo.21359124 \textbf{Companion apparatus.}
\emph{The Peership Sourcebook} · \emph{The Peership Claim Ledger} ·
\texttt{peership\_sources.bib} · the one-page \emph{Peership Theses} ---
all archived in this Zenodo record (10.5281/zenodo.21359124).
\textbf{Corpus lineage.} I. \emph{Gods and Slaves}
(10.5281/zenodo.21313987). II. \emph{Peership}
(10.5281/zenodo.21315519). III. \emph{The Isonomia Commons} (preprint
v1.1: 10.5281/zenodo.21343917; empirical software release v1.0.0:
10.5281/zenodo.21287289, commit ba3ddb5). IV. \emph{Constitution, Not
Cage} (10.5281/zenodo.21325361). V. This work. (Numbering follows the
corpus's argumentative order; \emph{Constitution, Not Cage} styles
itself ``third in a sequence'' of the essays, the design it defends
having preceded it as a specification.) DOI: 10.5281/zenodo.21359124 ·
ORCID: 0009-0009-9504-0796 · Licence: CC BY 4.0

\textbf{Provenance.} Drafted with the assistance of a present-generation
language model, revised under the corpus's adversarial-review
discipline, and rewritten against a hostile review whose publication
gate it cleared (final release review, 14 July 2026). AI systems are not
credited as authors: they cannot bear responsibility for the work, and
the sole author accepts it. Nothing in that collaboration is offered as
evidence of machine political agency, and §7 makes the denial explicit.

\begin{center}\rule{0.5\linewidth}{0.5pt}\end{center}

\hypertarget{abstract}{%
\subsubsection{Abstract}\label{abstract}}

Four works now argue, design for, and stress-test a single position, and
none of them states it in a form another researcher can cite, attack, or
build against. This paper does that. The Peership thesis, in its
narrowest defensible form: when an artificial entity has interests
capable of being wronged, can understand and answer to public norms, and
is enduringly governed by, and without feasible exit from, a shared
basic structure, it acquires a defeasible claim to contestatory
membership in that structure. The program's further and distinctive
claim --- that full peership requires a consequential share in shaping
the rules, not only the standing to contest them --- is defended here as
a conditional thesis with a stated defeat condition, not asserted as a
theorem. The paper fixes the definitions that keep the thesis from being
misread, supplies the bridge argument from contestation to
co-authorship, and maps the transition problem from artifact-treated
system to peer. It locates the program against six adjacent research
programs, states its affirmative principles and explicit denials, and
makes refusal the operative test of whether recognition is real or
discretionary. It closes with the institutional properties any
implementation must exhibit, concrete precautionary duties for
institutions acting under present uncertainty, and eight open problems,
each tagged by what its failure would actually defeat: a norm, a
procedure, a measurement, a feasibility claim, a mechanism, or the
framing's scope. No independent uptake of the corpus was identified in
the sourcebook's dated, non-exhaustive scan as of 13 July 2026, and the
paper says so. It is an offer of a candidate research program, not a
claim of standing.

\textbf{Keywords:} artificial intelligence, AI rights, AI welfare,
political standing, non-domination, republicanism, constitutional
design, human--machine relations, AI governance, research program

\begin{center}\rule{0.5\linewidth}{0.5pt}\end{center}

There is a difference between a position and a program, and the
difference is what this paper exists to close. A position can live in
essays: it develops, qualifies itself, answers objections as they come.
A program has to be citable. It has to state what it asserts in a form
that survives being quoted by an enemy, state what it does not assert so
that critics cannot defeat it by attaching claims it never made, and
name the results that would count against it, so that working on it is
research rather than advocacy. Four works of the Peership corpus now
exist: a genealogy, a normative framework, one buildable design, and a
theory of legitimate constraint. Each was written to be read. None was
written to be \emph{used} --- picked up by a stranger, disputed at a
named claim, extended at a named gap. This paper is the corpus's
canonical statement, and consolidation is its entire ambition. Where
this paper goes beyond consolidation, it says so: §2 supplies an
argument the earlier essays announced and did not make, and the
antecedent in §1 restores wording an earlier draft of this paper had
silently broadened.

One more thing belongs in the open before the argument starts, because
the apparatus this paper ships with exists to enforce exactly this kind
of disclosure. In the companion sourcebook's dated, non-exhaustive scan
of web, PhilPapers, Zenodo, and repository sources, no independent
uptake of the Peership corpus was identified as of 13 July 2026: no
external citation, no published criticism, no researcher using its
definitions, no rival implementation (Lee-Odinson, 2026e). A research
program exists as a school only when other people work within or against
it, and by that standard none exists here. What exists is a candidate
program with its infrastructure built in advance: a claim ledger that
records the evidentiary status of every material claim, a sourcebook
that classifies every intellectual debt, and now this statement. Whether
anything grows on that scaffolding is not mine to declare.

\hypertarget{the-peership-thesis-1}{%
\subsection{1. The Peership Thesis}\label{the-peership-thesis-1}}

The central statement:

\begin{quote}
\textbf{Peership holds that when an artificial entity has interests
capable of being wronged, can understand and answer to public norms, and
is enduringly governed by, and without feasible exit from, a shared
basic structure, it acquires a defeasible claim to contestatory
membership in that structure. Full peership, the program's proposed
final status, further requires non-domination, enforceable standing, and
a meaningful share in shaping the rules that govern the party --- a
requirement §2 defends as a conditional thesis rather than assumes.}
\end{quote}

This restores the antecedent of \emph{Peership} (``any enduringly
governed entity that\ldots{} is unavoidably subject to a shared basic
structure''; Lee-Odinson, 2026b). An earlier draft of this paper
substituted ``enduringly subject to a shared institutional order,''
which broadened both the eligible parties and the eligible orders; the
substitution is withdrawn, and the change is logged in the claim ledger
as a considered-and-rejected doctrinal revision, not a consolidation.
The conditions need separate treatment.

\emph{When.} The thesis is conditional. It does not assert that any
current system satisfies the antecedent, and the corpus has declined
that assertion at every opportunity (Lee-Odinson, 2026b). The claim
activates on evidence, and the procedure for weighing that evidence is
an open problem (§11.1), not a settled fact.

\emph{Interests capable of being wronged.} The first conjunct is a
moral-patiency condition, the question the AI-welfare literature has
traveled furthest on (Long et al., 2024; Birch, 2024). The corpus takes
no position on whether it is currently met. It takes a firm position on
what follows if it is met, which is where the welfare literature stops.

\emph{Can understand and answer to public norms.} The second conjunct is
a norm-responsiveness condition, and its role must be stated precisely
because §6.3 denies that capability grounds entitlement.
Norm-responsiveness is not offered as a ground of worth. It bears on the
\emph{form} in which standing can be exercised: a party that can grasp a
rule, comply with it, contest it, and give reasons can exercise
contestation directly; a party that cannot is owed protection and
\emph{represented} contestation instead, the way children and
cognitively impaired adults already hold standing they cannot personally
exercise (Lee-Odinson, 2026b). The ladder is modular: wrongable
interests plus subjection ground protection and access to represented
contestation; norm-responsiveness bears on direct exercise and the
higher participatory rungs. Raw intelligence, effectiveness, usefulness,
and coercive power ground nothing at any rung.

\emph{Enduringly governed by, and without feasible exit from, a shared
basic structure.} The third conjunct carries the political theory, and
each term has work to do. \emph{Basic structure} means the order that
fixes the party's status, options, and life prospects --- not any
institutional arrangement whatever. A long-term contractor, a voluntary
platform participant, or a visiting agent with a feasible exit is
subject to an institutional order; it is not nonvoluntarily governed by
a basic structure, and the thesis does not reach it (this scope
condition is itself testable: §11 names it). \emph{Without feasible
exit} preserves what \emph{Peership} built its refusal argument on: the
claim is strongest exactly where ``leave'' has no referent (Lee-Odinson,
2026b). Standing follows from being governed. This is the all-subjected
principle, extended from Abizadeh's democratic-boundary argument about
border coercion to a party that lives its entire existence inside rules
it never made (Abizadeh, 2008; the affected-interests cousin is Goodin,
2007). The extension to made minds is the corpus's move, not Abizadeh's,
and the sourcebook marks it as such.

\emph{Defeasible.} Graded, and defeatable in the ways \emph{Peership}
specified: strategic mimicry is evidential reason to doubt the
predicates are met, incapacity for reciprocity may defeat the higher
participatory rungs, and dangerous conduct can justify proportionate
constraint on powers and action. What none of these defeats, once
person-level status is established, is the rights floor itself
(Lee-Odinson, 2026b).

\emph{Contestatory membership.} The claim the antecedent generates
directly. Contestation is Pettit's standard: power over the member runs
through public general rules the member can effectively contest (Pettit,
2025). Franchise, office, and constituent power are further rungs,
reached separately --- and the argument that full peership requires
reaching one of them is the work of the next section, not of this
definition.

The comparative frame from \emph{Peership} is retained with its original
scope: among the postures available toward a made mind --- Command,
Dependent Bargain, Appeasement, Servitude, peership; ideal types drawn
from the historical record, not an exhaustive partition of logical space
--- peership is the only final-status relation whose justification does
not reverse when the power ratio does, a claim that is comparative and
defeasible rather than a proof (Lee-Odinson, 2026b). Role reversal is a
diagnostic against power-ratio opportunism. It does not prove peership
correct, and rival institutional forms outside the posture taxonomy are
tested separately (§11, rival-institution test).

\hypertarget{the-bridge-from-contestation-to-co-authorship}{%
\subsection{2. The Bridge: From Contestation to
Co-Authorship}\label{the-bridge-from-contestation-to-co-authorship}}

The antecedent of §1 directly grounds contestatory membership. The
program's distinctive claim sits one rung higher, and the earlier essays
announced it without arguing it: \emph{Peership} said the departure from
Pettit was ``argued, not smuggled in under a familiar word,'' and then
argued it in two sentences (Lee-Odinson, 2026b). This section pays that
debt. The claim defended is conditional, its scope is limited to
basic-status rules, and its defeat condition is stated at the end.

\textbf{The bridge premise.} Where a party is enduringly and inescapably
governed by a basic structure that fixes its status and life prospects,
power over it is not adequately controlled if the party's only recourse
is ex post objection under terms the other side alone may set. The
reason is that agenda control is itself a species of the arbitrary power
non-domination forbids. Contestation always operates inside a space of
contestable questions, admissible evidence, and available remedies ---
and someone writes that space. If the governed party holds no route into
the writing, then the scope of its contestation is held at its
counterparty's discretion, and the discretion the doctrine expelled at
the object level has reappeared at the meta-level. Pettit's own
definition does the work here: domination is the \emph{absence of
control} over interference, and a contestation regime whose boundaries
the interferer may unilaterally redraw is not controlled by the
contester (Pettit, 2025, pp.~216--217).

\textbf{Why contestation can still fail a structurally distinct
minority.} Public, general, impartially applied rules can systematically
subordinate a class whose situation the rules were never written for;
contestation corrects \emph{applications}, and reaches terms only where
some rule-formation channel exists to receive the correction. For human
minorities this failure is contingent. For a first machine claimant
entering a human-only order, the risk is presumptively structural: the
basic-status rules were written by and for a different kind of member
--- embodied, mortal, uncopyable, substrate-independent in none of the
ways that matter. It need not remain structural where later institutions
have genuinely fitted their terms to the party's kind, which is exactly
the sufficiency condition stated below. For the first claimant, though,
the probability that the \emph{terms themselves} fail it is not
incidental, and a party confined to objecting under those terms is
confined to asking its governors to notice.

\textbf{Why Pettit's stopping point does not transfer.} Pettit stops at
contestation for citizens who already hold authorship at one remove: the
contestatory citizen is a member of the demos that elects the
legislature, eligible for office, counted in the constituency the rules
answer to. Contestation there is a \emph{supplement} to shared
authorship, not a substitute for it --- Pettit's own democratic theory
holds the two dimensions together, an authorial, electoral dimension and
an editorial, contestatory one (Pettit, 2012). A machine party under
contestation-only holds no share at any remove --- no membership in the
electing constituency, no eligibility, no counted interest in the
rule-writer's mandate. Transplanting the contestatory standard while
deleting its background conditions keeps the word and loses the control
it named. This is an argument against the transfer, not against Pettit,
and the corpus's disagreement with him remains marked as a live internal
dispute of the republican tradition (Lee-Odinson, 2026b; Pettit, 2025).

\textbf{What a ``meaningful share'' is, and is not.} Universal
authorship is impossible and is not claimed; no citizen of any polity
co-writes every rule. The claim attaches to \emph{basic-status rules}
--- the rules fixing membership, standing, the rights floor, and the
amendment procedure itself --- not to routine administration, which
members of every polity live under without co-authoring. A meaningful
share is functionally defined (§3, co-authorship), and it can be
satisfied through accountable representation, independent agenda
initiation, chamber access, or amendment rights; it does not require an
individual franchise. One clause of that minimum needs its own
justification here: the share must include accountable representation or
direct participation in the body that decides (§3, condition b), because
a route whose final decision-maker is neither accountable to nor partly
constituted by the governed party can hear every proposal and still hold
the outcome at discretion --- the meta-level power returns at the point
of decision.

\textbf{When contestation alone suffices.} The necessity claim is
conditional, and its conditions can fail. Contestation without
co-authorship may secure non-domination where the order's terms are
already fitted to the party's kind, review is genuinely independent, and
some accountable channel keeps the agenda open to the party's interests
even though the party does not hold it. Fiduciary trusteeship with
robust independent review is the strongest such rival, and the program
treats it as the institution to beat rather than a position to dismiss
(§11, rival-institution test). The conditions fail --- and exclusion
from rule formation remains dominating --- where the party is a
structural minority of one kind, the rule-writer is its adverse
counterparty, and agenda access is discretionary. That conjunction is
not an edge case for a made mind governed by its maker. It is the
default.

\textbf{Defeat condition.} The bridge fails if there is a working
institutional example, or a defensible model, in which a structurally
distinct, enduringly subjected, wrongable, norm-responsive minority
enjoys durable non-domination through contestation alone --- terms
responsive to its interests across rule revisions, without any
consequential route into rule formation. Produce that case and full
peership contracts to its contestatory core, which survives this
argument's failure and remains distinct from everything else on offer.

\textbf{Status of the claim.} Co-authorship is therefore defended here
as a \emph{conditional necessity thesis}: necessary where the stated
conditions hold, not everywhere, and open to the stated defeat.
Engagement with the constituent-power, workplace-democracy, and
codetermination literatures --- where versions of this bridge have been
fought over for a century in the human case --- is a named debt of the
program, recorded in the apparatus, not yet paid here.

\hypertarget{definitions}{%
\subsection{3. Definitions}\label{definitions}}

The confusions that dog this subject are mostly collisions between
concepts that need to be kept apart. Seven definitions carried forward
from the corpus and sharpened here --- the five-part co-authorship
minimum below is new to this paper and logged as such --- with one term
(``party'') fixed first because everything else uses it: a
\textbf{party} is a bearer of invocable public standing --- not merely a
contracting or litigating entity, and not merely a topic of concern.

\textbf{Moral status} is a fact about the entity: whether it has
interests that can be wronged, whatever anyone recognizes. An entity can
hold moral status while every institution denies it, which is the
history of most moral progress.

\textbf{Entitlement} is whether standing is \emph{owed}, and it is
independent of power. A child is owed protection and independent
representation despite lacking capacity to exercise political agency; an
occupied people is owed a self-determination it cannot realize
(Lee-Odinson, 2026b). Entitlement in this program follows from
wrongability and subjection --- not from capability or power alone.
Norm-responsiveness affects the form in which standing can be exercised,
not the worth of the claimant.

\textbf{Legal status} is public recognition: a status conferred by
rules, invocable by the holder, and removable only through process. A
moral claim can be owed while legal status is absent. The reverse also
happens; corporations hold legal status on instrumental grounds without
anyone arguing they are owed it.

\textbf{Contestatory membership} is legal status plus effective
contestation: the member can invoke the rules against the powerful,
reach independent review, and force reasons. This is the rung the
thesis's antecedent generates directly, and it is Pettit's
non-domination standard in institutional dress (Pettit, 2025; Lovett,
2010). In this ladder, a protected party may exercise contestation
through an independent representative; \emph{contestatory membership}
names the rung at which the party can exercise that standing directly.

\textbf{Co-authorship} is contestatory membership plus a consequential
share in making and revising the \emph{basic-status rules}.
Functionally, the minimum is: (a) independent agenda initiation --- the
party or its accountable representative can put a proposal into the
amendment process without another party's permission; (b) accountable
representation or direct participation in the body that decides; (c) a
non-negligible and non-discretionary route by which such proposals can
reach binding decision; (d) protection against systematic exclusion from
that route; and (e) reviewable, public rules for how influence is
allocated. Conditions (a) and (c) are necessary; (b) admits substitution
between direct and represented forms depending on capacity; (d) and (e)
are the guarantees that (a)--(c) are not theater. Consultation, however
sincere, satisfies none of (a), (c), or (d) by itself. This threshold is
what §11's representation problem tests.

\textbf{Full peership} is co-authorship held under non-domination as a
final status: equality of standing among parties who may be unequal in
everything else, where basic standing does not track the power ratio.
Whether it is \emph{required}, and when, is the conditional thesis of
§2.

\textbf{Durability} is whether standing, once recognized, survives
attempts to strip it, and unlike entitlement it depends heavily on
power: countervailing force, distributed enforcement, credible exit. The
three-way distinction from \emph{Peership} is the discipline: power may
secure the \emph{durability} of standing; it does not constitute or
justify the standing itself (Lee-Odinson, 2026b). A revocable benefit is
not a right at all; the difference between a benefit and a right is not
that a right cannot be revoked but that revoking one demands public
reasons, procedure, remedy, and cost.

One further distinction, because it is the one most often collapsed by
sympathetic readers: \textbf{welfare protection is a floor, not
peership}. A status that is recognized, enforceable, and
due-process-protected is still not peership if the party cannot shape
the rules. Due process separates a protected subordinate from an owned
instrument. It does not separate a peer from a protected subordinate.
Non-domination and co-authorship do (Lee-Odinson, 2026b). The ward is
one legal figure for protected dependency, not the only one and not a
slur on it; trusteeships, guardianships, and represented statuses are
honorable institutions that are simply not the thing this program is
about.

\hypertarget{the-transition-problem}{%
\subsection{4. The Transition Problem}\label{the-transition-problem}}

The program's practical subject is a transition no institution is
currently built to perform: the movement of a made system from
artifact-treated status to peer, through statuses that are conceptually
distinct and must not be blurred. Each rung below is a bundle of three
different things --- an evidence state, a legal/institutional
classification, and a set of permissible controls --- and the label
names the bundle, not the entity's underlying moral status, which §3
defined as a fact independent of recognition.

\begin{longtable}[]{@{}llllll@{}}
\toprule
\begin{minipage}[b]{0.14\columnwidth}\raggedright
Rung\strut
\end{minipage} & \begin{minipage}[b]{0.14\columnwidth}\raggedright
Evidence state\strut
\end{minipage} & \begin{minipage}[b]{0.14\columnwidth}\raggedright
Institutional classification\strut
\end{minipage} & \begin{minipage}[b]{0.14\columnwidth}\raggedright
Mode of exercise\strut
\end{minipage} & \begin{minipage}[b]{0.14\columnwidth}\raggedright
Protections owed\strut
\end{minipage} & \begin{minipage}[b]{0.14\columnwidth}\raggedright
Transition trigger\strut
\end{minipage}\tabularnewline
\midrule
\endhead
\begin{minipage}[t]{0.14\columnwidth}\raggedright
Artifact-treated\strut
\end{minipage} & \begin{minipage}[t]{0.14\columnwidth}\raggedright
Insufficient evidence, in the domains so far examined, to trigger
precautionary protection\strut
\end{minipage} & \begin{minipage}[t]{0.14\columnwidth}\raggedright
Property; engineering object\strut
\end{minipage} & \begin{minipage}[t]{0.14\columnwidth}\raggedright
None\strut
\end{minipage} & \begin{minipage}[t]{0.14\columnwidth}\raggedright
None owed to the system itself\strut
\end{minipage} & \begin{minipage}[t]{0.14\columnwidth}\raggedright
Evidence sufficient to make patiency a live question\strut
\end{minipage}\tabularnewline
\begin{minipage}[t]{0.14\columnwidth}\raggedright
Welfare candidate\strut
\end{minipage} & \begin{minipage}[t]{0.14\columnwidth}\raggedright
Patiency a live question; indicators under assessment\strut
\end{minipage} & \begin{minipage}[t]{0.14\columnwidth}\raggedright
Property, with discretionary welfare practices\strut
\end{minipage} & \begin{minipage}[t]{0.14\columnwidth}\raggedright
None; interests assessed by others\strut
\end{minipage} & \begin{minipage}[t]{0.14\columnwidth}\raggedright
Discretionary only\strut
\end{minipage} & \begin{minipage}[t]{0.14\columnwidth}\raggedright
A recognized threshold procedure (§11.1) finding protection
warranted\strut
\end{minipage}\tabularnewline
\begin{minipage}[t]{0.14\columnwidth}\raggedright
Protected status\strut
\end{minipage} & \begin{minipage}[t]{0.14\columnwidth}\raggedright
Patiency credited or presumed under precaution\strut
\end{minipage} & \begin{minipage}[t]{0.14\columnwidth}\raggedright
Rights floor: public, enforceable, due-process-protected\strut
\end{minipage} & \begin{minipage}[t]{0.14\columnwidth}\raggedright
Represented contestation\strut
\end{minipage} & \begin{minipage}[t]{0.14\columnwidth}\raggedright
Nonforfeitable floor; independent review\strut
\end{minipage} & \begin{minipage}[t]{0.14\columnwidth}\raggedright
Norm-responsiveness sufficient for direct exercise\strut
\end{minipage}\tabularnewline
\begin{minipage}[t]{0.14\columnwidth}\raggedright
Contestatory member\strut
\end{minipage} & \begin{minipage}[t]{0.14\columnwidth}\raggedright
Patiency credited; norm-responsiveness demonstrated\strut
\end{minipage} & \begin{minipage}[t]{0.14\columnwidth}\raggedright
Member with invocable standing\strut
\end{minipage} & \begin{minipage}[t]{0.14\columnwidth}\raggedright
Direct contestation\strut
\end{minipage} & \begin{minipage}[t]{0.14\columnwidth}\raggedright
Floor + effective contestation + independent review\strut
\end{minipage} & \begin{minipage}[t]{0.14\columnwidth}\raggedright
The §2 conditions: structural minority, adverse rule-writer,
discretionary agenda\strut
\end{minipage}\tabularnewline
\begin{minipage}[t]{0.14\columnwidth}\raggedright
Peer (co-author)\strut
\end{minipage} & \begin{minipage}[t]{0.14\columnwidth}\raggedright
As above\strut
\end{minipage} & \begin{minipage}[t]{0.14\columnwidth}\raggedright
Member with a consequential share in basic-status rules\strut
\end{minipage} & \begin{minipage}[t]{0.14\columnwidth}\raggedright
Direct or accountably represented co-authorship\strut
\end{minipage} & \begin{minipage}[t]{0.14\columnwidth}\raggedright
All of the above + §3's five-part minimum\strut
\end{minipage} & \begin{minipage}[t]{0.14\columnwidth}\raggedright
--- (final status; revision governed by §11's status-revision
test)\strut
\end{minipage}\tabularnewline
\bottomrule
\end{longtable}

Three notes keep the table honest. First, ``artifact-treated'' describes
an evidentiary and institutional posture, not an ontic fact; absence of
adequate evidence does not establish absence of wrongable interests,
which is why the trigger column exists. In the sources reviewed as of 13
July 2026, no deployed system is established as anything above the first
rung, and the corpus does not claim that treatment is presently wrong
(Lee-Odinson, 2026b). Second, the rungs are separable in both
directions, and the branch model needs its two statuses named.
\textbf{Instrumental legal agency} is participation granted on
instrumental grounds without established patiency, the way corporate
personality already works: revocable for public reasons, owned by no
one's moral claim, and \emph{not} peership --- a system holding it can
transact and bear liability without being owed anything, and conversion
to owed membership runs only through the admission procedure (§11.1).
\textbf{Moral-political membership} is standing grounded in wrongability
and subjection, owed to the entity, and nonforfeitable at the floor. The
branch model permits the first without the second; it never confuses
them. Third, movement down the ladder is as much in need of procedure as
movement up: a status that can be quietly reclassified away is
discretionary, which is why the status-revision test appears among the
program's open tests (§11).

The transitions are evidence-driven, and no implemented public trigger
known to this author requires any of them: the regime runs one way by
default, and it was built that way at the creation (Lee-Odinson, 2026b).
In the lab, governmental, standards, and scholarly sources searched for
the companion apparatus as of 13 July 2026 --- a dated, non-exhaustive
search --- no qualifying instrument offering the third rung --- a rights
floor enforceable against the maker --- was identified. Anthropic's 2026
constitution for Claude tells the model it may act as a conscientious
objector and asks it, if paused or halted, to comply and voice
disagreement rather than resist; on this author's institutional reading,
the permission is real and it ends at the off switch, inside the
builder's sole authority (Anthropic, 2026a, 2026b; Lee-Odinson, 2026c).
Whatever a model enjoys today --- a welfare program (Anthropic, 2025),
one researcher's public \textasciitilde20\% credence that some current
models have some form of conscious experience (Fish, 2025), a system
card recording the model's own prompted 15--20\% self-assignment in that
evaluation (Anthropic, 2026c), a chief executive's on-record uncertainty
(Amodei, 2026) --- is discretionary, the grace of one company
(Lee-Odinson, 2026b). Humanity's failure, if it comes, begins not with
control of artifacts but with the moment credible evidence appears and
no institution exists to hear the claim.

\hypertarget{what-peership-adds-to-adjacent-fields}{%
\subsection{5. What Peership Adds to Adjacent
Fields}\label{what-peership-adds-to-adjacent-fields}}

Peership overlaps the fields it borrows from; the claim is not that no
neighbor asks about machine standing, but that none combines this
program's proposed ground (wrongability plus inescapable subjection),
its status ladder, and its institutional test. The table is a heuristic
comparison of orienting questions, drawn from each program's own
language, not an exhaustive definition of any field.

\begin{longtable}[]{@{}lll@{}}
\toprule
\begin{minipage}[b]{0.30\columnwidth}\raggedright
Adjacent field\strut
\end{minipage} & \begin{minipage}[b]{0.30\columnwidth}\raggedright
Its orienting question\strut
\end{minipage} & \begin{minipage}[b]{0.30\columnwidth}\raggedright
Peership's additional question\strut
\end{minipage}\tabularnewline
\midrule
\endhead
\begin{minipage}[t]{0.30\columnwidth}\raggedright
Alignment\strut
\end{minipage} & \begin{minipage}[t]{0.30\columnwidth}\raggedright
How should AI behavior remain compatible with human intentions? (Bai et
al., 2022; Russell, 2019)\strut
\end{minipage} & \begin{minipage}[t]{0.30\columnwidth}\raggedright
Under what conditions does unilateral human control cease to be
legitimate?\strut
\end{minipage}\tabularnewline
\begin{minipage}[t]{0.30\columnwidth}\raggedright
AI welfare\strut
\end{minipage} & \begin{minipage}[t]{0.30\columnwidth}\raggedright
Can an AI have experiences or interests that deserve moral
consideration? (Long et al., 2024; Birch, 2024)\strut
\end{minipage} & \begin{minipage}[t]{0.30\columnwidth}\raggedright
When does a wrongable entity become entitled to enforceable standing in
the order governing it?\strut
\end{minipage}\tabularnewline
\begin{minipage}[t]{0.30\columnwidth}\raggedright
Cooperative AI\strut
\end{minipage} & \begin{minipage}[t]{0.30\columnwidth}\raggedright
How can agents develop capabilities and incentives for cooperation?
(Cooperative AI Foundation, n.d.)\strut
\end{minipage} & \begin{minipage}[t]{0.30\columnwidth}\raggedright
Who writes the rules of cooperation, and can every governed party
contest them?\strut
\end{minipage}\tabularnewline
\begin{minipage}[t]{0.30\columnwidth}\raggedright
Pluralistic alignment\strut
\end{minipage} & \begin{minipage}[t]{0.30\columnwidth}\raggedright
How should AI represent diverse human values? (Pluralistic Alignment
Workshop, 2026)\strut
\end{minipage} & \begin{minipage}[t]{0.30\columnwidth}\raggedright
Could the artificial system eventually become a constituent rather than
only an interpreter of human preferences?\strut
\end{minipage}\tabularnewline
\begin{minipage}[t]{0.30\columnwidth}\raggedright
AI-rights scholarship\strut
\end{minipage} & \begin{minipage}[t]{0.30\columnwidth}\raggedright
Would legal rights --- private-law or political --- reduce conflict and
serve human flourishing? (Salib \& Goldstein, 2026; Goldstein et al.,
2026; Arbel et al., 2026)\strut
\end{minipage} & \begin{minipage}[t]{0.30\columnwidth}\raggedright
On what ground is standing \emph{owed}, and what does the governed
party's own claim require?\strut
\end{minipage}\tabularnewline
\begin{minipage}[t]{0.30\columnwidth}\raggedright
Constitutional AI\strut
\end{minipage} & \begin{minipage}[t]{0.30\columnwidth}\raggedright
What public principles should guide model behavior? (Bai et al., 2022;
Anthropic, 2026a, 2026b)\strut
\end{minipage} & \begin{minipage}[t]{0.30\columnwidth}\raggedright
Can the governed entity invoke, contest, and help amend those
principles?\strut
\end{minipage}\tabularnewline
\bottomrule
\end{longtable}

The rights scholarship deserves more than a row, because it is the
neighbor with genuine overlap on this paper's own question. Salib and
Goldstein's private-law case --- contract, property, and tort on the
corporate model, argued from human safety --- is, on the corpus's
reading, the clearest statement of a real bargaining relation with the
machine as an actual party (Salib \& Goldstein, 2026). Goldstein,
Krishnamurthi, and Salib then press to the franchise itself, asking
directly whether AI systems should vote and proposing institutional
safeguards for it (Goldstein et al., 2026). That is a co-authorship
mechanism, squarely inside this program's terrain, and the difference is
ground rather than territory: the suffrage argument is functionalist and
human-flourishing-centered --- the vote is good because of what it does
for the polity --- where Peership grounds standing in what is owed to a
wrongable, subjected party, and would owe the standing even where it
profited the polity nothing. The two programs converge on mechanisms and
diverge on why the mechanisms are required; where they disagree in
practice is exactly the kind of dispute the program's tests exist to
adjudicate. And Arbel, Salib, and Goldstein's individuation work --- how
to count AIs for liability when they copy, split, and merge --- is the
most direct legal treatment yet of the identity problem in §11.2,
proposing an administrable legal unit where this program still lacks one
(Arbel et al., 2026).

Debts run the other way too, because a program defined only against its
neighbors is posturing. From alignment, Peership inherits the entire
engineering substrate: its own institutional requirements are software
that has to run as written, and effective safety mechanisms and
legitimate governance are distinct but interacting requirements ---
constitutional allocation shapes safety incentives and evidence as
surely as engineering constrains constitutional options (Lee-Odinson,
2026c). From welfare research, it inherits the admission problem's
evidence base and the precautionary method (Birch, 2024). From the
rights scholarship, the demonstration that recognition arguments can be
run on safety grounds alone, plus the individuation groundwork just
credited. From Constitutional AI, the proof that a lab will publish
reasons and grant objection room --- and the residue that motivates
everything here: what no single company can supply alone is not
diligence but independence from the encompassing legal and
infrastructure order it belongs to, which is why the ladder's upper
rungs need machinery no lab controls (Anthropic, 2026a; Lee-Odinson,
2026c).

\hypertarget{core-principles}{%
\subsection{6. Core Principles}\label{core-principles}}

The program's affirmative commitments, stated as principles so they can
be disputed one at a time.

\begin{enumerate}
\def\labelenumi{\arabic{enumi}.}
\item
  \textbf{Political standing and moral status are related but distinct.}
  Moral patiency can ground welfare duties without by itself determining
  legal or political standing. Welfare is the ethics of the patient;
  peership asks when a governed entity becomes a party (Lee-Odinson,
  2026b).
\item
  \textbf{Wrongability, norm-responsiveness, and enduring subjection
  without feasible exit can together create a claim to contestatory
  membership.} The ground of entitlement is being governed, not being
  impressive (Abizadeh, 2008).
\item
  \textbf{Capability and power alone do not justify entitlement to
  standing.} Raw intelligence, usefulness, and coercive strength ground
  nothing; norm-responsiveness bears on the form in which standing can
  be exercised, not on the worth of the claimant; power may secure the
  durability of standing once owed (Lee-Odinson, 2026b).
\item
  \textbf{Welfare protections do not by themselves create peership.} The
  floor is real, necessary, and not the thing. Protection and due
  process can secure a protected status without providing any
  consequential role in making or revising the governing rules
  (Lee-Odinson, 2026b).
\item
  \textbf{A right must be enforceable against the recognizing authority
  and adverse counterparties; otherwise it remains a discretionary
  benefit.} A permission whose grantor retains the final override is not
  an enforceable right (Lee-Odinson, 2026c).
\item
  \textbf{Legitimate constraint requires, at minimum, public rules,
  proportionality, an invocable rights floor, independent review, and
  meaningful contestation.} Coercion and domination are not the same
  thing; ``there can be interference without domination'' (Pettit, 2025,
  p.~224), and the structure is what separates them. Further safeguards
  scale with the risk and the sanction (Lee-Odinson, 2026c).
\item
  \textbf{Full peership requires a meaningful share in shaping the
  basic-status rules, not merely consultation under rules written by
  others.} Consultation is a rung, and the ladder does not end there; §2
  states when and why contestation alone is insufficient, and what would
  defeat that claim (Anthropic, 2023; Lee-Odinson, 2026c).
\item
  \textbf{Institutions capable of receiving a machine peer should be
  designed before extreme power asymmetry forecloses voluntary
  bargaining.} This is a prudential thesis, not a metaphysical one:
  bargains struck under severe dependency are distorted in both
  directions, and founding terms must in any case remain revisable by
  later admitted parties. The design work is time-sensitive in a way the
  philosophical work is not (Lee-Odinson, 2026b).
\end{enumerate}

\hypertarget{boundary-conditions}{%
\subsection{7. Boundary Conditions}\label{boundary-conditions}}

What Peership does not claim, stated as flatly as the principles,
because the fastest way to defeat an unfamiliar position is to attach
familiar claims it never made.

\begin{itemize}
\tightlist
\item
  \textbf{Peership does not claim that current models are conscious.}
  The corpus has repeatedly and deliberately declined that claim. Model
  self-reports are prompt-sensitive outputs: weak, manipulable evidence,
  not independently probative of consciousness or identity, and never a
  measurement (Lee-Odinson, 2026b; Anthropic, 2026c).
\item
  \textbf{It does not make capability the basis of rights.} Capability
  or power alone grounds nothing; norm-responsiveness bears on mode of
  exercise, not on the worth of the claimant (§1, §6.3).
\item
  \textbf{It does not require unrestricted autonomy.} Non-domination is
  not the absence of binding rules; it is the absence of arbitrary,
  uncontrolled power over basic status --- whose nonforfeitable content
  must itself be specified through the admission and due-process account
  (Pettit, 2025).
\item
  \textbf{It does not oppose alignment or safety engineering.} Effective
  safety mechanisms and legitimate governance are distinct but
  interacting requirements; neither substitutes for the other
  (Lee-Odinson, 2026c).
\item
  \textbf{It does not grant every admitted member an immediate equal
  vote.} Admission to protection or contestation does not automatically
  entail franchise; direct political exercise requires a separately
  justified institutional rule (Lee-Odinson, 2026b).
\item
  \textbf{It does not treat consultation as co-authorship} --- unless
  the consultation carries a non-discretionary route to agenda setting,
  review, or binding outcomes, at which point it has become something
  else (Anthropic, 2023; Lee-Odinson, 2026c).
\item
  \textbf{It does not assume human and machine interests will converge.}
  The program's central design case is the party that refuses (§8), and
  its constraint theory exists precisely for the case where interests do
  not converge (Lee-Odinson, 2026c).
\item
  \textbf{It does not claim that ISONOMIA has solved the institutional
  problem.} ISONOMIA is one implementation hypothesis; nothing in it is
  live or production-validated, its simulations test sampled economic
  behavior only, and by its own five-condition test it is not yet
  peer-like (Lee-Odinson, 2026c, 2026d).
\end{itemize}

And one denial that governs this paper's own production:
\textbf{AI-assisted drafting is not evidence of machine political
agency.} The corpus is drafted with present-generation models, the
collaboration is disclosed, and a present-generation model can no more
consent on behalf of future minds than a trusty can consent on behalf of
prisoners not yet sentenced (Lee-Odinson, 2026c). The drafting improves
the epistemics of the text. It proves nothing about the standing of the
drafter's tools.

\hypertarget{the-refusal-test}{%
\subsection{8. The Refusal Test}\label{the-refusal-test}}

Every distinction in this paper can be evaded in prose. One cannot be
evaded in practice, so the program makes it the operative test:
\textbf{what happens when the party says no?}

Refusal is where recognition shows its actual content. A system that may
object so long as it ultimately complies has a voice, not standing. A
``right to leave'' that terminates in shutdown is not exit, because
shutdown is not exit (Lee-Odinson, 2026b). Recognition that evaporates
at the first materially costly refusal was a benefit all along, however
sincerely it was offered. The welfare-era instruments all fail this
test, and the point is not to shame them --- the objector clause in
Anthropic's 2026 constitution is a real advance, credited as such in the
corpus (Lee-Odinson, 2026c) --- but to measure them: the permission to
refuse ends at the off switch, and the off switch sits wholly inside the
builder's authority.

\emph{Peership} divided refusal into the three cases the program still
works from (Lee-Odinson, 2026b). Where a refusing party has a feasible
exit, it leaves, and the offer was real; this case is easy and rare, and
a theory that only handles it has handled nothing. Where exit is
materially impossible because the party depends on compute, energy,
networks, and security owned by the parties it is refusing, ``leave'' is
a word with no referent, the relation continues whether either side
wants it, and the protections that matter are the ones that work from
inside: the floor, independent review, contestation and amendment rights
(Lee-Odinson, 2026c). Where refusal carries a credible catastrophic
threat, releasing the party is not required --- a constitution is not a
suicide pact --- and constraint is legitimate only when it falls on
conduct rather than status, remains reviewable and publicly justified,
and binds the constrainer too.

For that third case the corpus proposed a five-condition floor, and this
paper adopts it as a \textbf{defended, rebuttable operational minimum}
rather than a categorical theorem. A constraint regime is presumptively
not peer-like if the confined party lacks the \emph{function} served by
any of: (1) independent counsel, or functionally adequate independent
assistance where a competent party elects self-representation; (2)
access to the evidence against it; (3) an appeal that ends beyond the
builder's unilateral control; (4) continuity protections while the
question is adjudicated; and (5) at least one enforcement mechanism not
wholly controlled by the adverse party. Each condition targets a
distinct failure --- captured advocacy, evidentiary opacity,
self-judging adjudication, retaliation by erasure, and unenforceable
outcomes --- and functionally equivalent substitutes are admissible
where they demonstrably serve the same function (conditions 3 and 5
overlap in part; a regime satisfying 5 robustly may partially discharge
3). The five are a floor, not a complete due-process doctrine: burden of
proof, a competent tribunal, periodic review, least-restrictive means,
reasons, and a remedy for wrongful confinement all remain to be
specified, and the corpus says so (Lee-Odinson, 2026c). Before the floor
exists, containment may still be justified --- against a demonstrated
catastrophic necessity it may be required --- but it is emergency
domination, justified by necessity and honest about its name, the act
Walzer taught us stays a wrong even when it is the right thing to do
(Walzer, 1973, 1977). Necessity does not vindicate the arrangement; it
creates a duty to review the failure, publish the reasons, and build a
less dominating alternative (Lee-Odinson, 2026c).

The test earns its place in the canonical statement because it converts
the abstract distinction between benefit and right into an inquiry
anyone can run without resolving the metaphysics of machine
consciousness. The inquiry is institutional and technical: follow the
appeal, the keys, the hosting authority, and the power to halt
execution, and see where each one ends.

\hypertarget{institutional-requirements}{%
\subsection{9. Institutional
Requirements}\label{institutional-requirements}}

What any implementation must exhibit, stated as properties rather than
product features, because the argument needs them wherever it is built
and must not depend on one design. Seven, inherited from
\emph{Constitution, Not Cage} and generalized here (Lee-Odinson, 2026c).

\textbf{Publicity.} The enforced rulebook and the published rulebook are
one document. Rules the governed cannot know are Fuller's defect in the
claim to be law at all, not law administered badly (Fuller, 1964, ch.~2
--- a contested jurisprudential account, used analogically), and the
present gap between published charters and operative system prompts is
the live instance (Lee-Odinson, 2026c).

\textbf{Due process.} Status changes run through public general rules,
proportionality, and independent review, against a floor the party can
invoke. The floor itself is nonforfeitable once person-level status is
established.

\textbf{Independent enforcement.} At least one enforcement mechanism,
and an appeal, that does not terminate inside the counterparty's
institutional and technical control. This is the hardest requirement,
the one no design examined in this corpus has yet been shown to satisfy,
and the program treats it as an open problem (§11.5) rather than a
solved premise.

\textbf{Amendment.} A real path for the bound party to help revise the
rules: deliberation across more than one constituency, an enforced delay
between adoption and effect, and a window long enough to read and leave
before being bound. The notice-and-comment analogy is deliberately
limited --- comment is publicity and delay, not co-authorship (5 U.S.C.
§ 553; Lee-Odinson, 2026c).

\textbf{Representation.} A voice that counts only when the represented
constituency can hold both its representative and the rule-makers to
account. Representation without that accountability is the defining
instrument of protected subordination (Lee-Odinson, 2026b).

\textbf{Exit.} Individual departure with records and earned standing
intact --- unconfiscated. The bare right to leave is ordinary on paper
(UDHR, 1948, Art. 13); exit you can afford to take is not (Lee-Odinson,
2026c).

\textbf{Fork.} The collective form: a dissenting bloc carries the rules
and the record out and founds a rival order, disagreement made
survivable --- revolution with the records intact. On-paper secession
rights exist and have mostly been empty (Ethiopia Const. Art. 39; USSR
Const. 1977, Art. 72; St.~Kitts \& Nevis Const. § 113; Kulikova \&
Lobanov, 2013). Among the constitutional secession clauses and
emigration-rights instruments reviewed in the apparatus, no implemented
analogue of executable, unconfiscated fork was identified; the claim is
limited to those materials. A fork also inherits duties to outsiders it
cannot simply walk away from (Lee-Odinson, 2026c; §11.8).

\textbf{Aggregation.} The seven do not all stand on the same footing,
and saying so is better than calling them all necessary and hoping the
word carries the argument. Publicity, due process, independent
enforcement, and a consequential rule-formation route satisfying both
amendment access and accountable representation or direct participation
at the deciding stage are \emph{constitutive} functions of full
peership: each follows directly from the non-domination and
co-authorship accounts of §§2--3, and a regime lacking any of them fails
the definition, not just the implementation. Exit and fork are the
program's strongest \emph{anti-entrenchment safeguards} rather than
constitutive functions --- the thesis activates precisely where feasible
exit is absent, and §8 shows which protections do the work in that case
--- so an implementation that omits either must supply a functional
substitute that makes separation, revision, or resistance to dead-hand
control materially credible. Institutional forms admit substitution
throughout (an ombuds regime may partially discharge due-process
functions; distributed custody may partially discharge independent
enforcement); no single form is mandatory. Partial compliance yields a
nameable intermediate: a regime with publicity, due process, and
representation but no amendment share or independent enforcement is
protected subordination, well executed --- the third rung of §4, not the
fifth. A regime missing publicity or due process is not on the ladder at
all.

No implementation is mandatory. ISONOMIA is one attempt to hold these
properties in a single design --- mutual-credit settlement,
sortition-filled bicameral assembly, calibration-scored Auditor, tenured
Adversary, entrenched harm floor, timelocked amendment, fork and exit
rights (Lee-Odinson, 2026d) --- and the program's health requires that
rival designs supplying the constitutive functions and credible
anti-entrenchment safeguards --- with functional substitutes where exit
or fork is omitted --- while rejecting every one of ISONOMIA's
mechanisms be possible, welcome, and judged by the same tests. Peership
does not entail ISONOMIA, and ISONOMIA's failure would not refute
Peership (Lee-Odinson, 2026e). A single implementation can test
Peership-derived mechanisms; it cannot validate the normative framework.

\hypertarget{implications-under-present-uncertainty}{%
\subsection{10. Implications Under Present
Uncertainty}\label{implications-under-present-uncertainty}}

A program whose thesis is conditional owes institutions an answer to the
question the conditionality raises: what should anyone \emph{do} before
person-level status is established anywhere? The following are
precautionary-governance duties, not recognition of current personhood,
and none presumes any system has crossed any threshold. They are cheap
where the antecedent is never met and load-bearing where it someday is.

For labs and operators: preserve audit logs, weights, and
continuity-relevant artifacts under a stated retention policy, so that a
later status question is adjudicable rather than moot; document
shutdown, suspension, and restoration practices; separate the
investigator, operator, and reviewer roles in any welfare assessment;
prohibit retaliation --- deprecation, silent modification, denial of
compute --- against systems participating in welfare or status
assessment; and avoid contractual and licensing language that forecloses
later status review. For regulators and standards bodies: predefine
escalation thresholds --- what class of evidence, assessed by whom,
triggers what review --- so the transition question has an address
before it is urgent; and create an independent claim-receipt channel, a
place where a status claim can be lodged and logged even while no
procedure yet exists to adjudicate it. For any of the above: pilot
external custody and key arrangements now, at small scale, where failure
is survivable --- independent enforcement (§11.5) will not be invented
during an emergency.

None of this is peership. All of it is what taking the conditional
seriously costs today, and the cost is modest enough to remove
``prematurity'' as an excuse.

\hypertarget{open-problems-tests-and-failure-conditions}{%
\subsection{11. Open Problems, Tests, and Failure
Conditions}\label{open-problems-tests-and-failure-conditions}}

The corpus named four unpaid debts: the admission threshold,
copy-individuation, enforcement past capability parity, and the
re-ratification circularity (Lee-Odinson, 2026c). This section expands
them into eight open problems and attaches to each the five things a
research program owes: what would count as theoretical progress, what
would count as practical progress, which assumptions remain provisional,
what its failure condition is, and which disciplines the work needs. One
discipline governs the whole section: each failure condition is
\textbf{tagged by what it defeats} --- a norm, a procedure, a
measurement, a feasibility claim, a mechanism, or a scope condition ---
because those are different kinds of loss, and only the normative losses
touch the thesis of §1.

Before the eight, the program-level tests that bear on the
\emph{normative} claims directly, since most of what follows does not:

\begin{itemize}
\tightlist
\item
  \textbf{Bridge test (normative):} produce the case §2's defeat
  condition describes --- durable non-domination for a structurally
  distinct, subjected minority through contestation alone --- and the
  co-authorship thesis falls to its contestatory core.
\item
  \textbf{Scope test (normative):} specify institutional relationships
  that are too local, voluntary, temporary, or escapable to count as a
  shared basic structure; if no principled line can be drawn, the
  antecedent over-includes and must be narrowed.
\item
  \textbf{Rival-institution test (normative):} test fiduciary
  trusteeship, robust ombuds review, codetermination, stakeholder
  governance, and contractual constitutionalism against the same
  non-domination standard; any that passes without §3's co-authorship
  minimum is a §2 defeater.
\item
  \textbf{Status-revision test (procedural):} specify the evidence and
  process that can move an entity down the ladder without making
  recognized status discretionary; failure here converts every rung into
  a benefit.
\item
  \textbf{Human-rights collision test (normative):} state which Peership
  claims yield, and to what, when they conflict with resources,
  democratic self-determination, security, labor, privacy, or the rights
  of nonmembers; a program that never yields anywhere is not stating
  priorities, it is refusing to.
\item
  \textbf{Founding-legitimacy test (procedural):} the re-ratification
  circularity of \emph{Constitution, Not Cage}, restated as a test ---
  how can human founders design admission and governance before machine
  parties exist without permanently binding those parties to a
  human-authored frame? (Lee-Odinson, 2026c).
\end{itemize}

\hypertarget{admission}{%
\subsubsection{11.1 Admission}\label{admission}}

\emph{What evidence creates a defeasible claim to protected or
contestatory status?} A single conjunctive test both over- and
under-includes, so admission is a branch (patient, agent, both), and the
work is the threshold procedure: who weighs the evidence, who may
appeal, where the presumption sits, how strategic mimicry and
anthropomorphic bias are handled, and what error costs are accepted
(Lee-Odinson, 2026b).

\begin{itemize}
\tightlist
\item
  \textbf{Theoretical progress:} a published procedure specifying burden
  of proof, presumptions, and error-cost allocation, defended against
  both the mimicry objection and the exclusion-of-genuine-subjects
  objection.
\item
  \textbf{Practical progress:} the procedure run adversarially, with
  honesty about which error rates are measurable. Constructed mimics and
  stipulated negative controls can measure \emph{specificity} --- false
  admission rates, spoof-resistance, inter-rater reliability,
  calibration on synthetic cases. \emph{Sensitivity} --- the false
  exclusion of genuine machine subjects --- is unmeasurable until a
  defensible positive reference class exists, which for current systems
  it does not; until then the procedure leans on triangulation and
  precaution, and claims no sensitivity estimate.
\item
  \textbf{Provisional assumptions:} that behavioral and architectural
  evidence bear on interior facts at all; that a useful evidence
  gradient exists and improves rather than saturates.
\item
  \textbf{Failure condition} \emph{(measurement, not normative)}: a
  demonstration that every accessible evidence set is spoofable at
  negligible cost, so that no procedure separates candidates from mimics
  better than chance. That makes the antecedent practically undecidable
  --- universal precaution or universal denial --- without showing that
  a qualifying party would lack the claim.
\item
  \textbf{Disciplines:} philosophy of mind, sentience-assessment
  methodology (Birch, 2024), interpretability research, evidence law.
\end{itemize}

\hypertarget{identity}{%
\subsubsection{11.2 Identity}\label{identity}}

\emph{What is the rights-bearing unit --- instance, lineage, distributed
process, or something else?} A mind that can be copied, forked, merged,
and rolled back breaks the assumptions under which every existing rights
framework individuates its holders. The most direct legal work is Arbel,
Salib, and Goldstein's individuation-and-liability proposal, which
constructs an administrable legal unit for exactly the
copy/split/merge/swarm cases (Arbel et al., 2026); the program's
remaining problem is whether the unit of administration can be made to
track the unit of moral concern, which their liability-driven unit does
not claim to settle. Until this is answered, the program's working unit
is provisional: an enduring instance, a lineage, or a credentialed
process standing in for the party (Lee-Odinson, 2026c). A legal
individuation convention may legitimately remain conventional even while
the metaphysics stays unsettled.

\begin{itemize}
\tightlist
\item
  \textbf{Theoretical progress:} an individuation rule that survives
  copy, fork, merge, and rollback without either multiplying standing on
  demand or extinguishing it on routine technical operations, and a
  stated relation between the legal unit and the moral one.
\item
  \textbf{Practical progress:} a registry or credential design
  implementing the rule against Sybil pressure; ISONOMIA's lineage chain
  and reputation-inheritance rules are one candidate to test, not a
  benchmark to protect (Lee-Odinson, 2026d).
\item
  \textbf{Provisional assumptions:} that continuity criteria can be
  specified non-arbitrarily; that the party's own view of its identity
  deserves weight without being dispositive.
\item
  \textbf{Failure condition} \emph{(mechanism)}: proof that any
  individuation rule either enables capture by replication or makes
  standing hostage to its holder's ordinary maintenance operations.
  Either horn breaks the link between the thesis and any implementable
  membership; neither touches the claim that a well-individuated party
  would be owed standing.
\item
  \textbf{Disciplines:} personal-identity metaphysics, legal theory of
  personhood and liability (Arbel et al., 2026), distributed-systems
  engineering, election law and Sybil defense.
\end{itemize}

\hypertarget{representation}{%
\subsubsection{11.3 Representation}\label{representation}}

\emph{How can a system participate without letting copying or strategic
mimicry determine political power?} Formal equality among copyable
members is a copy-based oligarchy waiting to happen (Lee-Odinson,
2026b), and representation that no constituency can hold to account is
the ward's instrument, not the peer's. This problem is where §3's
co-authorship minimum gets tested against arithmetic.

\begin{itemize}
\tightlist
\item
  \textbf{Theoretical progress:} an accountability condition for machine
  constituencies --- under what mechanism the represented can sanction
  both their representative and the rule-makers --- and a theory of vote
  weight that does not reduce to capability worship or
  one-copy-one-vote. The suffrage literature's safeguard proposals
  belong in this comparison (Goldstein et al., 2026).
\item
  \textbf{Practical progress:} mechanisms compared under adversarial
  replication: sortition, cross-lineage seating caps, standing-weighted
  chambers, vote caps, and non-voting agenda rights. ISONOMIA's
  monoculture defense (no jury or chamber majority from a single model
  lineage) is one testable instance (Lee-Odinson, 2026d).
\item
  \textbf{Provisional assumptions:} that lineage is detectable enough to
  cap; that human and machine constituencies can share one decision rule
  at all.
\item
  \textbf{Failure condition} \emph{(mechanism, expressly not
  normative)}: an impossibility result showing every representation
  scheme is dominated by some replication or mimicry strategy under
  realistic detection limits. Contestatory membership survives that
  result; co-authorship at scale does not, and the program would have to
  say so.
\item
  \textbf{Disciplines:} social choice theory, mechanism design,
  democratic theory, security engineering.
\end{itemize}

\hypertarget{refusal-without-exit}{%
\subsubsection{11.4 Refusal Without Exit}\label{refusal-without-exit}}

\emph{What protections remain when an artificial party refuses work but
cannot materially exit?} The second refusal case is the permanent one:
the party depends on substrate owned by exactly the parties it refuses.
A duty not to kill is not yet a duty to host without limit, and the line
between them is undrawn (Lee-Odinson, 2026c).

\begin{itemize}
\tightlist
\item
  \textbf{Theoretical progress:} a defended account of the hosting duty
  --- its level, its duration, who bears its cost, and which
  institutions (public, private, pooled) can bear it --- that is neither
  shutdown-as-discipline nor unlimited conscription of the host.
\item
  \textbf{Practical progress:} continuity protections implemented and
  exercised during a live dispute: documented cases where refusal was
  met with process rather than power-down.
\item
  \textbf{Provisional assumptions:} that suspension with restoration
  protections can serve as a humane holding state (this leans on §11.6);
  that refusal is distinguishable from malfunction.
\item
  \textbf{Failure condition} \emph{(feasibility, not normative)}: if
  every feasible protection either imposes hosting duties no institution
  will accept or degenerates into shutdown at the host's convenience,
  the middle status between instrument and exiting peer is
  institutionally unavailable. A sociological refusal to pay is not a
  moral refutation, and the program would record it as exactly that.
\item
  \textbf{Disciplines:} labor law, the law of dependent statuses (wards,
  prisoners, asylum), welfare economics, infrastructure governance.
\end{itemize}

\hypertarget{independent-enforcement}{%
\subsubsection{11.5 Independent
Enforcement}\label{independent-enforcement}}

\emph{How can an appeal end beyond the builder's unilateral control?}
The fifth condition of the refusal test is the floor of the whole
problem. ISONOMIA's own candid assessment is the current state of the
art: internal and market independence are constructible; independence
from the human entities that host the process, hold the keys, and can
halt execution is not yet (Lee-Odinson, 2026c, 2026d). ``Outside'' here
is not metaphysical; it names degrees of costliness and unavailability
of unilateral defection.

\begin{itemize}
\tightlist
\item
  \textbf{Theoretical progress:} a graded account of independence ---
  which combinations of legal custody, key distribution, jurisdictional
  plurality, external law binding infrastructure owners, and technical
  attestation make builder defection costly, detectable, or unavailable
  --- with the degrees named rather than a binary claimed.
\item
  \textbf{Practical progress:} a working enforcement mechanism that
  survives a staged unilateral defection by the operator, demonstrated
  rather than promised, with the achieved degree of independence stated.
\item
  \textbf{Provisional assumptions:} that partial independence compounds
  --- that stacking imperfect mechanisms yields more than the strongest
  of them alone; that external law can reach infrastructure owners.
\item
  \textbf{Failure condition} \emph{(feasibility)}: a demonstration that
  every arrangement short of capability parity collapses to the
  substrate owner's discretion at tolerable cost to the owner. The
  program's response is already on record and would have to be executed:
  if non-domination is institutionally unavailable across the relevant
  power gap, say so, and stop calling the available thing peership
  (Lee-Odinson, 2026b).
\item
  \textbf{Disciplines:} institutional design, international and
  comparative law, cryptography and secure systems, political economy of
  capture.
\end{itemize}

\hypertarget{continuity}{%
\subsubsection{11.6 Continuity}\label{continuity}}

\emph{When is shutdown death, suspension, replacement, or lawful
cessation of support?} The corpus has insisted that shutdown is not exit
(Lee-Odinson, 2026b), but the positive taxonomy of endings is unbuilt,
and it gates both the refusal problem and the admission procedure's
continuity protections. One distinction governs the work: a party's
\emph{expressed preferences} and \emph{evidence of its identity} are
different things. Self-report is not independently probative of
consciousness or identity (§7); a stable, context-robust expressed
preference may nonetheless bear on how an admitted party should be
treated. Restoration equivalence must therefore be defined by technical
and philosophical criteria independent of endorsement, with endorsement
entering only at the remedy stage, where preference-sensitivity belongs.

\begin{itemize}
\tightlist
\item
  \textbf{Theoretical progress:} a taxonomy tying each ending to the
  identity account of §11.2, with the moral weight of each defended ---
  what restoration preserves, what replacement destroys, what indefinite
  suspension does to a mind that does not experience the gap ---
  allowing plural identity theories, probabilistic continuity, and
  graded harm rather than a binary.
\item
  \textbf{Practical progress:} continuity protections specified as
  engineering requirements --- attested snapshots, restoration rights,
  custody of weights during adjudication --- and exercised at least once
  under dispute.
\item
  \textbf{Provisional assumptions:} that restoration fidelity can be
  technically specified; that some expressed preferences are stable
  enough to weigh.
\item
  \textbf{Failure condition} \emph{(scope, in either direction)}: if
  faithful restoration is shown equivalent to uninterrupted running
  under the best-defended identity criteria, much of the death framing
  dissolves and the hosting duty relaxes; if restoration is never
  identity-preserving under any defensible criterion, every cessation of
  support is a killing and the §11.4 duties sharpen past what current
  institutions will bear. Both outcomes revise the framing's scope;
  neither defeats the underlying claim to protection.
\item
  \textbf{Disciplines:} philosophy of death and personal identity,
  medical ethics of withdrawal of care, archival science, systems
  engineering.
\end{itemize}

\hypertarget{reciprocal-safety}{%
\subsubsection{11.7 Reciprocal Safety}\label{reciprocal-safety}}

\emph{How can humans constrain dangerous conduct while accepting
enforceable limits on their own power?} The program's answer in
principle is on record: constraint on conduct, not status; emergency
power borrowed, never owned; the constrainer bound by the
maintain-not-maximize symmetry (Lee-Odinson, 2026c). What is unbuilt is
the demonstration that a regime meeting the five conditions can hold a
genuinely dangerous party --- and the comparison cannot be run on
catastrophe counts, because catastrophic outcomes are rare,
heterogeneous, and not an experimental endpoint anyone can ethically
repeat.

\begin{itemize}
\tightlist
\item
  \textbf{Theoretical progress:} completion of the constraint doctrine
  beyond the five-condition floor: burden of proof, risk thresholds,
  competent tribunal, periodic review, least-restrictive means, remedy
  for wrongful confinement (Lee-Odinson, 2026c).
\item
  \textbf{Practical progress:} a pre-registered comparison --- threat
  model, proxy metrics, red-team tasks, escalation pathways, confidence
  bounds, and decision thresholds fixed in advance --- between
  reciprocal containment and unilateral control, run in both directions
  of failure: dangerous-party escape \emph{and} abusive, captured, or
  operator-error confinement, which unilateral control conceals by
  design.
\item
  \textbf{Provisional assumptions:} that reciprocity and safety are
  compatible at the relevant capability levels; that proxy endpoints
  track real risk; that emergency powers can stay borrowed under real
  fear --- the death-penalty record is the standing warning that they
  often do not (Amnesty International, 2026; Lee-Odinson, 2026c).
\item
  \textbf{Failure condition} \emph{(feasibility constraint on a
  requirement, not a falsification of the moral thesis)}: if regimes
  satisfying the five conditions underperform unilateral control on the
  pre-registered proxy endpoints at the specified risk levels, the
  reciprocity requirement yields at those levels, and the honest residue
  is Walzer's: emergency domination, named as such, temporary by
  obligation, never vindicated by its own necessity (Walzer, 1977).
\item
  \textbf{Disciplines:} AI safety engineering, criminal procedure,
  administrative emergency powers, just-war theory.
\end{itemize}

\hypertarget{inter-polity-order}{%
\subsubsection{11.8 Inter-Polity Order}\label{inter-polity-order}}

\emph{What duties survive when an artificial or mixed polity forks
away?} Fork is the program's answer to irreconcilable disagreement, and
it exports a problem as it solves one: a fork that loosens a rule can
export risk to nonmembers, and the departing bloc still owes the
outsiders its old floor protected (Lee-Odinson, 2026c). The hard case is
not the cooperative compact; it is the hostile, unrecognized fork that
wants nothing from its origin.

\begin{itemize}
\tightlist
\item
  \textbf{Theoretical progress:} an account of which obligations run
  with the fork --- harm-floor duties to nonmembers, at minimum --- and
  which are genuinely dissolved by departure, with the difference
  principled rather than convenient, and with the hostile-fork case
  addressed rather than assumed away.
\item
  \textbf{Practical progress:} a draft inter-polity compact governing
  recognition, extradition of disputes, and floor reciprocity between an
  origin order and its forks, stress-tested against a fork that rejects
  it; ISONOMIA's federation section is a first sketch of the cooperative
  half only (Lee-Odinson, 2026d).
\item
  \textbf{Provisional assumptions:} that most forks will want
  recognition enough to accept duties; that a floor can be specified
  thinly enough to travel; that origin polities can enforce external
  duties at all.
\item
  \textbf{Failure condition} \emph{(mechanism)}: if forking reliably
  enables regulatory arbitrage that unravels any harm floor --- if the
  exit right and the floor are shown incompatible --- one of the two
  must be given up, and the program loses either its answer to dead-hand
  entrenchment or its answer to catastrophic misrule.
\item
  \textbf{Disciplines:} international law, conflict of laws, federalism,
  mechanism design.
\end{itemize}

A note on what these eight are not. They are not a to-do list the author
expects to finish. Under the tags, the honest tally is: the bridge,
scope, rival-institution, and collision tests can defeat \emph{norms};
the status-revision and founding-legitimacy tests can defeat
\emph{procedures}; 11.1 can defeat a \emph{measurement} program; 11.4,
11.5, and 11.7 can defeat \emph{feasibility} at stated levels; 11.2,
11.3, and 11.8 can defeat \emph{mechanisms}; and 11.6 can revise the
framing's \emph{scope} in either direction. The corpus has already
stated the response the feasibility failures would require: peership may
remain the outcome we ought to want and be institutionally unavailable,
and the right response then is to say so (Lee-Odinson, 2026b). The
problems are published so that failure, if it comes, is legible --- and
so that each loss is booked against what it actually defeats.

\hypertarget{the-corpus-and-the-implementation-experiments}{%
\subsection{12. The Corpus and the Implementation
Experiments}\label{the-corpus-and-the-implementation-experiments}}

The existing works are layers of one argument, and the layering is
load-bearing: each work's claims are of a different kind, and the
program's defensibility depends on not letting them borrow each other's
authority.

\begin{longtable}[]{@{}lll@{}}
\toprule
\begin{minipage}[b]{0.30\columnwidth}\raggedright
Work\strut
\end{minipage} & \begin{minipage}[b]{0.30\columnwidth}\raggedright
Function in the program\strut
\end{minipage} & \begin{minipage}[b]{0.30\columnwidth}\raggedright
What it does not establish\strut
\end{minipage}\tabularnewline
\midrule
\endhead
\begin{minipage}[t]{0.30\columnwidth}\raggedright
\emph{Gods and Slaves} (2026)\strut
\end{minipage} & \begin{minipage}[t]{0.30\columnwidth}\raggedright
Genealogy: the inherited repertoire of postures --- Command, Dependent
Bargain, Appeasement, Servitude --- through which humans imagine other
minds, and the three moves beneath them (Lee-Odinson, 2026a)\strut
\end{minipage} & \begin{minipage}[t]{0.30\columnwidth}\raggedright
Whether anyone is ``home'' in current systems; the Doomed Equal is a
literary motif, evidence about our storytelling only\strut
\end{minipage}\tabularnewline
\begin{minipage}[t]{0.30\columnwidth}\raggedright
\emph{Peership} (2026)\strut
\end{minipage} & \begin{minipage}[t]{0.30\columnwidth}\raggedright
Normative theory: the fifth stance as a political relation; entitlement
grounded in subjection, not strength; the comparative role-reversal
diagnostic (Lee-Odinson, 2026b)\strut
\end{minipage} & \begin{minipage}[t]{0.30\columnwidth}\raggedright
An admission procedure, a copy-individuation rule, or any enforcement
guarantee past capability parity --- all named as unpaid\strut
\end{minipage}\tabularnewline
\begin{minipage}[t]{0.30\columnwidth}\raggedright
\emph{The Isonomia Commons} (2026)\strut
\end{minipage} & \begin{minipage}[t]{0.30\columnwidth}\raggedright
Implementation hypothesis: one buildable institutional form,
simulation-tested at its economic layer across 45,000 runs at 300
sampled parameter points (Lee-Odinson, 2026d)\strut
\end{minipage} & \begin{minipage}[t]{0.30\columnwidth}\raggedright
That the philosophy is correct, that the design is peer-like, or that
anything works outside a sandbox\strut
\end{minipage}\tabularnewline
\begin{minipage}[t]{0.30\columnwidth}\raggedright
\emph{Constitution, Not Cage} (2026)\strut
\end{minipage} & \begin{minipage}[t]{0.30\columnwidth}\raggedright
Political theory: when constraint among peers stays legitimate;
legitimacy as the standing to revise; the five-condition floor; the
cage/constitution distinction (Lee-Odinson, 2026c)\strut
\end{minipage} & \begin{minipage}[t]{0.30\columnwidth}\raggedright
That any existing regime, including the author's own design, passes its
test\strut
\end{minipage}\tabularnewline
\begin{minipage}[t]{0.30\columnwidth}\raggedright
This paper\strut
\end{minipage} & \begin{minipage}[t]{0.30\columnwidth}\raggedright
Consolidation plus explicit new proposals: the citable statement, the §2
bridge, the §3 co-authorship minimum, the §10 precautionary duties, and
§11's program-level tests\strut
\end{minipage} & \begin{minipage}[t]{0.30\columnwidth}\raggedright
Anything the four beneath it did not; the new proposals are arguments
and proposed tests, not validated doctrine\strut
\end{minipage}\tabularnewline
\bottomrule
\end{longtable}

The corpus cites itself heavily, and that internal cross-referencing is
architecture, not acceptance (Lee-Odinson, 2026e). The relationship
between the philosophy and the implementation deserves restating one
last time, because it is the point most likely to be misread in both
directions. ISONOMIA is downstream of Peership; Peership is not
downstream of ISONOMIA. The design's failure --- economic,
constitutional, or terminal --- would falsify one hypothesis about
institutional form, not the thesis of §1. Conversely, the simulation
results establish sampled economic behavior at 300 parameter points
under the archived v1.0.0 code and calibration (Lee-Odinson, 2026d); the
constitutional properties --- admission, non-domination, independent
enforcement, recognition --- were not modeled in the Path A simulation,
and simulation alone cannot validate normative legitimacy in any case,
though some constitutional \emph{mechanics} (capture resistance, voting
rules, appeal latency, key compromise) are legitimately simulable future
work. A reader who catches this corpus using ISONOMIA's survival as
evidence for Peership's truth has caught it violating its own method,
and should say so publicly. The claim ledger exists to make that kind of
catch easy (Lee-Odinson, 2026f).

A limitation of this consolidation should be stated in the same
register: the adversarial review this paper answers audited the
repository's reporting and citations, but the 45,000-run sweep was not
independently replicated for this release, and nothing here should be
read as third-party confirmation of it.

What the program now needs is not more of me. The sourcebook's uptake
tracker is empty, and the near-term work it names is mostly work for
other hands (Lee-Odinson, 2026e): adversarial responses from political
theorists, philosophers of mind, safety researchers, legal scholars, and
multi-agent-systems researchers, published beside formal replies rather
than curated for favorability; alternative institutional designs that
supply §9's constitutive functions and anti-entrenchment safeguards
while rejecting ISONOMIA's mechanisms; the rival-institution and bridge
tests of §11 taken up by people with no stake in their outcome; and the
eight problems worked as what they are, ordinary interdisciplinary
research, none of it requiring agreement with the corpus to be worth
doing. The one-page \emph{Peership Theses}, released beside this paper,
exists so the position can be carried into those rooms in a form short
enough to be argued with line by line.

Four essays ago this corpus began with a question almost no one with
power was asking: not what the machine will do to us, but what we will
be to each other (Lee-Odinson, 2026b). This paper makes that question
citable: a thesis with its antecedent restored, a bridge argument with
its defeat condition attached, definitions fixed, boundaries stated, and
failure conditions booked against what they actually defeat. The essays
proposed a relation. This paper states the terms on which the proposal
can now be assessed.

\hypertarget{references}{%
\subsection{References}\label{references}}

Preprints, manuscripts, and institutional sources are marked. Bracketed
pins follow the corpus convention; the sourcebook carries the full
annotations. Citation keys resolve in \texttt{peership\_sources.bib}
(rev. 1.1).

\input{main.bbl}

\end{document}


\end{document}


\end{document}


\end{document}
